% FGA Explained, Part 4. Original PDF 191--246, printed pp.179--234.
% Native reconstruction. Coverage and review boundaries: evidence/ch08-review.md.
% Source PDF 191 / printed p.179; PDF 192 is blank.
\part[形式几何中的 Grothendieck 存在定理]{形式几何中的 Grothendieck 存在定理\\[1em]{\large 附 Jean-Pierre Serre 的一封信}\\[1em]{\large Luc Illusie}}
% Source PDF 193 / printed p.181.
\chapter{形式几何中的 Grothendieck 存在定理}\label{ch08:chapter}
\section*{引言}
本讲义的主题是 Grothendieck 在 Séminaire Bourbaki 上所作的第 182 号报告 [Gro95a]. 该文中的大部分材料已在 [EGAIII1]、[EGAIII2] 和 [SGA1] 中得到详尽处理. 我们的目的是给出一个导论, 解释关键定理的证明, 讨论典型应用, 并在必要时补入较新的进展. 核心结果是固有态射的形式上同调与代数上同调之间的比较定理, 以及形式层的存在定理. 我们先在 §8.1 中回顾形式概形的若干基本术语 (始终限于局部 Noether 的情形, 这对我们的目的已经足够), 然后在 §8.2 和 §8.4 中给出证明的要点. §8.3 重新讨论 [EGAIII2, 7] 中的若干问题: 基变换公式与上同调平坦性. 我们认为, 导出范畴, 特别是完美复形的使用, 能使叙述简化. 不过, 后文并不依赖这一节. §8.5 把前几节的 Grothendieck 定理与形变理论的基本结果相结合, 主要在光滑情形下讨论形式提升或代数提升的存在性的一些应用. 最后, §8.6 讨论 Serre 在 [Ser61] 中给出的例子: 特征 $p>0$ 上不能提升到特征零的射影光滑概形.

我十分感谢 Serre 就他在 [Ser61] 中的例子与我交谈, 寄给我 Mumford 致他的未发表信件 [Mum] 的副本, 并允许我在此收入他写给我的一封信, 其中解决了 [Ser61] 中遗留的一个问题. Esnault、Messing、Raynaud 和 Serre 阅读了本讲义的初稿, 提出了多处订正与修改意见. 我衷心感谢他们, 也感谢暑期学校的学生提出的大量问题和意见.

\emph{约定.} 我们用符号 $\lim$ (相应地, $\operatorname{colim}$) 表示射影极限 (相应地, 归纳极限).

\section{局部 Noether 形式概形}\label{ch08:sec81}
\paragraph{8.1.1.}
\emph{进 Noether 环 (adic Noetherian ring)} 是指赋有满足下述性质的拓扑的 Noether 环 $A$: 存在一个理想 $I$, 使其各次幂 $I^n$ ($n\geq0$) 组成 $A$ 中 $0$ 的一组基本邻域, 而且 $A$ 对此拓扑是分离且完备的; 即 $A$ 是离散环 $A_n=A/I^{n+1}$ ($n\geq0$) 的射影极限. $A$ 的一个\emph{定义理想 (ideal of definition)} 是满足这一性质的理想 $I$; 等价地, 它是一个开理想且其各次幂趋于 $0$. 若 $I$ 是定义理想, 则称 $A$ 是 $I$ \emph{进的}, 此拓扑称为 $I$ \emph{进拓扑}, $A$ 按 $I$ 的各次幂所作的滤过称为 $I$ \emph{进滤过}.
% Source PDF 194 / printed p.182.
若 $I$ 是定义理想, 则理想 $J$ 是定义理想的充要条件是存在正整数 $p,q$, 使得 $J\supset I^p\supset J^q$.

与一个 $I$ 进 Noether 环 $A$ 相联系的是一个带拓扑的赋环空间
\begin{equation}\mathcal X=(\operatorname{Spf}A,\mathcal O_{\mathcal X}),\tag{8.1.1.1}\label{ch08:eq8111}\end{equation}
其定义如下. 对 $n\in\mathbb N$, 令 $X_n=\Spec A_n$ ($A_n=A/I^{n+1}$). 这些概形构成 $\Spec A$ 的一个递增闭子概形序列 (转移映射为闭幂零浸入)
\[
X_0=\Spec A_0\longrightarrow X_1\longrightarrow\cdots\longrightarrow X_n\longrightarrow\cdots.
\]
它们有相同的底空间 $\mathcal X$, 称为 $A$ 的\emph{形式谱 (formal spectrum)}. 注意, $I$ 包含在 $A$ 的根基中 [EGAI, Ch.~0, 7.1.10]; 即对所有 $x\in I$, $1-x$ 都可逆, 这意味着 $\mathcal X$ 作为 $\Spec A$ 的一个闭子空间, 包含它的所有闭点. $\Spec A$ 中任何包含 $\mathcal X$ 的开子集都等于 $\Spec A$. 环层 $\mathcal O_{\mathcal X}$ 定义为 $\mathcal X$ 上各层 $\mathcal O_{X_n}$ 的逆极限, 并赋予自然拓扑, 使得对 $\mathcal X$ 的任何开子集 $U$, 有 $\Gamma(U,\mathcal O_{\mathcal X})=\lim_n\Gamma(U,\mathcal O_{X_n})$, 其中 $\Gamma(U,\mathcal O_{X_n})$ 带离散拓扑. 特别地, $\Gamma(\mathcal X,\mathcal O_{\mathcal X})=A$; 对 $f\in A$, 若 $D(f)=\mathcal X_f$ 是 $\mathcal X$ 中 $f$ 在 $A_0$ 中的像可逆的开子集, 则 $\Gamma(\mathcal X_f,\mathcal O_{\mathcal X})=A_{\{f\}}$, 即完备分式环 $\lim_n S_f^{-1}A/S_f^{-1}I^{n+1}$. 茎
\[
\mathcal O_{\mathcal X,x}=\operatorname{colim}_{f\in A,\,f(x)\ne0}A_{\{f\}}
\]
是局部环 (但不完备).

带拓扑的赋环空间 (8.1.1.1) 只依赖于作为拓扑环的 $A$. 把 $I$ 换成另一个定义理想并不改变它. 空间 $\mathcal X$ 是由 $A$ 的开素理想组成的 $\Spec A$ 的子空间, 而 $\mathcal O_{\mathcal X}$ 是各层 $(A/J)^\sim$ 的逆极限, 其中 $J$ 遍历 $A$ 的定义理想.

\emph{仿射 Noether 形式概形 (affine Noetherian formal scheme)} 是指同构于形如 (8.1.1.1) 的带拓扑的赋环空间. \emph{局部 Noether 形式概形 (locally Noetherian formal scheme)} 是指每一点都有一个开邻域为仿射 Noether 形式概形的带拓扑赋环空间. 若其底空间为 Noether 空间, 则称它是 \emph{Noether 的}. 局部 Noether 形式概形之间的一个\emph{态射} $f:\mathcal X\to\mathcal Y$ 是指一个局部且连续的赋环空间态射: 局部是说对每个 $x\in\mathcal X$, 映射 $\mathcal O_{\mathcal Y,f(x)}\to\mathcal O_{\mathcal X,x}$ 都是局部同态; 连续是说对每个仿射开集 $V\subset\mathcal Y$, 映射 $\Gamma(V,\mathcal O_{\mathcal Y})\to\Gamma(f^{-1}(V),\mathcal O_{\mathcal X})$ 都连续. 局部 Noether 形式概形以显然的方式组成一个范畴.

与通常概形的情形一样, 可以验证: 若 $\mathcal Y=\operatorname{Spf}(A)$ 是 Noether 仿射形式概形 (下文通常从记号中略去环层), 而 $\mathcal X$ 是任意局部 Noether 形式概形, 则
\begin{equation}\Hom(\mathcal X,\mathcal Y)=\Hom_{\mathrm{cont}}(A,\Gamma(\mathcal X,\mathcal O_{\mathcal X})),\tag{8.1.1.2}\label{ch08:eq8112}\end{equation}
其中 $\Hom_{\mathrm{cont}}$ 表示连续环同态的集合. 特别地, 若 $\mathcal X$ 仿射且其环为 $B$, 则
\[\Hom(\mathcal X,\mathcal Y)=\Hom_{\mathrm{cont}}(A,B).\]

\paragraph{8.1.2.}
设 $\mathcal X=\operatorname{Spf}A$ 是一个仿射 Noether 形式概形, $I$ 是 $A$ 的一个定义理想. 设 $M$ 是有限型 $A$ 模. 与 $M$ 相联系的是 $X=\Spec A$ 上的凝聚模 $\widetilde M$. 类似地, 与 $M$ 相联系的还有 $\mathcal X$ 上的模 $M^\triangle$, 定义如下. 对 $n\in\mathbb N$, 如 8.1.1 中那样令 $X_n=\Spec A/I^{n+1}$. 置
\[M^\triangle=\lim_n\widetilde M_n,\]
% Source PDF 195 / printed p.183.
其中 $M_n=M/I^{n+1}M$. 容易验证, $M^\triangle$ 不依赖于 $I$ 的选择, 函子 $M\mapsto M^\triangle$ 正合,
\[\Gamma(\mathcal X,M^\triangle)=M,\]
并且 $M^\triangle$ 的形成与张量积和内部 $\mathcal Hom$ 交换. 关键在于: 若
\[i:\mathcal X\longrightarrow X\]
是由底拓扑空间的包含映射以及环层上的典范映射 $\mathcal O_X\to\mathcal O_{\mathcal X}$ 定义的自然态射, 那么由于 $M$ 有限型, Krull 定理给出
\[M^\triangle=i^*\widetilde M.\]
因为对任何 $f\in A$, $A_{\{f\}}$ 都是进 Noether 环, 所以 $\mathcal O_{\mathcal X}$ 是凝聚环层, $M^\triangle$ 是凝聚模, 而 $\mathcal X$ 上的凝聚模恰为形如 $M^\triangle$ 的模, 其中 $M$ 是有限型 $A$ 模.

\paragraph{8.1.3.}
局部 Noether 形式概形可以更方便地描述为一个递增幂零加厚链的余极限, 而且在实际中通常正是这样出现的. 所谓\emph{加厚 (thickening)}, 是指其理想 $I$ 为幂零元理想的概形闭浸入 $X\to X'$; 此时概形 $X$ 与 $X'$ 有相同的底空间. 若 $X'$ 是 Noether 的, 则 $X$ 也是 Noether 的且 $I$ 幂零; 反过来, 若 $X$ 是 Noether 的, $I/I^2$ 是凝聚 $\mathcal O_X$ 模, 则 $X'$ 是 Noether 的 [EGAI, Ch.~0, 7.2.6, 10.6.4]. 若 $X'$ 是 Noether 的, 则 $X'$ 仿射当且仅当 $X$ 仿射 [EGAI, 6.1.7]. 若 $I^{n+1}=0$, 则称该加厚为 $n$ \emph{阶}加厚.

设 $\mathcal X$ 是局部 Noether 形式概形. 由仿射情形的讨论可知, $\mathcal O_{\mathcal X}$ 是凝聚环层, 而 $\mathcal X$ 上的凝聚模恰为有限呈示模; 等价地, 它们在任何仿射开集 $\mathcal U=\operatorname{Spf}A$ 上都具有 $M^\triangle$ 的形式, 其中 $M$ 为有限型 $A$ 模.

$\mathcal X$ 的一个\emph{定义理想}是 $\mathcal O_{\mathcal X}$ 的一个凝聚理想 $\mathcal I$, 使得对任意 $x\in\mathcal X$, 存在 $x$ 的仿射邻域 $\mathcal U=\operatorname{Spf}A$, 使 $\mathcal I|\mathcal U$ 具有 $I^\triangle$ 的形式, 其中 $I$ 是 $A$ 的一个定义理想. 凝聚理想 $\mathcal I$ 是定义理想, 当且仅当赋环空间 $(\mathcal X,\mathcal O_{\mathcal X}/\mathcal I)$ 是以 $\mathcal X$ 为底空间的概形. $\mathcal X$ 的定义理想存在. 实际上, 存在一个最大的定义理想
\begin{equation}\mathcal T=\mathcal T_{\mathcal X},\tag{8.1.3.1}\label{ch08:eq8131}\end{equation}
它是使 $(\mathcal X,\mathcal O_{\mathcal X}/\mathcal I)$ 约化的唯一一个定义理想 $\mathcal I$. 若 $\mathcal U=\operatorname{Spf}A$ 是仿射开子集, 则 $\mathcal T|\mathcal U=T^\triangle$, 其中 $T$ 是由 $A$ 的拓扑幂零元 $a$ 组成的理想, 即 $a$ 在 $A/I$ 中的像幂零. 若 $\mathcal I$ 是 $\mathcal X$ 的定义理想, 则对任何 $n\geq1$, $\mathcal I^n$ 也是定义理想. 若 $\mathcal X$ 是 Noether 的, 则如仿射情形一样, 当 $\mathcal I$ 是 $\mathcal X$ 的定义理想而 $\mathcal J$ 是凝聚理想时, $\mathcal J$ 是定义理想的充要条件是存在正整数 $p,q$, 使 $\mathcal J\supset\mathcal I^p\supset\mathcal J^q$.

固定 $\mathcal X$ 的一个定义理想 $\mathcal I$. 对 $n\in\mathbb N$, 赋环空间 $(\mathcal X,\mathcal O_{\mathcal X}/\mathcal I^{n+1})$ 是一个局部 Noether 概形 $X_n$, 并且有递增加厚链
\begin{equation}X_\bullet=(X_0\longrightarrow X_1\longrightarrow\cdots\longrightarrow X_n\longrightarrow\cdots),\tag{8.1.3.2}\label{ch08:eq8132}\end{equation}
其余极限 (在局部 Noether 形式概形范畴中) 为 $\mathcal X$: 各加厚在底空间上诱导恒等映射, 这些底空间都等于 $\mathcal X$ 的底空间,
% Source PDF 196 / printed p.184.
并且作为拓扑环有
\[\mathcal O_{\mathcal X}=\lim_n\mathcal O_{X_n}\]
(在任何仿射开集 $U$ 上, $\Gamma(U,\mathcal O_{X_n})$ 带离散拓扑). 设 $J_n$ 为 $X_0$ 在 $X_n$ 中的理想, 即 $J_n=\ker(\mathcal O_{X_n}\to\mathcal O_{X_0})$. 则对 $m\leq n$, $X_m$ 在 $X_n$ 中的理想为 $J_n^{m+1}$ (特别地, $J_n^{n+1}=0$), $J_1$ 是 $X_0$ 上的凝聚模, 而且 $J_n=\mathcal I/\mathcal I^{n+1}$.

反过来, 考虑满足下列条件的赋环空间序列 (8.1.3.2):
\begin{enumerate}
\item[(i)] $X_0$ 是局部 Noether 概形;
\item[(ii)] 底拓扑空间的映射均为同胚, 借助它们将各底空间等同后, 环映射 $\mathcal O_{X_{n+1}}\to\mathcal O_{X_n}$ 满射;
\item[(iii)] 若 $J_n=\ker(\mathcal O_{X_n}\to\mathcal O_{X_0})$, 则对 $m\leq n$, 有 $\ker(\mathcal O_{X_n}\to\mathcal O_{X_m})=J_n^{m+1}$;
\item[(iv)] $J_1$ 作为 $\mathcal O_{X_0}$ 模是凝聚的.
\end{enumerate}
那么带拓扑的赋环空间 $\mathcal X=(X_0,\lim_n\mathcal O_{X_n})$ 是局部 Noether 形式概形; 若 $\mathcal I=\ker(\mathcal O_{\mathcal X}\to\mathcal O_{X_0})=\lim J_n$, 则 $\mathcal I$ 是 $\mathcal X$ 的定义理想, 且 $\mathcal I^{n+1}=\ker(\mathcal O_{\mathcal X}\to\mathcal O_{X_n})$.

验证是直接的 [EGAI, 10.6.3--10.6.5]: 归结到 $X_0$ 为环 $A_0$ 的仿射概形的情形, 此时每个 $X_n$ 自动为环 $A_n$ 的仿射 Noether 概形, 而 $\mathcal X=\operatorname{Spf}A$, 其中 $A=\lim_n A_n$.

\paragraph{8.1.4.}
设 $\mathcal X$ 是局部 Noether 形式概形, $\mathcal I$ 为 $\mathcal X$ 的定义理想, 并考虑相应的加厚链 (8.1.3.2). 对 $m\leq n$, 以
\[u_{mn}:X_m\longrightarrow X_n,\qquad u_n:X_n\longrightarrow\mathcal X\]
记典范态射. 若 $E$ 是 $\mathcal X$ 上的凝聚模, 则 $E_n:=u_n^*E$ 是 $X_n$ 上的凝聚模, 这些模组成一个逆系, 转移映射 $E_n\to E_m$ 是 $\mathcal O_{X_n}$ 线性的并诱导同构 $u_{mn}^*E_n\xrightarrow\sim E_m$, 而 $E=\lim_n E_n$. 反过来, 设 $F_\bullet=(F_n,f_{mn}:F_n\to F_m)$ 是 $\mathcal O_{X_n}$ 模的逆系, 转移映射 $f_{mn}$ 是 $\mathcal O_{X_n}$ 线性的. 若每个 $F_n$ 都凝聚, 且转移映射 $f_{mn}$ 诱导同构 $u_{mn}^*F_n\xrightarrow\sim F_m$, 则称 $F_\bullet$ \emph{凝聚}. 若 $F_\bullet$ 凝聚, $F:=\lim_nF_n$ 是相应的 $\mathcal O_{\mathcal X}$ 模, 则 $F$ 凝聚, 且 $F_\bullet$ 典范同构于逆系 $(u_n^*F)$. 从 $\mathcal X$ 上的凝聚层范畴到凝聚逆系 $(F_n)$ 的范畴 $\operatorname{Coh}(X_\bullet)$ 的函子
\begin{equation}\operatorname{Coh}(\mathcal X)\longrightarrow\operatorname{Coh}(X_\bullet),\qquad E\longmapsto(u_n^*E)\tag{8.1.4.1}\label{ch08:eq8141}\end{equation}
是一个等价. 对如上的 $E=\lim_nE_n\in\operatorname{Coh}(\mathcal X)$, $E$ 的支集是闭的 (因为 $E$ 凝聚), 并且与 $E_0$ 的支集一致. 由平坦性判别法 [Bou61, III, §5, th.~1] (的一个特例), $E$ 平坦 (等价地, 为有限型局部自由模) 当且仅当对所有 $n$, $E_n$ 都为有限型局部自由模.

\paragraph{8.1.5.}
设 $f:\mathcal X\to\mathcal Y$ 是局部 Noether 形式概形的态射, $\mathcal J$ 为 $\mathcal Y$ 的一个定义理想. 因为 $\mathcal J\subset\mathcal T_{\mathcal Y}$, $f$ 的连续性蕴含理想 $f^*(\mathcal J)\mathcal O_{\mathcal X}$ 包含在 $\mathcal T_{\mathcal X}$ 中 (8.1.3.1). 固定一个满足 $f^*(\mathcal J)\mathcal O_{\mathcal X}\subset\mathcal I$ 的定义理想 $\mathcal I$ (例如 $\mathcal I=\mathcal T_{\mathcal X}$), 并如 (8.1.3.2) 那样考虑分别由 $\mathcal I$ 和 $\mathcal J$ 定义的归纳系 $X_\bullet,Y_\bullet$. 于是, 因为 $f^*(\mathcal J^{n+1})\mathcal O_{\mathcal X}\subset\mathcal I^{n+1}$, $f$ 诱导一个归纳系态射
\begin{equation}f_\bullet:X_\bullet\longrightarrow Y_\bullet,\tag{8.1.5.1}\label{ch08:eq8151}\end{equation}
% Source PDF 197 / printed p.185.
即一些使下列方块交换的概形态射 $f_n:X_n\to Y_n$:
\begin{equation}
\begin{tikzcd}X_m\arrow[r]\arrow[d,"f_m"']&X_n\arrow[d,"f_n"]\\Y_m\arrow[r]&Y_n\end{tikzcd}
\tag{8.1.5.2}\label{ch08:eq8152}\end{equation}
而 $f$ 是各态射 $f_n$ 的余极限, 其特征在于使下列方块交换:
\begin{equation}
\begin{tikzcd}X_n\arrow[r,"u_n"]\arrow[d,"f_n"']&\mathcal X\arrow[d,"f"]\\Y_n\arrow[r,"u_n"']&\mathcal Y\end{tikzcd}
\tag{8.1.5.3}\label{ch08:eq8153}\end{equation}
容易验证 [EGAI, 10.6.8], $f\mapsto f_\bullet$ 给出了从满足 $f^*(\mathcal J)\mathcal O_{\mathcal X}\subset\mathcal I$ 的 $\mathcal X$ 到 $\mathcal Y$ 的态射集合, 到形如 (8.1.5.1) 的归纳系态射集合的双射.

一般说来, $f^*(\mathcal J)\mathcal O_{\mathcal X}$ 并不是 $\mathcal X$ 的定义理想. 若它是定义理想, 则称 $f$ 为\emph{进态射 (adic morphism)} (并称 $\mathcal X$ 为 $\mathcal Y$ \emph{进形式概形}). 此时可取 $\mathcal I=f^*(\mathcal J)\mathcal O_{\mathcal X}$, 而方块 (8.1.5.2) 均为笛卡尔方块. 反过来, 任何使方块 (8.1.5.2) 为笛卡尔方块的归纳系态射 (8.1.5.1), 都定义一个从 $\mathcal X$ 到 $\mathcal Y$ 的进态射.

设 $f:\mathcal X\to\mathcal Y$ 为进态射, $E$ 是 $\mathcal X$ 上的凝聚层. 则下列条件等价:
\begin{enumerate}
\item[(i)] $E$ 在 $\mathcal Y$ 上\emph{平坦} (或称 $\mathcal Y$ \emph{平坦}); 即对 $\mathcal X$ 的每一点 $x$, 茎 $E_x$ 在 $\mathcal O_{\mathcal Y,f(x)}$ 上平坦;
\item[(ii)] 用 (8.1.5.3) 中的记号, 对所有 $n\geq0$, $E_n=u_n^*E$ 都是 $Y_n$ 平坦的;
\item[(iii)] $E_0$ 是 $Y_0$ 平坦的, 且自然的 (满射) 映射
\[\operatorname{gr}^n\mathcal O_{\mathcal Y}\otimes_{\operatorname{gr}^0\mathcal O_{\mathcal Y}}\operatorname{gr}^0E\longrightarrow\operatorname{gr}^nE,\]
对所有 $n\geq0$ 都是同构, 其中伴随分次 $\operatorname{gr}$ 是相对于 $\mathcal J$ 进滤过取的.
\end{enumerate}
这是平坦性判别法 [Bou61, III, §5, th.~2, prop.~2] 的推论. 当 $E=\mathcal O_{\mathcal X}$ 满足上述等价条件时, 就称 $f$ \emph{平坦}.

\paragraph{8.1.6.}
设 $X$ 为局部 Noether \emph{概形}, $X'$ 为 $X$ (的底空间) 的一个闭子集. 选取 $\mathcal O_X$ 的一个凝聚理想 $I$, 使 $I$ 定义的 $X$ 的闭子概形以 $X'$ 为底空间. 这样的理想存在, 实际上存在最大的一个, 由在 $X'$ 上为零的 $\mathcal O_X$ 的局部截面组成; 对这个理想, $X'$ 具有约化概形结构. 考虑以 $X'$ 为共同底空间的 (局部 Noether) 概形归纳系
\[X_0\longrightarrow X_1\longrightarrow\cdots\longrightarrow X_n\longrightarrow\cdots,\]
其中 $X_n$ 是由 $I^{n+1}$ 定义的 $X$ 的闭子概形. 它满足 8.1.3 中的条件 (i)--(iv), 因而其余极限
\begin{equation}X_{/X'}:=\operatorname{colim}_n X_n\tag{8.1.6.1}\label{ch08:eq8161}\end{equation}
是一个以 $X'$ 为底空间的局部 Noether 形式概形, 称为 $X$ 沿 $X'$ 的\emph{形式完备化 (formal completion)}. 在不致混淆时, 有时也简记为 $\widehat X$. 容易验证, $X_{/X'}$ 不依赖于理想 $I$ 的选择.
% Source PDF 198 / printed p.186.
事实上, $\mathcal O_{\widehat X}$ 是各环 $\mathcal O_X/J$ 的逆极限, 其中 $J$ 遍历 $\mathcal O_X$ 中使 $\mathcal O_X/J$ 的支集为 $X'$ 的所有凝聚理想 (在 $X$ 的任何 Noether 开子集上, $I$ 的各次幂组成一个共尾系). 若 $X$ 仿射, $X=\Spec A$ 且 $I=\widetilde J$, 则 $\widehat X=\operatorname{Spf}\widehat A$, 其中 $\widehat A$ 是 $A$ 的 $J$ 进完备化.

闭浸入 $i_n:X_n\to X$ 定义一个赋环空间态射
\begin{equation}i_X:\widehat X\longrightarrow X\tag{8.1.6.2}\label{ch08:eq8162}\end{equation}
(或记为 $i$), 它是平坦的, 且对 $X$ 上的任何凝聚层 $F$, 自然映射
\begin{equation}i^*F\longrightarrow F_{/X'}:=\lim_n i_n^*F\tag{8.1.6.3}\label{ch08:eq8163}\end{equation}
是一个同构. 当 $X=\Spec A$, $F=\widetilde M$, $M$ 是有限型 $A$ 模时, 有 $F_{/X'}=\widehat M^\triangle$ (8.1.2), 其中 $\widehat M$ 是 $A$ 的 $I$ 进完备化 $\widehat A$ 上的模 $\lim_n(M/I^{n+1}M)$. 上述断言来自 Krull 定理: 若 $A$ 是 Noether 环, $J$ 是 $A$ 的一个理想, 则 $J$ 进完备化 $\widehat A$ 在 $A$ 上平坦, 而对任何有限型 $A$ 模 $M$, 有 $\widehat M=M\otimes\widehat A$. 在不致混淆时, 有时把 $F_{/X'}$ 写成 $\widehat F$. 注意, 若 $F$ 不凝聚, 则 (8.1.6.3) 一般不是同构. 类似地可以验证, 伴随映射
\begin{equation}F\longrightarrow i_*i^*F\tag{8.1.6.4}\label{ch08:eq8164}\end{equation}
的核由那些在 $X'$ 的某个邻域中为零的 $F$ 的截面组成.

设 $f:X\to Y$ 是局部 Noether 概形的一个态射, $X'$ (相应地, $Y'$) 为 $X$ (相应地, $Y$) 的闭子集, 满足 $f(X')\subset Y'$. 选取凝聚理想 $J\subset\mathcal O_X$, $K\subset\mathcal O_Y$, 分别定义以 $X'$ 与 $Y'$ 为底空间的闭子概形, 并满足 $f^*(K)\mathcal O_X\subset J$ (例如, 可取 $K$ 为任何定义以 $Y'$ 为底空间的闭子概形的理想, 并取 $J$ 为在 $X'$ 上为零的 $\mathcal O_X$ 的截面所组成的理想). 于是 $f$ 诱导归纳系态射
\[f_\bullet:X_\bullet\longrightarrow Y_\bullet,\]
其中 $X_n$ (相应地, $Y_n$) 如上定义. 由 8.1.5 中说明的对应, $f_\bullet$ 给出态射
\begin{equation}\widehat f:X_{/X'}\longrightarrow Y_{/Y'},\tag{8.1.6.5}\label{ch08:eq8165}\end{equation}
它不依赖于 $J,K$ 的选择, 称为 $f$ 到完备化 $X_{/X'}$ 和 $Y_{/Y'}$ 的\emph{延拓}. 这一态射包含在交换方块
\begin{equation}
\begin{tikzcd}X_{/X'}\arrow[r,"i"]\arrow[d,"\widehat f"']&X\arrow[d,"f"]\\Y_{/Y'}\arrow[r,"i"']&Y\end{tikzcd}
\tag{8.1.6.6}\label{ch08:eq8166}\end{equation}
中, 其中水平映射为典范态射 (8.1.6.2). 当 $X'=f^{-1}(Y')$ 时, 可取 $J=f^*(K)\mathcal O_X$, 所有方块
\[\begin{tikzcd}X_n\arrow[r]\arrow[d,"f_n"']&X\arrow[d,"f"]\\Y_n\arrow[r]&Y\end{tikzcd}\]
% Source PDF 199 / printed p.187, first paragraph completes §8.1.
都是笛卡尔方块, 因而方块 (8.1.5.2) 也一样, 所以 $\widehat f$ 是进态射.

\section{比较定理}\label{ch08:sec82}
\paragraph{8.2.1.}
设 $f:X\to Y$ 是局部 Noether 概形的一个态射, $Y'$ 是 $Y$ 的闭子集, $X'=f^{-1}(Y')$. 记 $\widehat X=X_{/X'}$, $\widehat Y=Y_{/Y'}$. 若 $F$ 是 $\mathcal O_X$ 模, 则方块 (8.1.6.6) 定义基变换映射 (见 8.2.19)
\begin{equation}i^*R^qf_*F\longrightarrow R^q\widehat f_*(i^*F)\tag{8.2.1.1}\label{ch08:eq8211}\end{equation}
(对所有 $q\in\mathbb Z$), 它们是 $\mathcal O_{\widehat Y}$ 模的映射. 因为函子 $i_*$ 正合, 8.2.19 中给出的一般定义在这里归结如下: 与 (8.2.1.1) 伴随的映射
\[R^qf_*F\longrightarrow i_*R^q\widehat f_*(i^*F)=R^qf_*(i_*i^*F)\]
是对伴随映射 $F\to i_*i^*F$ 施行 $R^qf_*$ 所得到的映射. 若 $F$ 凝聚, 则由 (8.1.6.3) 可将 $i^*F$ 与 $\widehat F=F_{/X'}$ 等同; 类似地, 若 $R^qf_*F$ 凝聚, 可将 $i^*R^qf_*F$ 与 $(R^qf_*F)_{/Y'}$ 等同: 当 $F$ 凝聚而 $f$ 固有时, 由固有态射的有限性定理 [EGAIII1, 3.2.1, 3.2.4] 可知这一点成立 (更一般地, $f$ 有限型且 $F$ 的支集在 $Y$ 上固有也足够; 后一条件意指存在 $X$ 的一个在 $Y$ 上固有的闭子概形 $Z$, 其底空间就是 $F$ 的支集 [EGAII, 5.4.10]). 此时 (8.2.1.1) 可以重写为
\begin{equation}(R^qf_*F)^\wedge\longrightarrow R^q\widehat f_*\widehat F.\tag{8.2.1.2}\label{ch08:eq8212}\end{equation}
另一方面, 若 $F$ 凝聚, 则取 $\mathcal X=\widehat X$, $\mathcal Y=\widehat Y$ 时的方块 (8.1.5.3) 定义 $\mathcal O_{Y_n}$ 线性的基变换映射
\[u_n^*R^q\widehat f_*\widehat F\longrightarrow R^q(f_n)_*F_n,\]
其中 $F_n=u_n^*\widehat F=i_n^*F$ (用 (8.1.6.2) 中的记号). 由伴随, 可将这些映射视为 $\mathcal O_{\widehat Y}$ 线性映射
\[R^q\widehat f_*\widehat F\longrightarrow R^q(f_n)_*F_n\]
(它们由对伴随映射 $\widehat F\to(u_n)_*F_n$ 施行 $R^q\widehat f_*$ 得到), 因而定义 $\mathcal O_{\widehat Y}$ 线性映射
\begin{equation}R^q\widehat f_*\widehat F\longrightarrow\lim_n R^q(f_n)_*F_n.\tag{8.2.1.3}\label{ch08:eq8213}\end{equation}
注意, 更一般地, 对 $F\in D^+(X,\mathcal O_X)$ 也定义了基变换映射 (8.2.1.1), 它由 $D^+(\widehat Y,\mathcal O_{\widehat Y})$ 中的基变换映射
\begin{equation}i^*Rf_*F\longrightarrow R\widehat f_*i^*F\tag{8.2.1.4}\label{ch08:eq8214}\end{equation}
在各层 $\mathcal H^q$ 上诱导.

\begin{origstatement}{定理}{8.2.2}\label{ch08:th822}
设 $f:X\to Y$ 是 Noether 概形的有限型态射, $Y'$ 是 $Y$ 的闭子集, $X'=f^{-1}(Y')$, $\widehat f:\widehat X\to\widehat Y$ 是 $f$ 到 $X$ 和 $Y$ 沿 $X'$ 与 $Y'$ 的形式完备化的延拓. 设 $F$ 是 $X$ 上的凝聚层, 其支集在 $Y$ 上固有. 则对所有 $q$, 典范映射 (8.2.1.2)、(8.2.1.3) 都是拓扑同构.
\end{origstatement}

\begin{origstatement}{注}{8.2.3}\label{ch08:rem823}
\begin{enumerate}
\item[(a)] 在 8.2.2 对 $f$ 的假设下, 对任何 $F\in D^+(X,\mathcal O_X)$, 若所有 $\mathcal H^iF$ 都凝聚且其支集在 $Y$ 上固有, 则 (8.2.1.4) 是同构.
% Source PDF 200 / printed p.188.
利用 $X$ 上凝聚层的有界导出范畴 $D^b(\operatorname{Coh}(X))$ 到 $D^b(X):=D^b(X,\mathcal O_X)$ 中由具有凝聚上同调的复形组成的满子范畴 $D^b(X)_{\mathrm{coh}}$ 的自然函子是一个等价 [SGA6, II 2.2.2.1], 可将 8.2.2 中的同构 (8.2.1.3) 推广到 $F\in D^b(X)_{\mathrm{coh}}$ 的情形. 我们略去细节.
\item[(b)] 考虑 $X$ 的一个以 $F$ 的支集为底空间的闭子概形 $Z$, 可将 8.2.2 归结到 $f$ 固有的情形.
\item[(c)] Grothendieck 的原始证明没有发表. 从 [Gro95a, p.~05] 可以推测, 它包含两步: (i) 在 $f$ 为射影态射时, 对 $q$ 作下降归纳 (当 $Y'$ 是一点时, 见 [Har77, III 11.1]); (ii) 在一般情形, 借助 Chow 引理与 Noether 归纳归结到射影情形. [EGAIII1, 4.1.7, 4.1.8] 中给出的证明采用了 Serre 的一个论证.
\item[(d)] 容易看出, 8.2.2 实际等价于下面的特殊情形.
\end{enumerate}
\end{origstatement}

\begin{origstatement}{推论}{8.2.4}\label{ch08:cor824}
在 8.2.2 的假设下, 设 $Y=\Spec A$, 其中 $A$ 是 Noether 环, $I$ 是 $A$ 的一个理想, 满足 $\operatorname{Supp}(\mathcal O_Y/\mathcal I)=Y'$, 这里 $\mathcal I=\widetilde I$. 令 $Y_n=\Spec(A/I^{n+1})$, $X_n=Y_n\times_YX$, $F_n=i_n^*F=F/I^{n+1}F$. 则对所有 $q$, 由 (8.2.1.2) 与 (8.2.1.3) 的复合定义的自然映射
\begin{equation}\varphi_q:H^q(X,F)^\wedge\longrightarrow\lim_nH^q(X,F_n),\tag{8.2.4.1}\label{ch08:eq8241}\end{equation}
以及由 (8.2.1.3) 定义的自然映射
\begin{equation}\psi_q:H^q(\widehat X,\widehat F)\longrightarrow\lim_nH^q(X,F_n),\tag{8.2.4.2}\label{ch08:eq8242}\end{equation}
都是拓扑同构.
\end{origstatement}

8.2.4 的证明 [EGAIII1, 4.1.7] 使用两个要素, 第一个是标准而初等的交换代数与同调代数, 第二个则深刻得多: (a) Artin--Rees 引理与 Mittag-Leffler 条件; (b) 固有态射的有限性定理 [EGAIII1, 3.2], 特别是其一个分次变体 [EGAIII1, 3.3.2].

我们将简要回顾 (a) 与 (b), 然后给出 8.2.4 证明的要点.

\paragraph{8.2.5. Artin--Rees 与 Mittag-Leffler.}
\paragraph{8.2.5.1.}
设 $A$ 是 Noether 环, $I$ 是 $A$ 的一个理想, $M$ 是一个有限生成 $A$ 模, 赋有由子模组成的下降滤过 $(M_n)_{n\in\mathbb Z}$. 若该滤过\emph{穷尽} (即存在 $n_1$ 使 $M_{n_1}=M$) 并满足下列两个条件, 则称 $(M_n)$ 为 $I$ \emph{良滤过}:
\begin{enumerate}
\item[(i)] 对所有 $n\in\mathbb Z$, $IM_n\subset M_{n+1}$ (意即由 $(M_n)$ 滤过的 $M$ 是由 $I$ 进滤过滤过的环 $A$ 上的滤过模);
\item[(ii)] 存在整数 $n_0$, 使对所有 $n\geq n_0$, $M_{n+1}=IM_n$.
\end{enumerate}
例如, 由 $n\leq0$ 时 $M_n=M$, $n\geq0$ 时 $M_n=I^nM$ 定义的 $M$ 的 $I$ 进滤过就是 $I$ 良滤过. 所有 $I$ 良滤过在 $M$ 上定义同一个拓扑, 即 $I$ 进拓扑.

假定条件 (i) 成立. 考虑分次环
\[A':=\bigoplus_{n\in\mathbb N}I^n,\]
% Source PDF 201 / printed p.189.
有时也写作 $\bigoplus I^nt^n$, 其中 $t$ 是不定元, 以明确表示 $I^n=I^nt^n$ 是 $A'$ 的第 $n$ 个分量; 再考虑 $A'$ 上的分次模
\[M'=\bigoplus_{n\in\mathbb N}M_n,\]
它有时也写作 $\bigoplus_{n\in\mathbb N}M_nt^n$. 一个基本观察 [Bou61, III, §3, th.~1] 是, 条件 (ii) 等价于
\begin{enumerate}\item[(ii$'$)] $M'$ 是有限生成 $A'$ 模.\end{enumerate}
这一观察的证明是直接的.

因为 $A'$ 是 Noether 的, 这立即蕴含经典的 \emph{Artin--Rees 定理}: 若 $N$ 是 $M$ 的子模, 则 $M$ 的 $I$ 进滤过在 $N$ 上诱导的滤过是 $I$ 良滤过; 换言之, 存在 $n_0\geq0$, 使对所有 $n\geq n_0$,
\[I^nM\cap N=I^{n-n_0}(I^{n_0}M\cap N).\]
(ii$'$) 蕴含 (ii) 是下面性质 [EGAII, 2.1.6] 的一个特例 (其证明同样是直接的):
\begin{enumerate}\item[(iii)] 设 $A$ 为交换环, $S=\bigoplus_{n\in\mathbb N}S_n$ 为 $S_0$ 上有限型且由 $S_1$ 生成的分次 $A$ 代数, $M=\bigoplus_{n\in\mathbb Z}M_n$ 为有限型分次 $S$ 模. 则存在 $n_0\in\mathbb N$, 使对所有 $n\geq n_0$, $M_{n+1}=S_1M_n$.\end{enumerate}

\paragraph{8.2.5.2.}
设 $A$ 是交换环, $M_\bullet=(M_n,u_{mn}:M_n\to M_m)$ 是由 $\mathbb N$ 编指标的 $A$ 模射影系. 我们作如下定义:
\begin{enumerate}
\item[(i)] 若转移映射 $u_{mn}$ 满射, 称 $M_\bullet$ 为\emph{严格的};
\item[(ii)] 若对每个 $m$, 存在 $n\geq m$ 使 $u_{mn}=0$, 即 $M_\bullet$ 定义的 pro 对象为零, 称 $M_\bullet$ 为\emph{本质为零的};
\item[(iii)] 若对每个 $m$, 存在 $n\geq m$, 使对所有 $n'\geq n$, 在 $M_m$ 中有 $\operatorname{Im}u_{mn'}=\operatorname{Im}u_{mn}$, 则称 $M_\bullet$ 满足 \emph{Mittag-Leffler 条件} (简称 ML).
\end{enumerate}
有时考虑下列更强的条件是有用的:
\begin{enumerate}
\item[(ii$'$)] 若存在整数 $r\geq0$, 使对所有 $n$, $u_{n,n+r}=0$, 则称 $M_\bullet$ 为 \emph{Artin--Rees 零} (简称 AR 零);
\item[(iii$'$)] 若存在整数 $r\geq0$, 使对所有 $m$ 和所有 $n'\geq m+r$, 有 $\operatorname{Im}u_{mn'}=\operatorname{Im}u_{m,m+r}$, 则称 $M_\bullet$ 满足 \emph{Artin--Rees--Mittag-Leffler 条件} (简称 ARML).
\end{enumerate}
关于 Mittag-Leffler 条件的讨论, 参见 [EGAIII1, Ch.~0, 13]. 这里只回顾两个基本的 (容易的) 事实:
\begin{enumerate}
\item[(a)] 若 $M_\bullet$ 本质为零, 则 $\lim_nM_n=0$.
\item[(b)] 函子 $M_\bullet\mapsto\lim_nM_n$ 左正合. 此外, 设
\[0\longrightarrow L_\bullet\longrightarrow M_\bullet\longrightarrow N_\bullet\longrightarrow0\]
是 $A$ 模逆系的一个正合列; 若 $L_\bullet$ 满足 ML, 则序列
\[0\longrightarrow\lim_nL_n\longrightarrow\lim_nM_n\longrightarrow\lim_nN_n\longrightarrow0\]
正合.
\end{enumerate}
更强的条件 (iii$'$) (有时称为\emph{一致 Mittag-Leffler 条件}) 与 Artin--Rees 定理有密切关系. 相关讨论见 [SGA5, V]. 术语 AR 零、ARML 即取自该处.

我们将用到一个关于 $H^q(X,-)$ 与逆极限交换的一般结果 [EGAIII1, Ch.~0, 13.3.1] 的一个 (很) 特殊的情形; 其证明是初等的:
% Source PDF 202 / printed p.190.
\begin{origstatement}{命题}{8.2.5.3}\label{ch08:prop8253}
设 $X$ 是概形, $(F_n)_{n\in\mathbb N}$ 是 $X$ 上拟凝聚层的一个逆系, 其转移映射满射. 假定对所有 $i\in\mathbb Z$, (作为 $\mathbb Z$ 模的) 逆系 $H^i(X,F_\bullet)$ 满足 ML. 则对所有 $i\in\mathbb Z$, 自然映射
\[H^i(X,\lim_nF_n)\longrightarrow\lim_nH^i(X,F_n)\]
是同构.
\end{origstatement}

\paragraph{8.2.6. 有限性定理.}
固有态射的基本有限性定理 [EGAIII1, 3.2.1] 断言: 若 $f:X\to Y$ 是固有态射, $Y$ 局部 Noether, $F$ 是 $X$ 上的凝聚层, 则对所有 $q\in\mathbb Z$, $Y$ 上的层 $R^qf_*F$ 都凝聚. 我们需要下面的变体 [EGAIII1, 3.3.1]:

\begin{origstatement}{定理}{8.2.6.1}\label{ch08:th8261}
设 $f:X\to Y$ 是固有态射, $Y$ 是 Noether 的. 设 $S=\bigoplus_{n\in\mathbb N}S_n$ 是在 $S_0$ 上有限型且由 $S_1$ 生成的拟凝聚分次 $\mathcal O_Y$ 代数. 设 $M=\bigoplus_{n\in\mathbb Z}M_n$ 是有限型拟凝聚分次 $f^*(S)$ 模. 则对所有 $q\in\mathbb Z$,
\[R^qf_*M:=\bigoplus_{n\in\mathbb Z}R^qf_*M_n\]
是有限型分次 $S$ 模, 并且存在整数 $n_0$, 使对任何 $n\geq n_0$,
\[R^qf_*M_n=S_{n-n_0}R^qf_*M_{n_0}.\]
\end{origstatement}
这里 $R^qf_*M$ 上的分次 $S$ 模结构来自乘法映射, 即下列复合:
\[\begin{aligned}
S_k\otimes R^qf_*M_n&\longrightarrow f_*f^*S_k\otimes R^qf_*M_n\\
&\longrightarrow R^qf_*((f^*S_k)\otimes M_n)\longrightarrow R^qf_*M_{n+k}.
\end{aligned}\]
8.2.6.1 的最后一个断言是第一个断言的推论 (利用 8.2.5.1 (iii)); 而第一个断言则来自将有限性定理应用于由笛卡尔方块
\[\begin{tikzcd}X\arrow[d,"f"']&\widetilde X\arrow[l]\arrow[d,"\widetilde f"]\\Y&\widetilde Y\arrow[l]\end{tikzcd}\]
定义的 (固有) 态射 $\widetilde f:\widetilde X\to\widetilde Y$, 以及 $\widetilde X$ 上的凝聚模 $\widetilde M$. 这里 $\widetilde Y=\Spec S$, $\widetilde X=\Spec f^*(S)$.

\begin{origstatement}{推论}{8.2.6.2}\label{ch08:cor8262}
在 8.2.4 的假设下, 令 $B:=\bigoplus_{n\in\mathbb N}I^n$. 则对所有 $q$, $\bigoplus_{n\in\mathbb N}H^q(X,I^nF)$ 是有限生成分次 $B$ 模, 并且存在 $n_0\geq0$, 使对所有 $n\geq n_0$, 有 $H^q(X,I^nF)=I^{n-n_0}H^q(X,I^{n_0}F)$.
\end{origstatement}

\paragraph{8.2.7. 8.2.4 的证明.}
与 Grothendieck 的原始证明不同, [EGAIII1, 4.1.7] 中给出的证明不对 $q$ 作下降归纳. 整个证明中整数 $q$ 都固定, 证明由对映射逆系
\begin{equation}H^q(F)\longrightarrow H^q(F_n),\tag{8.2.7.1}\label{ch08:eq8271}\end{equation}
% Source PDF 203 / printed p.191.
的细致分析组成, 其中为简洁起见写 $H^q=H^q(X,-)$. 映射 (8.2.7.1) 包含在短正合列
\[0\longrightarrow I^{n+1}F\longrightarrow F\longrightarrow F_n\longrightarrow0\]
所对应的上同调长正合列的一段中, 即
\begin{equation}\begin{aligned}
H^q(I^{n+1}F)\longrightarrow H^q(F)\longrightarrow H^q(F_n)\\
\longrightarrow H^{q+1}(I^{n+1}F)\longrightarrow H^{q+1}(F).
\end{aligned}\tag{8.2.7.2}\label{ch08:eq8272}\end{equation}
由 (8.2.7.2) 得到正合列
\begin{equation}0\longrightarrow R_n\longrightarrow H^q(F)\longrightarrow H^q(F_n)\longrightarrow Q_n\longrightarrow0,\tag{8.2.7.3}\label{ch08:eq8273}\end{equation}
其中
\[R_n=\operatorname{Im}(H^q(I^{n+1}F)\longrightarrow H^q(F)),\]
而
\[\begin{aligned}
Q_n&=\operatorname{Im}(H^q(F_n)\longrightarrow H^{q+1}(I^{n+1}F))\\
&=\ker(H^{q+1}(I^{n+1}F)\longrightarrow H^{q+1}(F)).
\end{aligned}\]
关键点如下:
\begin{enumerate}
\item[(1)] $H^q(F)$ 上的滤过 $(R_n)$ 为 $I$ 良滤过 (8.2.5.1); 特别地, $(R_n)$ 在 $H^q(F)$ 上定义的拓扑为 $I$ 进拓扑.
\item[(2)] 逆系 $Q_\bullet=(Q_n)$ 为 AR 零 (8.2.5.2 (ii$'$)).
\item[(3)] 逆系 $H^q(F_\bullet)=(H^q(F_n))$ 满足 ARML (8.2.5.2 (iii$'$)).
\end{enumerate}
先说明 (1)、(2)、(3) 蕴含 8.2.4. 考虑 (8.2.7.3) 定义的逆系正合列:
\begin{equation}0\longrightarrow H^q(F)/R_n\longrightarrow H^q(F_n)\longrightarrow Q_n\longrightarrow0.\tag{*}\label{ch08:eq827star}\end{equation}
由 (2), 有 $\lim_nQ_n=0$ (8.2.5.2 (a)). 因而由函子 $\lim_n$ 的左正合性得到同构
\begin{equation}\lim_nH^q(F)/R_n\xrightarrow\sim\lim_nH^q(F_n).\tag{**}\label{ch08:eq827stars}\end{equation}
由 (1), 从满射 $H^q(F)/I^{n+1}H^q(F)\to H^q(F)/R_n$ 得到的映射
\begin{equation}\begin{aligned}
H^q(F)^\wedge=\lim_nH^q(F)/I^{n+1}H^q(F)\longrightarrow\lim_nH^q(F)/R_n
\end{aligned}\tag{***}\label{ch08:eq827star3}\end{equation}
是同构. 把 (**) 与 (***) 合在一起, 就得到 (8.2.4.1) 为同构. 由定义,
\[\begin{aligned}
H^q(\widehat X,\widehat F)&=H^q(\widehat X,\lim F_n)\\
&=H^q(\widehat X,i_*\lim F_n)=H^q(X,\lim(i_n)_*F_n).
\end{aligned}\]
% Original printed p.191 has \widehat X in the third term; retained, review note records apparent error.
由 (3), 8.2.5.3 的假设成立, 因而 (8.2.4.2) 为同构.

还需证明 (1)、(2)、(3).

\emph{(1) 的证明.} 有 $R_{-1}=H^q(F)$. 包含关系
\[I^mR_n\subset R_{m+n}\]
来自自然映射
\[\bigoplus_{n\in\mathbb N}H^q(I^{n+1}F)t^n\longrightarrow\bigoplus_{n\in\mathbb N}H^q(F)t^n\]
是分次 $B$ 模映射这一事实, 其中 $B=\bigoplus_{n\in\mathbb N}I^nt^n$ (8.2.6.2). 对 $IF$ 应用 8.2.6.2, 可知 $\bigoplus_{n\in\mathbb N}H^q(I^{n+1}F)t^n$ 在 $B$ 上有限型, 故其商 $R:=\bigoplus_nR_n$ 也是如此. 利用 8.2.5.1 中条件 (ii) 与 (ii$'$) 的等价性, 即得 (1).

% Source PDF 204 / printed p.192.
\emph{(2) 的证明.} 这是最细致的一点. 再用 8.2.6.2, 模 $N:=\bigoplus_nH^{q+1}(I^{n+1}F)$ 在 $B$ 上有限生成. 因为 $B$ 是 Noether 的, $N$ 的 (分次) $B$ 子模 $Q:=\bigoplus_nQ_n$ 也有限生成, 因而存在 $r\geq0$, 使对所有 $n\geq r$, 有 $Q_{n+1}=IQ_n$. 因为 $Q_k$ 作为 $H^q(F_k)$ 的商被 $I^{k+1}$ 消灭 (作为 $A$ 模), 所以每个 $Q_n$ 都被 $I^{r+1}$ 消灭 (作为 $A$ 模). 现在, 对 $a\in I^p$, 从 $H^{q+1}(I^{n+1}F)$ 到 $H^{q+1}(I^{p+n+1}F)$ 的乘 $a$ 映射与从 $H^{q+1}(I^{p+n+1}F)$ 到 $H^{q+1}(I^{n+1}F)$ 的转移映射的复合, 就是 $H^{q+1}(I^{n+1}F)$ 中的乘 $a$ 映射. 因为对 $n\geq r$ 有 $Q_{n+r+1}=I^{r+1}Q_n$, 所以对所有 $n\geq r$, 转移映射 $Q_{n+r+1}\to Q_n$ 为零; 因而若取 $s=2r+1$, 则对所有 $n$, 转移映射 $Q_{n+s}\to Q_n$ 都为零.

\emph{(3) 的证明.} 这是 (2) 的形式推论. 正合列 (*) 的左项具有满射转移映射 (故显然满足 ARML), 右项为 AR 零, 因此它们都显然满足 ARML. 由下面引理 [SGA5, V 2.1.2] 的第二个断言, 中间项也满足 ARML; 引理的证明是初等的:

\begin{origstatement}{引理}{8.2.7.4}\label{ch08:lem8274}
设
\[0\longrightarrow L'_\bullet\longrightarrow L_\bullet\longrightarrow L''_\bullet\longrightarrow0\]
为 $A$ 模逆系的正合列. 若 $L_\bullet$ 满足 ARML, 则 $L''_\bullet$ 也满足; 若 $L'_\bullet$ 与 $L''_\bullet$ 都满足 ARML, 则 $L_\bullet$ 也满足.
\end{origstatement}

\begin{origstatement}{注}{8.2.8}\label{ch08:rem828}
\begin{enumerate}
\item[(a)] 8.2.7 中的性质 (2) 是 Deligne 为 Noether 概形的一个可紧化态射 $f:X\to S$ (即形如 $gi$, 其中 $g$ 固有而 $i$ 为开浸入) 构造从 $\operatorname{pro}D^b(X)_{\mathrm{coh}}$ 到 $\operatorname{pro}D^b(S)_{\mathrm{coh}}$ 的 $Rf_!$ 函子的关键技术点 [Har66, Appendix, Prop.~5]. 更精确地说, 在 8.2.4 的情形下, 若假定 $f$ 诱导从 $X-X'$ 到 $Y-Y'$ 的同构, 则该处证明的是对 $q>0$, 逆系 $H^q(I^{n+1}F)$ 为 AR 零; 这是因为对 $q>0$, $H^q(F)$ 被 $I$ 的某次幂消灭.
\item[(b)] 8.2.2 的证明表明, 若 $f:X\to Y$, $Y'$, $X'$ 如 8.2.2 中所述, $F$ 是 $X$ 上的凝聚层, 且对某个整数 $q$, 分次 $\mathcal O_Y$ 代数 $\bigoplus_nI^n$ 上的分次模 $\bigoplus_nR^qf_*(I^nF)$ 和 $\bigoplus_nR^{q+1}f_*(I^nF)$ 都有限生成, 则 $R^qf_*F$ 凝聚, 并且 (8.2.1.3) 与 (8.2.1.4) 为同构. 细节与例子见 [SGA2, IX]. 这一精细的比较定理是 Grothendieck 关于基本群与 Picard 群的 Lefschetz 型定理的关键工具 [SGA2, X, XI].
\end{enumerate}
\end{origstatement}
比较定理 8.2.2 有许多推论与应用. 我们只提其中几个. 下面的推论 (对 $r=0,1$) 是 §8.3 将讨论的 Grothendieck 存在定理证明的主要要素.
% Source cross-reference §8.3 retained; original existence theorem is in §8.4.

\begin{origstatement}{推论}{8.2.9}\label{ch08:cor829}
[EGAIII1, 4.5.1] 设 $A$ 是 Noether 环, $I$ 是 $A$ 的一个理想, $f:X\to Y$ 是有限型态射, $\widehat f:\widehat X\to\widehat Y$ 是如 8.2.4 中那样沿 $Y'=V(I)$ 与 $X'=f^{-1}(Y')$ 的完备化. 设 $F,G$ 是 $X$ 上的凝聚层, 其支集的交在 $Y$ 上固有. 则对所有 $r\in\mathbb Z$, $\operatorname{Ext}^r(F,G)$ 是有限型 $A$ 模, 而自然映射 $\operatorname{Ext}^r(F,G)\to\operatorname{Ext}^r(\widehat F,\widehat G)$ 诱导同构
\begin{equation}\operatorname{Ext}^r(F,G)^\wedge\xrightarrow\sim\operatorname{Ext}^r(\widehat F,\widehat G).\tag{8.2.9.1}\label{ch08:eq8291}\end{equation}
\end{origstatement}
% Source PDF 205 / printed p.193.
有
\[\operatorname{Ext}^r(F,G)=H^rR\Hom(F,G)=H^rR\Gamma(X,R\mathcal Hom(F,G)).\]
对 $F,G$ 的假设蕴含 $R\mathcal Hom(F,G)$ 的上同调层凝聚且其支集在 $Y$ 上固有. 因而由有限性定理, $R\Hom(F,G)=R\Gamma(X,R\mathcal Hom(F,G))$ 的上同调群在 $A$ 上有限生成, 并由注 8.2.3 (a), 基变换映射
\[R\Gamma(X,R\mathcal Hom(F,G))^\wedge\longrightarrow R\Gamma(\widehat X,R\mathcal Hom(F,G)^\wedge),\]
其中 $(-)^\wedge=i^*$, 是一个同构. 而且, 因为 $i$ 平坦,
\[R\mathcal Hom(F,G)^\wedge=R\mathcal Hom(\widehat F,\widehat G),\]
结论随之成立.

下面的推论在几何应用中十分有用:
\begin{origstatement}{推论}{8.2.10}\label{ch08:cor8210}
(形式函数定理, [EGAIII1, 4.2.1]) 设 $f:X\to Y$ 是局部 Noether 概形的固有态射, $y$ 是 $Y$ 的一点, $X_y=X\times_Y\Spec k(y)$ 是 $f$ 在 $y$ 处的纤维, $F$ 是 $X$ 上的凝聚层. 令 $F_n=F\otimes\mathcal O_y/\mathfrak m_y^{n+1}$, 视为 $X_n=X\times_Y\Spec\mathcal O_y/\mathfrak m_y^{n+1}$ 上的层. 则对所有 $q\in\mathbb Z$, 茎 $R^qf_*(F)_y$ 是有限型 $\mathcal O_y$ 模, 且自然映射
\begin{equation}\begin{aligned}
(R^qf_*(F)_y)^\wedge
&=\lim_n(R^qf_*(F)_y/\mathfrak m_y^{n+1}R^qf_*(F)_y)\\
&\longrightarrow\lim_nH^q(X_y,F_n)
\end{aligned}\tag{8.2.10.1}\label{ch08:eq82101}\end{equation}
是同构.
\end{origstatement}
映射 (8.2.10.1) 由基变换映射
\[(R^qf_*(F)_y/\mathfrak m_y^{n+1}R^qf_*(F)_y)\longrightarrow H^q(X_y,F_n)\]
定义, 其中右边把 $X_y$ 视为概形 $X_n$ 的底空间. 当 $y$ 为闭点时, 8.2.10 是 8.2.2 的特殊情形. 通过基变换 $\Spec\mathcal O_y\to Y$, 可归结到这一情形.

\paragraph{8.2.11.}
设 $f:X\to Y$ 是局部 Noether 概形的固有态射. 则 $f_*\mathcal O_X$ 是一个有限 $\mathcal O_Y$ 代数. 其谱 $Y'=\Spec f_*\mathcal O_X$ 是 $Y$ 上的一个有限概形, 而 $f_*\mathcal O_X$ 的恒等映射定义 $f$ 的分解
\[X\xrightarrow{f'}Y'\xrightarrow{g}Y,\]
其中 $f'$ 固有, $g$ 有限, 称为 $f$ 的 \emph{Stein 分解}. 其主要性质由下述定理描述:

\begin{origstatement}{定理}{8.2.12}\label{ch08:th8212}
(Zariski 连通性定理, [EGAIII1, 4.3.1]) 在 8.2.11 的假设和记号下, $f'_*\mathcal O_X=\mathcal O_{Y'}$, 且 $f'$ 的各纤维连通而非空.
\end{origstatement}
第一个断言直接来自定义. 对第二个断言, 先归结到 $Y'=Y$ 且 $y$ 是 $Y$ 的闭点的情形. 此时若 $\widehat X$ 为 $X$ 沿 $X_y$ 的完备化, 则由 8.2.10,
\[\widehat{\mathcal O}_{Y,y}=(f_*\mathcal O_X)_y^\wedge=H^0(X_y,\mathcal O_{\widehat X}),\]
它不可能是两个非零环的积.

特别地, 若 $Y'=Y$, 即 $f_*\mathcal O_X=\mathcal O_Y$, 则 $f$ 的各纤维连通而非空. 利用下文的基变换公式 8.3.3, 不难看出,
% Source PDF 206 / printed p.194.
它们实际上\emph{几何连通} (即连通且在任意域扩张后仍然连通) [EGAIII1, 4.3.4].

下面的推论是容易的, 细节见 [EGAIII1, 4.3, 4.4].

\begin{origstatement}{推论}{8.2.13}\label{ch08:cor8213}
在 8.2.12 的假设下, 对 $Y$ 的每一点 $y$, 纤维 $X_y$ 的连通分支与 $Y'_y$ 的点, 即有限 $\mathcal O_y$ 代数 $f_*(\mathcal O_X)_y$ 的极大理想, 一一对应.
\end{origstatement}
这是因为 $Y'_y$ 的底空间 $g^{-1}(y)$ 有限且离散.

\begin{origstatement}{推论}{8.2.14}\label{ch08:cor8214}
设 $f:X\to Y$ 是整 Noether 概形的固有满射态射, $Y$ 正规. 假定 $f$ 的一般纤维几何连通 (例如当 $f$ 双有理时即如此). 则 $f$ 的所有纤维都几何连通.
\end{origstatement}
设 $\zeta$ (相应地, $\eta$) 为 $X$ (相应地, $Y$) 的一般点 (于是 $f(\zeta)=\eta$). 关于一般纤维的假设意味着 $K=k(\eta)$ 在 $k(\zeta)$ 中的代数闭包 $K'$ 是 $K$ 的一个 (有限) 纯不可分扩张 [EGAIV2, 4.5.15]. 设 $y\in Y$. 因为 $\mathcal O_y$ 正规, 且 $K'$ 在 $K$ 上纯不可分, $\mathcal O_y$ 在 $K'$ 中的正规化 $A$ 是局部环, 而其剩余域扩张纯不可分 [Bou64, Ch.~5, §8.2, n\textsuperscript{o}3, Lemme 4]. 因为 $A$ 包含 $(f_*\mathcal O_X)_y$, 所以 $(f_*\mathcal O_X)_y$ 也具有同样性质. 因而由 8.2.12 (及其后的注), 纤维 $X_y$ 几何连通.

\begin{origstatement}{推论}{8.2.15}\label{ch08:cor8215}
在 8.2.12 的假设下, $X$ 的一点 $x$ 在其纤维中孤立, 即在 $f^{-1}(f(x))$ 中既开又闭, 当且仅当 $f'^{-1}(f'(x))=\{x\}$. 这些点组成的集合 $U$ 在 $X$ 中开, $U'=f'(U)$ 在 $Y'$ 中开, 而 $f':X\to Y'$ 诱导同构 $f'|_U:U\xrightarrow\sim f'(U)$.
\end{origstatement}
设 $y=f(x)$, $y'=f'(x)$. 因为 $g^{-1}(y)$ 有限且离散, $x$ 在 $f^{-1}(y)$ 中孤立当且仅当它在 $f'^{-1}(y')$ 中孤立. 因此可假定 $Y'=Y$, 即 $f_*\mathcal O_X=\mathcal O_Y$; 由 8.2.13, $x$ 在 $f^{-1}(y)$ 中孤立当且仅当 $f^{-1}(y)=\{x\}$. 对 $x\in U$, 分别选取 $x$ 和 $y$ 的仿射开邻域 $V=\Spec B$, $W=\Spec A$, 使 $f(V)\subset W$. 因为 $f$ 为闭映射, $f(X-V)$ 是 $Y$ 中不含 $y$ 的闭子集. 因而存在 $y$ 的一个形如 $W_s=\Spec A_s$ 的仿射开邻域, 其中 $s\in A$, 使 $f^{-1}(W_s)\subset V$. 此时 $f^{-1}(W_s)=V_s=\Spec B_s$. 因为 $f_*\mathcal O_X=\mathcal O_Y$, $f$ 在 $W_s$ 上诱导同构 $V_s\xrightarrow\sim W_s$.

\begin{origstatement}{推论}{8.2.16}\label{ch08:cor8216}
设 $f:X\to Y$ 是局部 Noether 概形的固有态射. 若 $f$ 拟有限 (即各纤维有限), 则 $f$ 有限.
\end{origstatement}

\begin{origstatement}{推论}{8.2.17}\label{ch08:cor8217}
(Zariski 主定理) 设 $f:X\to Y$ 是局部 Noether 概形的一个可紧化态射 (8.2.8 (a)) (例如一个拟射影态射 [EGAII, 5.3.2]). 若 $f$ 拟有限, 则 $f$ 可分解为 $f=gj$, 其中 $j:X\to Z$ 为开浸入, $g:Z\to Y$ 为有限态射.
\end{origstatement}
若 $Y$ 是 Noether 的, 则只要假定 $f$ 分离且有限型, 就可去掉 $f$ 可紧化的假设; 见 [EGAIV3, 8.12.6], 那里的证明不使用比较定理 8.2.2, 而依赖于更深刻的交换代数.

最后提及 8.2.13 的一个有用应用. 若 $X$ 为局部 Noether 概形, 用 $\pi_0(X)$ 表示其连通分支的集合.
% Source PDF 207 / printed p.195.
\begin{origstatement}{推论}{8.2.18}\label{ch08:cor8218}
设 $A$ 是 hensel Noether 局部环, $S=\Spec A$, $s$ 为其闭点, $X$ 是 $S$ 上的固有概形. 则自然映射
\[\pi_0(X_s)\longrightarrow\pi_0(X)\]
是双射.
\end{origstatement}
考虑结构态射 $f:X\to S$ 的 Stein 分解
\[X\xrightarrow{f'}S'\longrightarrow S.\]
有 $S'=\Spec A'$, 其中 $A'$ 是一个有限 $A$ 代数. 因为 $A$ 是 hensel 的, $A'$ 分解为一些局部 $A$ 代数 $A'_i$ 的积, 它们由 $S'_s$ 的点 $i$ 参数化. 令 $S'_i=\Spec A'_i$, $X_i=S'_i\times_SX$, 于是 $X$ 是各 $X_i$ 的不交并. 由 8.2.13, $X_i$ 在 $i$ 处的纤维 $(X_i)_i=f'^{-1}(i)$ 连通. 因为 $X_i$ 在 $S'_i$ 上固有而 $S'_i$ 为局部的, $X_i$ 的任何分支都不可能与其特殊纤维不交, 因而 $X_i$ 连通. 所以各 $X_i$ 就是 $X$ 的连通分支; 把一个分支对应到其特殊纤维, 便得到它们与 $X_s$ 的连通分支的一一对应.

\paragraph{8.2.19. 基变换映射.}
设
\begin{equation}
\begin{tikzcd}X'\arrow[r,"h"]\arrow[d,"f'"']&X\arrow[d,"f"]\\Y'\arrow[r,"g"']&X\end{tikzcd}
\tag{8.2.19.1}\label{ch08:eq82191}\end{equation}
% The original diagram prints X at lower right, while subsequent text uses Y; retained.
是一个赋环空间的交换方块, $F$ 是一个 $\mathcal O_X$ 模. 则存在典范的 $\mathcal O_{Y'}$ 模映射
\begin{equation}\gamma:g^*f_*F\longrightarrow f'_*h^*F,\tag{8.2.19.2}\label{ch08:eq82192}\end{equation}
称为\emph{基变换映射}, 它可用下列两种等价方式定义. 令 $a=gf'=fh$.
\begin{enumerate}
\item[(a)] 由 $g^*$ 与 $g_*$ 的伴随, 定义 $\gamma$ 等价于定义
\[\gamma_1:f_*F\longrightarrow g_*f'_*h^*F=a_*h^*F.\]
有 $a_*h^*F=f_*h_*h^*F$, 对伴随映射 $F\to h_*h^*F$ 施行 $f_*$, 就定义了 $\gamma_1$.
\item[(b)] 由 $f'^*$ 与 $f'_*$ 的伴随, 定义 $\gamma$ 等价于定义
\[\gamma_2:f'^*g^*f_*F=a^*f_*F\longrightarrow h^*F.\]
有 $a^*f_*F=h^*f^*f_*F$, 对伴随映射 $f^*f_*F\to F$ 施行 $h^*$, 就定义了 $\gamma_2$.
\end{enumerate}
这两个定义的等价性是一个非平凡事实, 由 Deligne 在更一般的背景下证明 [SGA4, XVII].

沿着同样的思路, 对所有 $q\in\mathbb Z$, 可定义典范映射
\begin{equation}\gamma:g^*R^qf_*F\longrightarrow R^qf'_*(h^*F),\tag{8.2.19.3}\label{ch08:eq82193}\end{equation}
同样称为基变换映射. 再次利用 $g^*$ 与 $g_*$ 的伴随, 这等价于定义
\[\gamma_1:R^qf_*F\longrightarrow g_*R^qf'_*(h^*F).\]
将 $\gamma_1$ 定义为下列两个映射的复合 $vu$:
\[u:R^qf_*F\longrightarrow R^qa_*(h^*F),\]
% Source PDF 208 / printed p.196.
\[v:R^qa_*(h^*F)\longrightarrow g_*R^qf'_*(h^*F).\]
映射 $u$ 是上同调上的经典函子性映射. 也就是说, 在 $D^+(X)$ 中有伴随映射
\[\alpha:F\longrightarrow Rh_*h^*F,\]
定义为复合 $F\to h_*h^*F\to h_*C(h^*F)$, 其中第一个映射为经典伴随映射, 第二个则由选取 $h^*F$ 的一个由 $h_*$ 无上同调模组成的分解 $h^*F\to C(h^*F)$ 给出. 对 $\alpha$ 施行 $Rf_*$, 得到映射
\[Rf_*(\alpha):Rf_*F\longrightarrow Rf_*Rh_*h^*F=Ra_*h^*F,\]
取上同调层即得 $u$. 换言之, 若 $V$ 是 $Y$ 的开子集, $U=f^{-1}(V)$, $U'=a^{-1}(V)=h^{-1}(U)$, 则 $R^qf_*F$ 是预层 $V\mapsto H^q(U,F)$ 的相伴层, $R^qa_*(h^*F)$ 是预层 $V\mapsto H^q(U',h^*F)$ 的相伴层, 而 $u$ 对应于函子性映射 $H^q(U,F)\to H^q(U',h^*F)$.

映射 $v$ 是谱序列
\[E_2^{ij}=R^ig_*R^jf'_*(h^*F)\Longrightarrow R^{i+j}a_*(h^*F)\]
的一个边缘同态 $H^q\to E_\infty^{0q}\to E_2^{0q}$. 更明确地说, 用上述记号且令 $V'=g^{-1}(V)$, $v$ 对应于映射
\[H^q(U',h^*F)\longrightarrow H^0(V',R^qf'_*(h^*F)),\]
它由把 $H^q(U',h^*F)$ 中的元素限制到开子集 $f'^{-1}(W)$ 得到, 其中 $W$ 在 $V'$ 中开.

在适当的上同调有限性假设下, 可以在 $D(Y')$ 中定义一个基变换映射
\begin{equation}Lg^*Rf_*F\longrightarrow Rf'_*Lh^*F,\tag{8.2.19.4}\label{ch08:eq82194}\end{equation}
它诱导 (8.2.19.3) (参见 [SGA4, XVII 4.1.5]). 然而, 当 (8.2.19.1) 为概形的笛卡尔方块, $F$ 是拟凝聚层时, 这一映射一般并无良好性质 (见 8.3.5).

\section{上同调平坦性}\label{ch08:sec83}
\paragraph{8.3.1.}
本节的结果不用于 §8.4, 它们是对 §8.2 的补充. 更确切地说, 按照 [EGAIII2, 7], 我们讨论如下问题: 在 8.2.10 的情形下, 若 $F$ 在 $Y$ 上平坦, 何时可以断言各个基变换映射
\[R^qf_*(F)_y/\mathfrak m_y^{n+1}R^qf_*(F)_y\longrightarrow H^q(X_y,F_n)\]
为同构? 更一般地, 何时可以断言 $R^qf_*(F)$ 的形成与任意基变换交换, 何时 $R^qf_*(F)$ 是有限型局部自由模? 正如 [SGA6, III] 所表明的, 使用导出范畴可简化 [EGAIII2, 7] 中的叙述. 其他叙述见 [Har77, III 12] 与 [Mum70, 5].

下文中, 若 $X$ 是概形, 用 $D(X)$ 表示 $\mathcal O_X$ 模范畴的导出范畴. 主要工具是下面的\emph{基变换公式}:
% Source PDF 209 / printed p.197.
\begin{origstatement}{定理}{8.3.2}\label{ch08:th832}
设
\begin{equation}
\begin{tikzcd}X'\arrow[r,"h"]\arrow[d,"f'"']&X\arrow[d,"f"]\\Y'\arrow[r,"g"']&Y\end{tikzcd}
\tag{8.3.2.1}\label{ch08:eq8321}\end{equation}
是概形的笛卡尔方块, 其中 $X,Y$ 拟紧且分离. 设 $F$ (相应地, $G$) 是 $X$ (相应地, $Y'$) 上的拟凝聚层. 假定 $F$ 与 $G$ 在 $Y$ 上 Tor 无关, 即对所有满足 $g(y')=f(x)$ 的点 $x\in X$, $y'\in Y'$, 以及所有 $q>0$, 有
\[\operatorname{Tor}_q^{\mathcal O_{Y,y}}(G_{y'},F_x)=0\]
(例如当 $F$ 或 $G$ 在 $Y$ 上平坦时即如此). 则在 $D(Y')$ 中有自然同构
\begin{equation}G\otimes_Y^LRf_*F\xrightarrow\sim Rf'_*(G\otimes_YF),\tag{8.3.2.2}\label{ch08:eq8322}\end{equation}
其中 $G\otimes_Y^LRf_*F:=G\otimes^LLg^*Rf_*F$, $G\otimes_YF=f'^*G\otimes h^*F$.
\end{origstatement}
当 $Y=Y'$ (相应地, $G=\mathcal O_{Y'}$) 时, 同构 (8.3.2.2) 称为\emph{投影同构} (相应地, \emph{基变换同构}). 当 $G=\mathcal O_{Y'}$ 时, 由 (8.3.2.2) 对 $q\in\mathbb Z$ 得到典范映射
\begin{equation}g^*R^qf_*F\longrightarrow R^qf'_*(h^*F).\tag{8.3.2.3}\label{ch08:eq8323}\end{equation}
这一映射是典范映射 $g^*R^qf_*F\to H^q(Lg^*Rf_*F)$ 与对 (8.3.2.2) 施行 $H^q$ 所得到的同构 $H^q(Lg^*Rf_*F)\xrightarrow\sim R^qf'_*(h^*F)$ 的复合. 从 (8.3.2.2) 的构造可知, 这一映射就是 (8.2.19.3) 中定义的基变换映射. 它一般不是同构. 这一问题将在 8.3.10--8.3.11 中讨论.

下列推论是最有用的特殊情形:
\begin{origstatement}{推论}{8.3.3}\label{ch08:cor833}
若笛卡尔方块 (8.3.2.1) 中的 $g$ 平坦, 则 (8.3.2.2) 给出基变换同构
\[g^*Rf_*F\xrightarrow\sim Rf'_*h^*F,\]
且所诱导的基变换映射 (8.3.2.3) 均为同构.
\end{origstatement}

\begin{origstatement}{推论}{8.3.4}\label{ch08:cor834}
设 $f:X\to Y$ 是拟紧分离概形之间的态射. 设 $y$ 是 $Y$ 的一点, 以 $X_y$ 记 $f$ 在 $y$ 处的纤维, 即 $\Spec k(y)\times_YX$, 设 $F$ 是 $X$ 上在 $Y$ 上平坦的拟凝聚层. 则 (8.3.2.2) 给出自然同构 (在 $k(y)$ 向量空间的导出范畴中)
\[k(y)\otimes_{\mathcal O_Y}^LRf_*F\xrightarrow\sim R\Gamma(X_y,\mathcal O_{X_y}\otimes_{\mathcal O_X}F).\]
\end{origstatement}

现在证明 8.3.2. 首先考虑 $X,Y,Y'$ 仿射, 其环分别为 $B,A,A'$ 的情形. 此时 $X'$ 为环 $B'=A'\otimes_AB$ 的仿射概形, 并且对某个 $B$ 模 $M$ 与 $A'$ 模 $N$, 有 $F=\widetilde M$, $G=\widetilde N$. 于是 $Rf_*F$ 由 $M$ 的底 $A$ 模 $M_{[A]}$ 表示, 而 $Rf'_*(G\otimes_YF)$ 由 $(N\otimes B')\otimes_{B'}(B'\otimes M)$ 的底 $A'$ 模 $(N\otimes_AM)_{[A']}$ 表示. 另一方面, $G\otimes_Y^LRf_*F$ 由 $N\otimes_A^LM_{[A]}:=N\otimes_{A'}^L(A'\otimes_A^LM_{[A]})$ 表示; 若 $P$ 是 $M_{[A]}$ 的平坦分解, 这可计算为 $N\otimes_AP$. Tor 无关假设意指对 $q>0$, 有 $\operatorname{Tor}_q^A(N,M_{[A]})=0$; 即自然映射
\begin{equation}N\otimes_A^LM_{[A]}\longrightarrow N\otimes_AM_{[A]}\tag{*}\label{ch08:eq832star}\end{equation}
% Source PDF 210 / printed p.198.
是同构 (在 $D(A')$ 中). 同构 (8.3.2.2) 就是 (*) 与 (平凡的) 同构
\begin{equation}N\otimes_AM_{[A]}\xrightarrow\sim(N\otimes_AM)_{[A']}\tag{**}\label{ch08:eq832star2}\end{equation}
的复合.

现在假定态射 $f$ 仿射 (但概形 $Y$ 不必仿射). 此时对某个拟凝聚 $\mathcal O_Y$ 代数 $B$, 有 $X=\Spec B$; 对某个拟凝聚 $B$ 模 $M$, 有 $F=\widetilde M$. 仍有 $Rf_*F=f_*F$, 它由 $M$ 的底 (拟凝聚) $\mathcal O_Y$ 模 $M_{[A]}$ 表示. 将前面的讨论应用于 $Y$ 的仿射开子集之上的 $Y'$ 的仿射开子集, 得到自然等同
\[\begin{aligned}
G\otimes_Y^LRf_*F&\xrightarrow\sim G\otimes g^*f_*F\\
&\xrightarrow\sim f'_*(G\otimes_YF)\xrightarrow\sim Rf'_*(G\otimes_YF).
\end{aligned}\]
其复合定义同构 (8.3.2.2).

在一般情形, 选取 $X$ 的一个有限仿射开覆盖 $\mathcal U=(U_i)_{i\in I}$ ($I=\{1,\ldots,r\}$). 因为 $X,Y$ 分离, 各 $U_i$ 的任意有限交 $U_{i_0\cdots i_n}=U_{i_0}\cap\cdots\cap U_{i_n}$ ($i_0<\cdots<i_n$) 都在 $Y$ 上仿射 [EGAII, 1.6.2]. 因而 (由 [EGAIII1, 1.4]) 有
\[Rf_*F=f_*\check C(\mathcal U,F),\]
其中 $\check C(\mathcal U,F)$ 是覆盖 $\mathcal U$ 的以 $F$ 为值的交错 Čech 复形. 由 $f$ 仿射时的讨论得到同构
\begin{equation}\begin{aligned}
G\otimes_Y^LRf_*F&\xrightarrow\sim G\otimes g^*f_*\check C(\mathcal U,F)\\
&\xrightarrow\sim f'_*\check C(\mathcal U',G\otimes_YF)\xrightarrow\sim Rf'_*(G\otimes_YF),
\end{aligned}\tag{***}\label{ch08:eq832star3}\end{equation}
其中 $\mathcal U'$ 为由各 $U_i$ 的逆像组成的 $X'$ 的覆盖. 容易验证, 复合 (***) 不依赖于 $\mathcal U$ 的选择. 我们就用这一复合作为同构 (8.3.2.2) 的定义.

(8.3.2.3) 与 (8.2.19.3) 的相容性留给读者验证.

\begin{origstatement}{注}{8.3.5}\label{ch08:rem835}
容易将 8.3.2 推广到 $Y$ 拟紧且 $f$ 拟紧拟分离的情形 (在论证的最后一部分, 各交 $U_{i_0\cdots i_n}$ 仅为拟紧的, 必须把 Čech 复形换成适当的 ``超 Čech'' 变体).

不过, 去掉 Tor 无关假设似乎很困难. 例如, 当 (8.3.2.1) 如 8.3.2 证明开头那样是仿射概形的笛卡尔方块, $F=\mathcal O_X$, 但不作 Tor 无关假设时, 确实有一个形如 (8.2.19.4) 的基变换映射, 即 $D(A')$ 中对应于映射
\[A'\otimes_A^LB\longrightarrow A'\otimes_AB\]
的映射; 但这一映射为同构当且仅当 $A'$ 与 $B$ 在 $A$ 上 Tor 无关.

为得到令人满意的形式体系, 必须使用同伦代数的一些工具, 例如环的导出张量积. 至今还没有写成的系统叙述.
\end{origstatement}

\paragraph{8.3.6.}
8.3.2 的主要应用是 $f$ 为 Noether 概形的固有态射, $F$ 是 $X$ 上在 $Y$ 上平坦的凝聚层的情形. 此时复形 $Rf_*F$ 具有良好性质, 即它是\emph{完美复形}, 而基变换公式 8.3.4 使我们能在 $Y$ 的一点 $y$ 附近分析各上同调层 $R^qf_*F$ 与基变换的相容性.

先回顾 $D(X)$ 中对象的若干基本有限性条件, 其中 $X$ 为局部 Noether 概形. [SGA6, I, II, III] 在更大的一般性下讨论了这些条件.
% Source PDF 211 / printed p.199.
主要有三个条件: 伪凝聚性、有限 Tor 维数和完美性; 最后一个是前两个的结合.

\paragraph{8.3.6.1. 伪凝聚性.}
若复形 $E\in D^b(X)$ 具有凝聚上同调 (即对所有 $q$, $H^q(E)$ 凝聚), 则称它为\emph{伪凝聚的 (pseudo-coherent)}. 通常用 $D^b(X)_{\mathrm{coh}}$ 表示 $D^b(X)$ 中由伪凝聚复形组成的满子范畴. 它是一个三角子范畴. 若 $X$ 仿射, $E\in D(X)$ 满足对 $q>a$, $H^q(E)=0$, 则 $E$ 伪凝聚当且仅当 $E$ 在 $D(X)$ 中同构于一个复形 $L\in D^b(X)$, 使对 $q>a$ 有 $L^q=0$, 且对所有 $q$, $L^q$ 都是有限型自由模. 特别地, 在任意局部 Noether 概形 $X$ 上, 伪凝聚复形在导出范畴中局部同构于一个由有限型自由 $\mathcal O$ 模组成的有上界复形; 对 $X$ 的任何一点 $x$, 茎 $E_x$ 作为 $\mathcal O_{X,x}$ 模复形, (在导出范畴 $D(\mathcal O_{X,x})$ 中) 同构于一个由有限型自由 $\mathcal O_{X,x}$ 模组成的有上界复形.
% Source prints D^b for L here; retained despite apparent boundedness issue.
当 $E$ 的各分量拟凝聚时, 通过简单的逐步构造即可证明上述断言 [EGAIII1, Ch.~0, 11.9.1]. 在一般情形, 必须先把 $E$ 替换为一个分量拟凝聚的有界复形, 这较为细致 (见 [Har66, II 7.19]、[SGA6, II 2.2.1]).

\paragraph{8.3.6.2. 有限 Tor 维数.}
设 $a,b\in\mathbb Z$, $a\leq b$. 若复形 $E\in D(X)$ 满足以下等价条件, 则称其 \emph{Tor 幅度在 $[a,b]$ 中}:
\begin{enumerate}
\item[(i)] $E$ 在 $D(X)$ 中同构于复形 $L$, 其中当 $q\notin[a,b]$ 时 $L^q=0$, 且对所有 $q$, $L^q$ 都平坦;
\item[(ii)] 对任意 $\mathcal O_X$ 模 $M$, 当 $q\notin[a,b]$ 时, 有 $H^q(M\otimes^LE)=0$.
\end{enumerate}
(i) 与 (ii) 等价的证明是直接的. 若存在区间 $[a,b]$, 使复形 $E$ 的 Tor 幅度在 $[a,b]$ 中, 则称 $E$ 具有\emph{有限 Tor 维数} (或\emph{有限 Tor 幅度}). $D(X)$ 中由有限 Tor 维数的复形组成的满子范畴是一个三角子范畴.

复形 $E$ 的 Tor 幅度在 $[a,b]$ 中, 当且仅当对所有 $x\in X$, 茎 $E_x$ 作为 $\mathcal O_{X,x}$ 模复形的 Tor 幅度在 $[a,b]$ 中, 即它在 $D(\mathcal O_{X,x})$ 中同构于一个次数集中在 $[a,b]$ 且每一次数都平坦的复形 $L$; 等价地, 对任何 $\mathcal O_{X,x}$ 模 $M$, 当 $q\notin[a,b]$ 时, 有 $H^q(M\otimes^LE_x)=0$.

\paragraph{8.3.6.3. 完美性.}
若复形 $E\in D^b(X)$ 伪凝聚且局部具有有限 Tor 维数, 则称它是\emph{完美的 (perfect)}. 若它伪凝聚且 Tor 幅度在 $[a,b]$ 中 (其中 $a,b\in\mathbb Z$, $a\leq b$), 则称其\emph{完美幅度在 $[a,b]$ 中}. \emph{严格完美复形}是指由有限型局部自由模组成的有界复形. $D^b(X)$ 中由完美复形组成的满子范畴是一个三角子范畴.

因为 $\mathcal O_X$ 模为有限型局部自由模当且仅当它凝聚且平坦, 由 8.3.6.1 和 8.3.6.2 可知, 复形 $E\in D^b(X)$ 完美当且仅当它在导出范畴中局部同构于一个严格完美复形. 同样地, 有下列有用的判别法:

\begin{origstatement}{命题}{8.3.6.4}\label{ch08:prop8364}
设 $x$ 是 $X$ 的一点, $E\in D^b(X)$ 是 $X$ 上的伪凝聚复形, 对某个区间 $[a,b]$, 当 $q\notin[a,b]$ 时 $H^q(E_x)=0$. 则下列条件等价:
\begin{enumerate}
\item[(i)] 在 $x$ 的某个邻域中, $E$ 的完美幅度在 $[a,b]$ 中;
\item[(ii)] $H^{a-1}(k(x)\otimes^LE)=0$;
% Source PDF 212 / printed p.200.
\item[(iii)] 在 $x$ 的某个邻域中, $E$ 在导出范畴中同构于一个由有限型自由 $\mathcal O$ 模组成且次数集中在 $[a,b]$ 的复形.
\end{enumerate}
\end{origstatement}
由 8.3.6.1, 可以假定 $E$ 的分量凝聚, 次数集中在 $[a,b]$, 且当 $q>a$ 时 $E^q$ 为有限型自由模. 我们需说明 (ii) 蕴含 $E^a$ 在 $x$ 附近为有限型局部自由模. 从正合列
\[0\longrightarrow E^{[a+1,b]}\longrightarrow E\longrightarrow E^a[-a]\longrightarrow0,\]
其中 $E^{[a+1,b]}$ 是 $E$ 在次数 $\geq a+1$ 的朴素截断, 得到
\[\operatorname{Tor}_1^{\mathcal O_{X,x}}(k(x),E_x^a)=0.\]
由标准平坦性判别法 [Bou61, III, §5, th.~3], 这蕴含 $E_x^a$ 为有限型自由模, 因而 $E^a$ 在 $x$ 的某个邻域中为有限型自由模.

\begin{origstatement}{推论}{8.3.6.5}\label{ch08:cor8365}
设 $x$ 是 $X$ 的一点, $q\in\mathbb Z$, $E\in D^b(X)$ 是 $X$ 上的伪凝聚复形, 满足对某个整数 $b$, 当 $i>b$ 时 $H^iE=0$. 对 $i\in\mathbb Z$, 令
\[\alpha^i(x):k(x)\otimes H^i(E)\longrightarrow H^i(k(x)\otimes^LE)\]
表示典范映射.

(a) 下列条件等价:
\begin{enumerate}
\item[(i)] $\alpha^q(x)$ 满射;
\item[(ii)] $\tau_{>q}E$ 在 $x$ 的某个邻域中的完美幅度在 $[q+1,b]$ 中.
\end{enumerate}
当这些条件满足时, 存在 $x$ 的开邻域 $U$, 使对所有 $y\in U$, $\alpha^q(y)$ 双射, 而且对 $U$ 上所有拟凝聚模 $M$, 自然映射
\[\alpha^q(M):M\otimes H^q(E)\longrightarrow H^q(M\otimes^LE)\]
双射.

(b) 假定 (a) (i) 成立. 则下列条件等价:
\begin{enumerate}
\item[(i)] $\alpha^{q-1}(x)$ 满射;
\item[(ii)] $H^q(E)$ 在 $x$ 的某个邻域中为有限型局部自由模.
\end{enumerate}
\end{origstatement}
这里, 若 $L$ 是 (某个 Abel 范畴中的) 复形, $i\in\mathbb Z$, 则 $\tau_{\geq i}L$ 表示 $L$ 在次数 $\geq i$ 的\emph{典范截断}, 定义为 $0\to L^i/dL^{i-1}\to L^{i+1}\to\cdots$, 并置 $\tau_{>i}=\tau_{\geq i+1}$.

现在证明 (a). 投影 $E\to\tau_{\geq q}E$ 诱导同构
\begin{equation}H^q(k(x)\otimes^LE)\xrightarrow\sim H^q(k(x)\otimes^L\tau_{\geq q}E).\tag{*}\label{ch08:eq8365star}\end{equation}
考虑典范的区别三角
\begin{equation}H^q(E)[-q]\longrightarrow\tau_{\geq q}E\longrightarrow\tau_{>q}E\longrightarrow.\tag{**}\label{ch08:eq8365star2}\end{equation}
结合 (*), 从 (**) 得到
\begin{equation}\operatorname{Coker}\alpha^q(x)=H^q(k(x)\otimes^L\tau_{>q}E).\tag{***}\label{ch08:eq8365star3}\end{equation}
于是 (i) 与 (ii) 的等价性来自 8.3.6.4. 假定这些条件成立. 只需证明 (a) 的最后一个断言. 设 $U$ 是 $x$ 的开邻域, 使 $\tau_{>q}E|U$ 的完美幅度在 $[q+1,b]$ 中. 设 $M$ 是 $U$ 上的拟凝聚层. 对 (**) 施行 $M\otimes^L-$, 并注意到当 $i\leq q$ 时 $H^i(M\otimes^L\tau_{>q}E)=0$, 即得 $\alpha^q(M)$ 双射. 再证明 (b). 因为 $\tau_{>q}E$ 在 $x$ 的某个邻域中的完美幅度在 $[q+1,b]$ 中, 三角 (**) 表明: $H^q(E)$ 在 $x$ 的某个邻域中为有限型局部自由模,
% Source PDF 213 / printed p.201.
当且仅当 $\tau_{\geq q}E$ 在 $x$ 的某个邻域中的完美幅度在 $[q,b]$ 中. 而由 (i), 这一条件等价于 $\alpha^{q-1}(x)$ 的满射性.

\paragraph{8.3.7.}
设 $f:X\to Y$ 是局部 Noether 概形的一个态射, $F$ 是 $X$ 上的凝聚层. 如前所述, 8.3.2 的主要应用涉及 $F$ 在 $Y$ 上平坦的情形. 按定义, $F$ 在 $Y$ 上平坦当且仅当对所有 $x\in X$, $\mathcal O_{X,x}$ 模 $F_x$ 在 $\mathcal O_{Y,y}$ 上平坦, 其中 $y=f(x)$. 利用平坦性判别法 [Bou61, III, §5, th.~3], 常可方便地将这一点表达如下: $F$ 在 $Y$ 上平坦, 当且仅当对所有 $y\in Y$, 自然的 (满射) 映射
\begin{equation}\operatorname{gr}^n\mathcal O_{Y,y}\otimes_{k(y)}\operatorname{gr}^0F_y\longrightarrow\operatorname{gr}^nF_y\tag{8.3.7.1}\label{ch08:eq8371}\end{equation}
对所有 $n\geq0$ 都是同构, 其中 $\operatorname{gr}$ 意指 $\mathcal O_{Y,y}$ 上 $\mathfrak m_y$ 进滤过的伴随分次 ($\mathfrak m_y$ 为极大理想), $F_y$ 为 $F$ 在 $\Spec\mathcal O_{Y,y}\times_YX$ 上的逆像. (8.3.7.1) 的双射性也等价于 $\operatorname{Tor}_1^{\mathcal O_{Y,y}}(k(y),F_y)=0$; 即自然映射 $\tau_{\geq-1}(k(y)\otimes^LF)\to k(y)\otimes F$ 为同构, 也等价于对每个 $n\geq0$, $F_y/\mathfrak m_y^{n+1}F$ 在 $\Spec\mathcal O_{Y,y}/\mathfrak m_y^{n+1}$ 上平坦.

\begin{origstatement}{定理}{8.3.8}\label{ch08:th838}
设 $f:X\to Y$ 是 Noether 概形的固有态射, $F$ 是 $X$ 上的凝聚层. 则 $Rf_*F$ 在 $Y$ 上伪凝聚 (8.3.6.1). 若 $F$ 在 $Y$ 上平坦, 则 $Rf_*F$ 完美 (8.3.6.3).
\end{origstatement}
第一个断言只是 Grothendieck 有限性定理 [EGAIII1, 3.2.1] (断言各层 $R^qf_*F$ 凝聚) 与 [EGAIII1, 1.4.12] (蕴含 $Rf_*F$ 属于 $D^b(Y)$) 的另一种表述. 为证明第二个断言, 可假定 $Y$ 仿射. 设 $N$ 为整数, 使对 $X$ 上所有拟凝聚层 $E$ 及所有 $q>N$, 都有 $R^qf_*E=0$ (可取 $N$, 使 $X$ 有一个由 $N+1$ 个仿射开子集组成的覆盖 [EGAIII1, 1.4.12]). 由 (8.3.2.2), 对任意拟凝聚 $\mathcal O_Y$ 模 $G$, 有
\[G\otimes^LRf_*F\xrightarrow\sim Rf_*(G\otimes_YF),\]
特别地, 当 $q\notin[0,N]$ 时,
\[H^q(G\otimes^LRf_*F)=0.\]
因此, 对 $Y$ 的任意一点 $y$ 和任何 $\mathcal O_{Y,y}$ 模 $M$, 当 $q\notin[0,N]$ 时, 有
\[H^q(M\otimes^L(Rf_*F)_y)=0.\]
由 8.3.6.2, 这意味着 $Rf_*F$ 的完美幅度在 $[0,N]$ 中.

\paragraph{8.3.9.}
在 8.3.8 的假设下, 设 $F$ 在 $Y$ 上平坦, $Y$ 仿射, $Y=\Spec A$, $Y'$ 为理想 $I$ 定义的 $Y$ 的闭子概形. 如 [EGAIII2, 7.4.8] 所说明的, $Rf_*F$ 的伪凝聚性结合基变换公式 8.3.2, 给出了映射 $\varphi_q$ (8.2.4.1) 为同构这一事实的另一证明 (在这一特殊情形下).

由 8.3.2, 对 $n\geq0$, 有
\begin{equation}A/I^{n+1}\otimes^LR\Gamma(X,F)\xrightarrow\sim R\Gamma(X,F_n).\tag{*}\label{ch08:eq839star}\end{equation}
映射 $\varphi_q$ 是下列映射的逆极限:
\[\varphi_{q,n}:A/I^{n+1}\otimes H^q(X,F)\longrightarrow H^q(X,F_n),\]
它由自然映射 $A/I^{n+1}\otimes H^q(X,F)\to H^q(A/I^{n+1}\otimes^LR\Gamma(X,F))$ 与同构 $H^q(*)$ 的复合得到. 因为 $Rf_*F$ 伪凝聚, $R\Gamma(X,F)$
% Source PDF 214 / printed p.202.
同构于一个由有限型自由 $A$ 模组成的有上界复形 $P$. 映射 $\varphi_{q,n}$ 可重写为
\[H^q(P)/I^{n+1}H^q(P)\longrightarrow H^q(P/I^{n+1}P).\]
一般地, 这些映射没有一个是同构, 但由 Artin--Rees 定理可知, 极限映射
\[\lim H^q(P)/I^{n+1}H^q(P)\longrightarrow\lim H^q(P/I^{n+1}P)\]
是同构.

\paragraph{8.3.10.}
设 $f:X\to Y$ 是分离 Noether 概形的一个固有态射, $F$ 是 $X$ 上在 $Y$ 上平坦的凝聚层. 设 $q\in\mathbb Z$. 若对任意态射 $g:Y'\to Y$, 基变换映射 (8.3.2.3)
\begin{equation}g^*R^qf_*F\longrightarrow R^qf'_*F'\tag{8.3.10.1}\label{ch08:eq83101}\end{equation}
都是同构, 则称 $F$ 在 $Y$ 上\emph{在次数 $q$ 上同调平坦}, 其中用 (8.3.2.1) 的记号, $F'=h^*F$. 当 $F=\mathcal O_X$ (即 $f$ 平坦) 时, 简称 $f$ 在次数 $q$ 上同调平坦.

若 $y$ 是 $Y$ 的一点, $X_y=\Spec k(y)\times_YX$ 为 $f$ 在 $y$ 处的纤维, 则映射 (8.3.10.1) 写成
\begin{equation}k(y)\otimes R^qf_*F\longrightarrow H^q(X_y,F/\mathfrak m_yF).\tag{8.3.10.2}\label{ch08:eq83102}\end{equation}
仿照 8.3.6.5 中的记号, 我们将它记为 $\alpha^q(y)$. 下面的判别法是将 8.3.6.5 应用于 $Y$ 上的伪凝聚复形 $E=Rf_*F$ 的直接推论 (参见 [EGAIII2, 7.8.4]、[Har77, III, 12.11]):

\begin{origstatement}{推论}{8.3.11}\label{ch08:cor8311}
设 $f:X\to Y$ 与 $F$ 如 8.3.10 中所述, $q\in\mathbb Z$, $y$ 为 $Y$ 的一点. 设整数 $b$ 满足当 $i>b$ 时 $R^if_*F=0$.

(a) 下列条件等价:
\begin{enumerate}
\item[(i)] 映射 $\alpha^q(y)$ (8.3.10.2) 满射;
\item[(ii)] $\tau_{>q}Rf_*F$ 在 $y$ 的某个邻域中的完美幅度在 $[q+1,b]$ 中.
\end{enumerate}
当这些条件满足时, 存在 $y$ 的开邻域 $U$, 使对所有 $z\in U$, $\alpha^q(z)$ 双射, 并使 $F|f^{-1}(U)$ 在 $U$ 上在次数 $q$ 上同调平坦.

(b) 假定 (a) (i) 成立. 则下列条件等价:
\begin{enumerate}
\item[(i)] $\alpha^{q-1}(y)$ 满射;
\item[(ii)] $R^qf_*F$ 在 $y$ 的某个邻域中为有限型局部自由模.
\end{enumerate}
\end{origstatement}

\begin{origstatement}{注}{8.3.11.1}\label{ch08:rem83111}
因为 $q=-1$ 时条件 (a) 对所有 $y$ 都显然满足, 所以得到 $Rf_*F$ 的完美幅度在 $[0,b]$ 中, 如 8.3.8 证明末尾已经观察到的.

另一方面, $q=b$ 时条件 (b) 对所有 $y$ 都显然满足. 因而 $F$ 在次数 $b$ 上同调平坦.
\end{origstatement}

\begin{origstatement}{注}{8.3.11.2}\label{ch08:rem83112}
(参见 [EGAIII1, 4.6.1]) 下列判别法十分有用: 若 $H^{q+1}(X_y,F/\mathfrak m_yF)=0$, 则 $\alpha_q(y)$ 满射; 特别地, 由 (b), 若 $H^1(X_y,F/\mathfrak m_yF)=0$, 则在 $y$ 的某个邻域中, $f_*F$ 为有限型局部自由模且与基变换交换.
\end{origstatement}
事实上, 若 $H^{q+1}(X_y,F/\mathfrak m_yF)=0$, 则形式函数定理 8.2.10 蕴含
\[(R^qf_*(F)_y)^\wedge=0,\]
% Source PDF 215 / printed p.203.
从而 $R^qf_*(F)_y=0$, 因此 $\tau_{>q}Rf_*(F)_y=\tau_{>q+1}Rf_*(F)_y$. 因为 (显然) $\alpha_{q+1}(y)=0$, 由 (a) 可知 $\tau_{>q}Rf_*F$ 的完美幅度在 $[q+2,b]$ 中, 因而当然也在 $[q+1,b]$ 中, 所以 $\alpha_q(y)$ 满射.
% The original has q (not q+1) in both preceding vanishing formulas; retained for review.

\paragraph{8.3.12.}
假定 8.3.11 的 (a) (i) 对所有 $y\in Y$ 都成立, 即 $\tau_{>q}Rf_*F$ 的完美幅度在 $[q+1,b]$ 中. 对偶
\[K=R\mathcal Hom(\tau_{>q}Rf_*F,\mathcal O_Y)\]
是完美复形, 其完美幅度在 $[-b,-q-1]$ 中. 令
\[Q:=H^{-q-1}(K).\]
则对任何拟凝聚 $\mathcal O_Y$ 模 $M$, 有自然同构
\begin{equation}R^{q+1}f_*(M\otimes_YF)\xrightarrow\sim\mathcal Hom(Q,M).\tag{8.3.12.1}\label{ch08:eq83121}\end{equation}
此外, $Q$ 的形成与任意基变换交换. (这就是所谓的\emph{交换性质}, 参见 [EGAIII2, 7.7.5, 7.7.6, 7.8.9].)

证明仍是 8.3.2 的简单应用. 由 8.3.2, 有
\begin{equation}R^{q+1}f_*(M\otimes_YF)=H^{q+1}(M\otimes^LRf_*F).\tag{*}\label{ch08:eq8312star}\end{equation}
投影 $Rf_*F\to\tau_{>q}Rf_*F$ 诱导同构
\begin{equation}H^{q+1}(M\otimes^LRf_*F)\xrightarrow\sim H^{q+1}(M\otimes^L\tau_{>q}Rf_*F).\tag{**}\label{ch08:eq8312star2}\end{equation}
令 $L:=\tau_{>q}Rf_*F$, 于是 $K=R\mathcal Hom(L,\mathcal O_Y)$. 因为 $L$ 完美, 有自然的双重对偶同构
\[L\xrightarrow\sim R\mathcal Hom(K,\mathcal O_Y),\]
它诱导同构
\begin{equation}M\otimes^LL\xrightarrow\sim R\mathcal Hom(K,M).\tag{***}\label{ch08:eq8312star3}\end{equation}
复合 (*)、(**) 和 $H^{q+1}(***)$, 得到
\[R^{q+1}f_*(M\otimes_YF)\xrightarrow\sim H^{q+1}R\mathcal Hom(K,M)=\mathcal Ext^{q+1}(K,M).\]
但因为 $K$ 的完美幅度在 $[-b,-q-1]$ 中, 即局部同构于一个由自由模组成且次数集中在 $[-b,-q-1]$ 的复形, 所以
\[\mathcal Ext^{q+1}(K,M)=\mathcal Hom(H^{-q-1}K,M),\]
这就给出 (8.3.12.1). $Q$ 与基变换相容的证明留给读者.

\paragraph{8.3.13.}
在 8.3.8 的假设下, $Rf_*F$ 的完美性蕴含 $Y$ 上函数
\[y\longmapsto\operatorname{rk}H^q(X_y,F/\mathfrak m_yF)\]
(对固定的 $q$), 以及
\[y\longmapsto\sum(-1)^q\operatorname{rk}H^q(X_y,F/\mathfrak m_yF)\]
的良好性质. 第一个函数\emph{上半连续}, 第二个函数\emph{局部常值}. 这来自 8.3.4. 验证留给读者. 关于这些问题的详细讨论, 见 [EGAIII2, 7.7, 7.9] 与 [SGA6, III].

% Source PDF 216 / printed p.204.
\section{存在定理}\label{ch08:sec84}
\paragraph{8.4.1.}
设 $A$ 是进 Noether 环 (8.1.1), $I$ 为 $A$ 的一个定义理想, $Y=\Spec A$, $Y_n=\Spec A/I^{n+1}$, $\widehat Y=\operatorname{colim}_nY_n=\operatorname{Spf}(A)$. 本节讨论如下问题: 给定一个进 Noether $\widehat Y$ 形式概形 $\mathcal Z=\operatorname{colim}_n Z_n$ (8.1.5), 何时可以断言存在 (唯一的) (适当的) $Y$ 上的局部 Noether 概形 $X$, 使它的 $I$ 进完备化 $\widehat X=\operatorname{colim}_nX_n$ ($X_n=X\times_YY_n$) 同构于 $\mathcal Z$? 这就是所谓的\emph{代数化 (algebraization)} 问题. 如同复解析几何中的类似问题 (Serre 的 GAGA), Grothendieck 的处理方法是先固定 $X$, 比较 $X$ 与 $\widehat X$ 上的凝聚层. 基本结果为下述定理:

\begin{origstatement}{定理}{8.4.2}\label{ch08:th842}
[EGAIII1, 5.1.4] 设 $X$ 是 $Y$ 上分离且有限型的 Noether 概形, $\widehat X$ 为其如 8.4.1 中所述的 $I$ 进完备化. 则函子 $F\mapsto\widehat F$ (8.1.6) 从 $X$ 上支集在 $Y$ 上固有的凝聚层范畴到 $\widehat X$ 上支集在 $\widehat Y$ 上固有的凝聚层范畴, 是一个等价.
\end{origstatement}
回顾: $\widehat X$ 上凝聚层 $\mathcal E$ 的支集是 $X_0$ 上 $E_0=\mathcal E\otimes\mathcal O_{X_0}$ 的支集 (8.1.4). 若它作为 $X_0$ 的闭子集在 $Y_0$ 上固有, 则称它在 $\widehat Y$ 上固有.

\paragraph{8.4.3. 8.4.2 的证明.}
设 $F,G$ 是 $X$ 上支集在 $Y$ 上固有的凝聚层. 由 8.2.9, $\Hom(F,G)$ 是有限型 $A$ 模, 因而对 $I$ 进拓扑分离且完备, 所以自然映射
\[\Hom(F,G)\longrightarrow\Hom(\widehat F,\widehat G)\]
是同构. 这证明函子 $(-)^\wedge$ 全忠实. 还需证明它本质满射. 证明分为几步. 我们只说明要点.

(a) \emph{射影情形.} 假定 $f:X\to Y$ 射影. 设 $L$ 是 $X$ 上的丰沛线丛. 若 $M$ 是 $\mathcal O_X$ 模 (相应地, $\mathcal O_{\widehat X}$ 模), $n\in\mathbb Z$, 如常用 $M(n)$ 记 $M\otimes L^{\otimes n}$ (相应地, $M\otimes\widehat L^{\otimes n}$). 关键是下面的结果, 它是 [EGAIII1, 5.2.4] 的特殊情形:

\begin{origstatement}{引理}{8.4.3.1}\label{ch08:lem8431}
设 $E$ 是 $\widehat X$ 上的凝聚层. 则存在非负整数 $m,r$ 及一个满同态
\[\mathcal O_{\widehat X}(-m)^r\longrightarrow E.\]
\end{origstatement}
假定 8.4.3.1 成立, 先说明如何证明此时的本质满射性. 设 $E$ 为 $\widehat X$ 上的凝聚层. 由 8.4.3.1, 对某些非负整数 $m_0,m_1,r_0,r_1$, 可以找到正合列
\[\mathcal O_{\widehat X}(-m_1)^{r_1}\xrightarrow{u}\mathcal O_{\widehat X}(-m_0)^{r_0}\longrightarrow E\longrightarrow0.\]
由 $(-)^\wedge$ 的全忠实性, 存在唯一态射 $v:\mathcal O_X(-m_1)^{r_1}\to\mathcal O_X(-m_0)^{r_0}$, 使 $u=\widehat v$. 令 $F:=\operatorname{Coker}v$. 则由 $(-)^\wedge$ (在 $X$ 上凝聚层范畴中) 的正合性, $E=\widehat F$.

现在证明 8.4.3.1. 因为 $L$ 丰沛, 所以 $L_0=\mathcal O_{X_0}(1)$ 在 $X_0$ 上也丰沛. 考虑分次 $\mathcal O_{Y_0}$ 代数
\[S=\operatorname{gr}_I\mathcal O_Y=\bigoplus_{n\in\mathbb N}\widetilde I^n/\widetilde I^{n+1}=\bigoplus_{n\in\mathbb N}\mathcal I^n/\mathcal I^{n+1},\]
及分次 $f_0^*(S)$ 模 $M=\operatorname{gr}_IE=\bigoplus_{n\in\mathbb N}\mathcal I^nE/\mathcal I^{n+1}E$, 其中 $\mathcal I=I^\triangle$. 因为 $\operatorname{gr}_0E$
% Source PDF 217 / printed p.205.
在 $X_0$ 上凝聚, 典范映射 $\operatorname{gr}_I\mathcal O_Y\otimes_{\operatorname{gr}_0\mathcal O_Y}\operatorname{gr}_0E\to\operatorname{gr}_IE$ 满射, 所以 $M$ 在 $f_0^*(S)$ 上有限型, 因而对应于 $X':=\Spec f_0^*(S)$ 上的凝聚模 $\widetilde M$. 因为 $\mathcal O_X(1)$ 在 $X'$ 上的逆像 $\mathcal O_{X'}(1)$ 丰沛, 对 $\widetilde M$, $\mathcal O_{X'}(1)$ 以及由 $f_0$ 经基变换 $Y'=\Spec S\to Y$ 得到的态射 $f':X'\to Y'=\Spec S$, 应用 Serre 消没定理 [EGAIII1, 2.2.1], 得到存在整数 $n_0$, 使对所有 $n\geq n_0$, 所有 $k\in\mathbb N$ 及所有 $q>0$, 有
\[H^q(X_0,\operatorname{gr}_kE(n))=0.\]
因而对所有 $n\geq n_0$ 和 $k\geq0$, 转移映射 $H^0(X_0,E_{k+1}(n))\to H^0(X_0,E_k(n))$ 满射, 从而典范映射
\[H^0(\widehat X,E(n))=\lim_kH^0(X_0,E_k(n))\longrightarrow H^0(X_0,E_0(n))\]
满射. 因为 $\mathcal O_{X_0}(1)$ 丰沛, 可假定 $n_0$ 已取到足够大, 使存在有限多个 $E_0(n_0)$ 的整体截面生成 $E_0(n_0)$. 把这些截面提升到 $H^0(\widehat X,E(n_0))$, 就得到映射
\[u:\mathcal O_{\widehat X}(-n_0)^r\longrightarrow E,\]
使 $u_0=u\otimes\mathcal O_{X_0}:\mathcal O_{X_0}(-n_0)^r\to E_0$ 满射. 由 Nakayama 引理 (因为 $I$ 包含在 $A$ 的根基中 (8.1.1)), 这蕴含 $u$ 满射.

\begin{origstatement}{注}{8.4.3.2}\label{ch08:rem8432}
上述证明表明, 若把 $\widehat X$ 换成一个满足 $X_0=\mathcal X\times_{\widehat Y}Y_0$ 固有的进 $\widehat Y$ 形式概形 $\mathcal X$, 并把 $\mathcal O_{\widehat X}(1)$ 换成可逆 $\mathcal O_{\mathcal X}$ 模 $L$, 满足 $L_0=L\otimes\mathcal O_{X_0}$ 丰沛, 则 8.4.3.1 的结论仍成立. 证明还表明, 在这些假设下, 存在整数 $n_0$, 使对所有 $n\geq n_0$, $\Gamma(\mathcal X,E(n))\to\Gamma(X_0,E_0(n))$ 满射, 这里仍用常用记号 $E(n)=E\otimes L^{\otimes n}$.
\end{origstatement}

(b) \emph{拟射影情形.} 假定存在开浸入 $j:X\to Z$, 其中 $Z$ 在 $Y$ 上射影. 设 $E$ 是 $\widehat X$ 上的凝聚层, 其支集 $T_0$ 在 $\widehat Y$ 上固有. 则零延拓 $\widehat j_!E$ 在 $\widehat Z$ 上凝聚, 因而由 (a), 它具有 $\widehat F$ 的形式, 其中 $F$ 是 $Z$ 上的凝聚层. $F$ 的支集 $T$ 包含在 $X$ 中 (因为 $X\cap T$ 在 $T$ 中开且包含 $T_0$, 故等于 $T$), 从而 $F=j_!j^*F$, $E=(j^*F)^\wedge$.

(c) \emph{一般情形.} 对 $X$ 作 Noether 归纳. 假定对 $X$ 的每个底空间严格包含于 $X$ 的底空间的闭子概形 $T$, $\widehat T$ 上所有支集在 $\widehat Y$ 上固有的凝聚层都可代数化, 即形如 $\widehat F$, 其中 $F$ 为 $T$ 上支集在 $Y$ 上固有的凝聚层; 要说明 $\widehat X$ 上每个支集在 $\widehat Y$ 上固有的凝聚层都可代数化. 主要工具是 Chow 引理 [EGAIII1, 5.6.1]: 假定 $X$ 非空, 可以找到态射
\[Z\xrightarrow{g}X\xrightarrow{f}Y,\]
使 $g$ 射影且满射, $fg$ 拟射影, 并且存在非空的开浸入 $j:U\to X$, 使 $g$ 在 $U$ 上诱导同构. 令 $T=X-U$ 赋予约化概形结构, $J$ 为 $T$ 在 $X$ 中的理想. 设 $E$ 为 $\widehat X$ 上支集在 $\widehat Y$ 上固有的凝聚层. 考虑正合列
\begin{equation}0\longrightarrow K\longrightarrow E\longrightarrow\widehat g_*\widehat g^*E\longrightarrow C\longrightarrow0.\tag{*}\label{ch08:eq843star}\end{equation}
% Source PDF 218 / printed p.206.
只需说明以下各点:
\begin{enumerate}
\item[(1)] $\widehat g_*\widehat g^*E$ 可代数化.
\item[(2)] $K,C$ 被 $\widehat J$ 的某个正次幂 $\widehat J^N$ 消灭, 因而可视为 $\widehat{T'}$ 上的凝聚层, 其中 $T'$ 为 $J^N$ 定义的 $T$ 的加厚.
\item[(3)] $K,C$ 作为 $\widehat{T'}$ 上的凝聚层可代数化.
\item[(4)] 对 $\widehat X$ 上支集在 $\widehat Y$ 上固有的凝聚层, 可代数化的性质对核、余核与扩张稳定.
\end{enumerate}
对 (1), 注意由情形 (b), $\widehat g^*E$ (其支集在 $\widehat Y$ 上固有, 因为 $g$ 固有) 可代数化. 随后由比较定理 8.2.2 得到 $\widehat g_*\widehat g^*E$ 可代数化.

为证明 (2), 可在 $\widehat X$ 上局部工作. 可以把 $X$ 换成 $\Spec B$, 其中 $B$ 是进 Noether 环且 $IB$ 是 $B$ 的定义理想. 此时对 $\Spec B$ 上的某个凝聚层 $F$, 有 $E=\widehat F$; 由 8.2.9, (2) 来自 $F\to g_*g^*F$ 的核与余核被 $J$ 的某个正次幂消灭这一事实.

根据 (2), (3) 来自 Noether 归纳假设.

在 (4) 中, 对核与余核的稳定性是直接的, 对扩张的稳定性来自 $r=1$ 时的 8.2.9.

这就完成了 8.4.2 的证明.

\paragraph{8.4.4.}
首先给出 8.4.2 对闭形式子概形、有限形式概形以及形式概形之间态射的代数化的应用. 我们需要一些定义.

(a) \emph{闭形式子概形.} 设 $\mathcal X$ 是局部 Noether 形式概形. 若 $\mathcal A$ 是 $\mathcal O_{\mathcal X}$ 的凝聚理想, 则由 $\mathcal O_{\mathcal X}/\mathcal A$ 的支集 $\mathcal Y$ (它是 $\mathcal X$ 的闭子集) 以及限制到 $\mathcal Y$ 的环层 $\mathcal O_{\mathcal X}/\mathcal A$ 组成的带拓扑赋环空间 $\mathcal Y$, 是一个在 $\mathcal X$ 上进的局部 Noether 形式概形 (8.1.5), 称为 $\mathcal A$ 定义的 $\mathcal X$ 的\emph{闭形式子概形}. 若 $\mathcal I$ 是 $\mathcal X$ 的一个定义理想 (8.1.3), $X_n=(\mathcal X,\mathcal O_{\mathcal X}/\mathcal I^{n+1})$, 于是 $\mathcal X=\operatorname{colim}_nX_n$, 则 $\mathcal Y=\operatorname{colim}_nY_n$, 其中 $Y_n$ 是 $X_n$ 的闭子概形, 满足 $\mathcal O_{Y_n}=\mathcal O_{X_n}\otimes_{\mathcal O_{\mathcal X}}\mathcal O_{\mathcal Y}$. 反过来, 任意归纳系态射 $Y_\bullet\to X_\bullet$, 若 $Y_n\to X_n$ 为闭子概形且 $Y_n=X_n\times_{X_{n+1}}Y_{n+1}$ (参见 8.1.5), 则定义 $\mathcal X$ 的一个闭形式子概形 $\mathcal Y=\operatorname{colim}_nY_n$, 满足 $X_n\times_{\mathcal X}\mathcal Y=Y_n$. 若 $\mathcal X$ 仿射, $\mathcal X=\operatorname{Spf}A$, 则对 $A$ 的一个理想 $\mathfrak a$, 有 $\mathcal A=\mathfrak a^\triangle$, $\mathcal Y=\operatorname{Spf}(A/\mathfrak a)$. 最后, 若 $X$ 是局部 Noether 概形, $\widehat X$ 是其沿一个闭子概形 $X_0$ 的完备化, 而 $Y$ 是 $X$ 的闭子概形, $\widehat Y$ 是其沿 $Y_0=X_0\times_XY$ 的完备化, 则 $\widehat Y$ 为 $\widehat X$ 的闭形式子概形.

(b) \emph{有限态射.} 设 $\mathcal X=\operatorname{colim}X_n$ 是如 (a) 中所述的局部 Noether 形式概形. 若局部 Noether 形式概形态射 $f:\mathcal Z\to\mathcal X$ 为进态射 (8.1.5), 且 $f_0:Z_0\to X_0$ 有限, 则称 $f$ \emph{有限} [EGAIII1, 4.8.2]. 由标准交换代数 [Bou61, III §8.2, 11] (或 [EGAI, Ch.~0, 7.2.9]), 这等价于局部地 $f$ 形如 $\operatorname{Spf}(B)\to\operatorname{Spf}(A)$, 其中 $B$ 在 $A$ 上有限且为 $IB$ 进的, $I$ 是 $A$ 的定义理想; 也等价于 $f$ 为进态射且每个 $f_n:Z_n=X_n\times_{\mathcal X}\mathcal Z\to X_n$ 都有限. 若 $f$ 有限, 则 $f_*\mathcal O_{\mathcal Z}$ 是一个有限 $\mathcal O_{\mathcal X}$ 代数 $\mathcal B$, 使对所有 $n$, $\mathcal O_{X_n}\otimes_{\mathcal O_{\mathcal X}}\mathcal B=f_*\mathcal O_{X_n}$. 若 $\mathcal X=\widehat X$, $X$ 如 (a) 中所述, $Z$ 是 $X$ 上的有限概形, 则 $Z$ 沿 $Z_0=X_0\times_XZ$ 的完备化 $\widehat Z$ 在 $\widehat X$ 上有限.
% Original prints f_*O_{X_n}; retained for review (expected Z_n).

% Source PDF 219 / printed p.207.
\begin{origstatement}{推论}{8.4.5}\label{ch08:cor845}
设 $X/Y$ 如 8.4.2 中所述. 则 $Z\mapsto\widehat Z$ 给出从 $X$ 的在 $Y$ 上固有的闭子概形集合到 $\widehat X$ 的在 $\widehat Y$ 上固有的闭形式子概形集合的双射 (8.4.4 (a)).
\end{origstatement}
非平凡之处是满射性. 设 $\mathcal Z=\operatorname{colim}Z_n$ 是 $\widehat X$ 的一个在 $\widehat Y$ 上固有的闭形式子概形. 它对应于 $\mathcal O_{\widehat X}$ 的一个凝聚商 $\mathcal O_{\mathcal Z}$, 其支集在 $\widehat Y$ 上固有. 由 8.4.2, 存在唯一凝聚 $\mathcal O_X$ 模 $F$, 使 $\widehat F=\mathcal O_{\mathcal Z}$. 问题在于代数化满射 $u:\mathcal O_{\widehat X}\to\mathcal O_{\mathcal Z}$. 不能应用 8.4.2, 因为 $\mathcal O_X$ 的支集不一定在 $Y$ 上固有. 不过, $F$ 的支集, 即 $\mathcal O_X$ 与 $F$ 的支集之交, 在 $Y$ 上固有. 由 8.2.9, 这足以保证映射 $\Hom(\mathcal O_X,F)\to\Hom(\mathcal O_{\widehat X},\mathcal O_{\mathcal Z})$ 双射. 因而存在唯一 $v:\mathcal O_X\to F$, 使 $\widehat v=u$. 因为 $v_0=u_0$ 满射, 所以 $v$ 也满射, 因此 $F=\mathcal O_Z$, 其中 $Z$ 是 $X$ 的一个在 $Y$ 上固有且满足 $\widehat Z=\mathcal Z$ 的闭子概形.

\begin{origstatement}{推论}{8.4.6}\label{ch08:cor846}
设 $X/Y$ 如 8.4.2 中所述. 则 $Z\mapsto\widehat Z$ 给出从在 $Y$ 上固有的有限 $X$ 概形范畴到在 $\widehat Y$ 上固有的有限 $\widehat X$ 形式概形范畴的等价 (8.4.4 (b)).
\end{origstatement}
通过 $Z\mapsto g_*\mathcal O_Z$ (相应地, $\mathcal Z\mapsto g_*\mathcal O_{\mathcal Z}$), 其中 $g$ 为结构态射, 第一个 (相应地, 第二个) 范畴反等价于有限且支集在 $Y$ (相应地, $\widehat Y$) 上固有的 $\mathcal O_X$ (相应地, $\mathcal O_{\widehat X}$) 代数的范畴. 若 $A,B$ 是支集在 $Y$ 上固有的有限 $\mathcal O_X$ 代数, $u:A\to B$ 是一个 $\mathcal O_X$ 模映射, 满足 $\widehat u$ 是 $\mathcal O_{\widehat X}$ 代数映射, 则由 8.4.2, $u$ 自动是 $\mathcal O_X$ 代数映射. 于是得到全忠实性. 若 $\mathcal A$ 是支集在 $\widehat Y$ 上固有的有限 $\mathcal O_{\widehat X}$ 代数, 则由 8.4.2, 存在支集在 $Y$ 上固有的凝聚 $\mathcal O_X$ 模 $A$, 使 $\widehat A=\mathcal A$ 作为 $\mathcal O_{\widehat X}$ 模成立. 再由 8.4.2, 给出 $\mathcal A$ 上代数结构的映射 $\mathcal A\otimes\mathcal A\to\mathcal A$ 与 $\mathcal O_{\widehat X}\to\mathcal A$, 唯一地代数化为赋予 $A$ 一个 $\mathcal O_X$ 代数结构的映射, 且作为 $\mathcal O_{\widehat X}$ 代数有 $\widehat A=\mathcal A$.

\begin{origstatement}{推论}{8.4.7}\label{ch08:cor847}
设 $X$ 是固有 $Y$ 概形, $Z$ 是 $Y$ 上分离且有限型的 Noether 概形. 则映射
\[\Hom_Y(X,Z)\longrightarrow\Hom_{\widehat Y}(\widehat X,\widehat Z),\qquad f\longmapsto\widehat f\]
双射. 特别地, 从固有 $Y$ 概形范畴到 $\widehat Y$ 形式概形范畴的函子 $X\mapsto\widehat X$ 全忠实.
\end{origstatement}
若 $\widehat f=\widehat g$, 则关于 (8.1.6.4) 的核的说明表明 (参见 [EGAI, 10.9.4]), 在 $X_0$ 的某个邻域中 $f=g$; 因为 $X\to Y$ 固有 (特别地, 是闭映射), 且 $I$ 包含在 $A$ 的根基中, 所以处处有 $f=g$. 为证明满射性, 对给定态射 $\widehat X\to\widehat Z$ 的图应用 8.4.5, 把图视为 $(X\times_YZ)^\wedge$ 的闭形式子概形.

\begin{origstatement}{注}{8.4.8}\label{ch08:rem848}
若在 8.4.7 中去掉 $X$ 固有的假设, 则结论一般不再成立. 例如, 若 $X=Z=\Spec A[t]$, 则 $\widehat X=\widehat Z=\operatorname{Spf}A\{t\}$, $\Hom_Y(X,Z)=A[t]$, 而
\[\Hom_{\widehat Y}(\widehat X,\widehat Z)=\Hom_{A,\mathrm{cont}}(A\{t\},A\{t\})=A\{t\},\]
其中 $A\{t\}=A[t]^\wedge$ 是限制形式幂级数 $\sum a_nt^n$ 的环, 即当 $n$ 趋于无穷时, $a_n$ 在 $I$ 进拓扑中趋于 $0$.
\end{origstatement}

% Source PDF 220 / printed p.208.
\paragraph{8.4.9.}
若 $X$ 是固有 $Y$ 概形, 则 $\widehat X$ 是在 $\widehat Y$ 上固有的 Noether 进 $\widehat Y$ 形式概形 (这里意指 $X_0=X\times_YY_0$ 在 $Y_0=\Spec(A/I)$ 上固有). 若 $\mathcal X$ 是一个固有进 $\widehat Y$ 形式概形, 且 $\mathcal X$ 可代数化, 即对某个固有 $Y$ 概形 $X$ 有 $\mathcal X\simeq\widehat X$, 则由 8.4.7, $X$ 唯一 (相差一个在 $\widehat X$ 上诱导恒等的唯一同构). 形变理论能构造出不可代数化的固有进 $\widehat Y$ 形式概形, 参见 8.5.24 (b). 然而, 下面这个极为有用的结果表明, 在射影形式情形这不可能发生.

\begin{origstatement}{定理}{8.4.10}\label{ch08:th8410}
[EGAIII1, 5.4.5] 设 $\mathcal X=\operatorname{colim}X_n$ 是一个固有进 $\widehat Y$ 形式概形, 其中 $X_n=\mathcal X\times_{\widehat Y}Y_n$. 设 $L$ 是一个可逆 $\mathcal O_{\mathcal X}$ 模, 满足 $L_0=L\otimes\mathcal O_{X_0}=L/IL$ 丰沛 (所以 $X_0$ 在 $Y_0$ 上射影). 则 $\mathcal X$ 可代数化; 若 $X$ 是满足 $\widehat X=\mathcal X$ 的固有 $Y$ 概形, 则存在唯一的 $X$ 上的线丛 $M$, 使 $L=\widehat M$, 并且 $M$ 丰沛 (特别地, $X$ 在 $Y$ 上射影).
\end{origstatement}
利用 8.4.3.2, 选取 $n$ 使:
\begin{enumerate}
\item[(i)] $L_0^{\otimes n}$ 很丰沛, 即对某个标准射影空间 $P_0=\mathbb P_{Y_0}^r$ 及闭浸入 $i_0:X_0\to P_0$, 有 $L_0^{\otimes n}=i_0^*\mathcal O_{P_0}(1)$;
\item[(ii)] $\Gamma(\mathcal X,L^{\otimes n})\to\Gamma(X_0,L_0^{\otimes n})$ 满射.
\end{enumerate}
利用 (ii), 将 $i_0$ 给出的典范满射 $u_0:\mathcal O_{X_0}^{r+1}\to L_0^{\otimes n}$ 提升为一个 $\mathcal O_{\mathcal X}$ 线性映射 $u:\mathcal O_{\mathcal X}^{r+1}\to L^{\otimes n}$. 由 Nakayama 引理, 每个 $u_k=u\otimes\mathcal O_{X_k}:\mathcal O_{X_k}^{r+1}\to L_k^{\otimes n}$ ($k\in\mathbb N$) 都满射, 因而对应于 $Y_k$ 概形态射 $i_k:X_k\to P_k=\mathbb P_{Y_k}^r$, 使 $L_k^{\otimes n}=i_k^*\mathcal O_{P_k}(1)$. 由 8.4.4 (b), $i_k$ 有限, 故由 Nakayama 引理为闭浸入. 这些闭浸入 $i_k$ 组成归纳系 $i_\bullet:X_\bullet\to P_\bullet$, 并具有形如 (8.1.5.2) 的笛卡尔方块, 因而定义一个闭形式子概形 (8.4.4 (a)) $i:\mathcal X\to\widehat P$, 其中 $\widehat P$ 为 $Y=\Spec A$ 上标准射影空间 $P=\mathbb P_Y^r$ 的完备化, 且 $L^{\otimes n}=i^*\mathcal O_{\widehat P}(1)$. 由 8.4.5, 存在 (唯一的) 闭子概形 $j:X\to P$, 使 $\widehat X=\mathcal X$. 此外, 由 8.4.2, 存在 $X$ 上 (唯一的) 线丛 $M$, 使 $L=\widehat M$. 因为 $L^{\otimes n}=i^*\mathcal O_{\widehat P}(1)$ 且 $(M^{\otimes n})^\wedge=\widehat M^{\otimes n}$, 得到 $(M^{\otimes n})^\wedge=(j^*\mathcal O_P(1))^\wedge$, 因而 (由 8.4.2) $M^{\otimes n}=j^*\mathcal O_P(1)$, 所以 $M$ 丰沛.

\begin{origstatement}{注}{8.4.11}\label{ch08:rem8411}
\begin{enumerate}
\item[(a)] §8.2、§8.4 中的主要定理与 Serre 关于代数几何和解析几何之间比较的结果 (GAGA) 类似. 见 [Ser56] 与 [SGA1, XII].
\item[(b)] Knutson 已将 §8.2、§8.4 的结果推广到代数空间 [Knu71, chap.~V]. 最近也已得到对叠 (Deligne-Mumford 叠或 Artin 叠 [Art74b]、[LMB00]) 的推广. Zariski 主定理 (8.2.17) 的推广见 [LMB00, 16.5]. Faltings [Fal03] 和 Olsson [Ols] 独立地给出了固有态射有限性定理到 Artin 叠的推广; 后者从 Artin 叠的一个 Chow 引理导出该结果. Olsson 在同一文献中还将 Grothendieck 的主要存在定理 8.4.2 推广到 Artin 叠.
\end{enumerate}
\end{origstatement}

\section{对提升问题的应用}\label{ch08:sec85}
\paragraph{8.5.1.}
设 $A$ 为局部 Noether 环, 极大理想为 $\mathfrak m$, 剩余域为 $k=A/\mathfrak m$. 令 $S=\Spec A$, 其闭点为 $s=\Spec k$. 提升问题的一个原型如下.
% Source PDF 221 / printed p.209.
给定 $s$ 上的有限型概形 $X_0$, 能否找到一个在 $S$ 上有限型且\emph{平坦}的概形 $X$ 来提升 $X_0$, 即使 $X_s\simeq X_0$? 例如, $A$ 可以是一个混合特征离散赋值环, $k$ 的特征为 $p>0$ 而分式域 $K$ 的特征为零, 于是如上的概形 $X$ 就提供了 ``$X_0$ 到特征零的提升'' (即一般纤维 $X_\eta$, $\eta=\Spec K$). 通常 $X_0$ 还满足一些附加假设 (例如固有、光滑等), 有时还带有附加结构 (例如群结构), 这些都要求在提升中保留. 为简单起见, 这里略去这些要求. Grothendieck 处理这一问题的策略包含几个步骤.

(1) 尝试将 $X_0$ 提升为一些 (平坦且有限型的) 概形 $X_n$ 的归纳系, 使 $X_{n+1}\times_{S_{n+1}}S_n=X_n$. 闭浸入 $S_n\to S_{n+1}$ 是一个 1 阶加厚 (8.1.3): 其理想 $\mathfrak m^{n+1}\mathcal O_{S_{n+1}}$ 被 $\mathfrak m$ 消灭, 因而更是平方为零的. 假定对 $m\leq n$ 已构造 $X_m$. 要将 $X_n$ 提升为 $S_{n+1}$ 上的 $X_{n+1}$, 通常会在 $X_0$ 的一个上同调群中遇到障碍; 当该障碍消失时, 这样的 $X_{n+1}$ 的同构类集合与 $X_0$ 的另一个上同调群一一对应. 在 $X_n$ 上诱导恒等的给定 $X_{n+1}$ 的自同构, 也可以由 $X_0$ 的一个适当上同调群描述. 这种研究是\emph{形变理论}的课题.

(2) 假定已经找到了如 (1) 中的归纳系 $X_\bullet$. 它定义 $S$ 在 $s$ 处的完备化 $\widehat S=\operatorname{Spf}\widehat A$ 上的一个进 (局部 Noether) 形式概形 $\mathcal X$, 它 (按定义) 在 $\widehat S$ 上平坦且有限型 (8.1.5). 下一个问题是把 $\mathcal X$ 在 $\Spec\widehat A$ 上代数化, 即寻找一个 $\Spec\widehat A$ 上的有限型概形 $X$, 使 $\widehat X=\mathcal X$. 这里可尝试应用 §4 的存在定理, 假定 $X_0$ 固有; 此时若代数化存在, 它就是唯一的. 主要工具 (即使不说是唯一的工具) 是 8.4.10. 为此, 首先必须假定 $X_0$ 射影. 设 $L_0$ 是 $X_0$ 上的丰沛可逆层. 若能选取这样的 $L_0$, 使它可提升到 $\mathcal X$, 即各 $X_n$ 上存在可逆层的射影系 $L_n$, 满足 $L_n=L_{n+1}\otimes\mathcal O_{X_n}$, 则由 8.4.10, 问题已获解决: 存在 $\Spec\widehat A$ 上的射影概形 $X$, 使 $\widehat X=\mathcal X$. 假定对 $m\leq n$ 已构造 $L_m$, 则将 $L_n$ 提升到 $L_{n+1}$ 有一个上同调障碍, 它类似于 (1) 中提到的障碍且与之密切相关.
% The source prints §4 here; retained rather than silently changed to §8.4.

(3) 找到 $\Spec\widehat A$ 上的 $X$ 后, 一般不能再进一步, 即不能把 $X$ 下降到 $S=\Spec A$. 但有时在模问题中, $X/\Spec\widehat A$ 具有万有形变性质. 若此外 $k$ 可分闭, 则 Artin 的\emph{逼近理论} [Art69a]、[Art74b] 通常允许把 $X$ 至少下降到 $S$ 的 hensel 化 $S^h$, 即找到 $S^h$ 上的 $Z$, 使 $X=\Spec\widehat A\times_{S^h}Z$. 在某种意义上, Artin 理论回答了 Grothendieck 在 [Gro95a, p.~15] 中提出的问题: ``要从关于一个局部环的完备化的已知结果, 过渡到关于这个局部环本身的相应结果, 就需要第四个‘基本定理’, 其最终表述尚待找到.''

本节回顾形变理论的基本事实, 并给出与 (1)、(2) 有关的问题的应用.

\subsection*{A. 向量丛的形变}
\paragraph{8.5.2.}
最简单的形变问题是\emph{向量丛}, 即有限秩局部自由层的形变问题. 设 $i:X_0\to X$ 是一个 1 阶加厚 (8.1.3), 由平方为零的理想 $I$ 定义. 设 $E_0$ 是 $X_0$ 上的向量丛. 我们希望把 $E_0$ ``形变'' (或 ``延拓'') 到 $X$ 上, 即找到 $X$ 上的向量丛 $E$,
% Source PDF 222 / printed p.210.
使 $\mathcal O_{X_0}\otimes E=E/IE=E_0$. 更精确地, $E_0$ 在 $X$ 上的一个\emph{形变}, 是指由 $X$ 上的向量丛 $E$ 及一个 $\mathcal O_X$ 线性映射 $E\to i_*E_0$ 组成的二元组, 此映射诱导同构 $i^*E\xrightarrow\sim E_0$. 所谓形变的\emph{态射} $u:E'\to E$, 是指满足 $i^*u=\id_{E_0}$ 的态射 $u$. 这样的态射自动是同构.

\begin{origstatement}{定理}{8.5.3}\label{ch08:th853}
设 $i:X_0\to X$ 如 8.5.2 中所述.

(a) 设 $E,F$ 是 $X$ 上的向量丛, $E_0=i^*E$, $F_0=i^*F$, $u_0:E_0\to F_0$ 为 $\mathcal O_{X_0}$ 线性映射. 延拓 $u_0$ 的 $\mathcal O_X$ 线性映射 $u:E\to F$ 的存在性有一个障碍
\[o(u_0,i)\in H^1(X_0,I\otimes\mathcal Hom(E_0,F_0)).\]
当 $o(u_0,i)=0$ 时, 延拓 $u_0$ 的 $u$ 的集合是 $H^0(X_0,I\otimes\mathcal Hom(E_0,F_0))$ 下的仿射空间.

(b) 设 $E_0$ 是 $X_0$ 上的向量丛. 存在一个障碍
\[o(E_0,i)\in H^2(X_0,I\otimes\mathcal End(E_0)),\]
其消失是 $E_0$ 在 $X$ 上存在形变 $E$ 的充要条件. 当 $o(E_0,i)=0$ 时, $E_0$ 在 $X$ 上的形变集合是 $H^1(X_0,\mathcal End(E_0)\otimes I)$ 下的仿射空间, 而给定形变 $E$ 的自同构群通过 $a\mapsto a-\id$ 与 $H^0(X_0,\mathcal End(E_0)\otimes I)$ 等同.
\end{origstatement}
证明是初等的, 在更一般的背景下 (赋环空间或拓扑斯) 也适用. 首先证明 (a). 第二个断言是清楚的. 而且, $u_0$ 的延拓 $u$ 局部存在. 因而得到 $X_0$ 上 $I\otimes\mathcal Hom(E_0,F_0)$ 下的一个挠子 $P$, 其在 $X_0$ 的开子集 $U$ 上的截面是 $u_0|U$ 的 $\mathcal O_X$ 线性延拓. $P$ 在 $H^1(X_0,I\otimes\mathcal Hom(E_0,F_0))$ 中的类就是障碍 $o(u_0,i)$.

为证明 (b), 首先为简单起见, 假定 $X_0$ (或等价地, $X$) 分离. 选取 $(\mathcal U=(U_i)_{i\in K},(E_i)_{i\in K})$, 其中 $\mathcal U$ 是 $X_0$ 的仿射开覆盖, $E_i$ 是 $E_0|U_i$ 到 $X\cap U_i$ 的一个形变. 因为 $X_0$ 分离, $U_{ij}=U_i\cap U_j$ 仿射, 所以由 (a) 可找到同构 $g_{ij}:E_i|U_{ij}\xrightarrow\sim E_j|U_{ij}$ (在 $X_0$ 上诱导恒等). 这样的同构是唯一的, 相差加上一个 $h_{ij}\in H^0(U_{ij},I\otimes\mathcal End(E_0))$. 于是
\[(i,j,k)\longmapsto c_{ijk}=g_{ij}g_{ik}^{-1}g_{jk}-\id\in H^0(U_{ijk},I\otimes\mathcal End(E_0))\]
是 $\mathcal U$ 的以 $I\otimes\mathcal End(E_0)$ 为值的一个 2 余循环, 它是余边界 $dh$, $h=(h_{ij})$, 当且仅当可以修改 $g_{ij}$, 使其成为 $(E_i)$ 的黏合资料, 换言之, 当且仅当 $E_0$ 可在 $X$ 上形变. 因此, $c$ 的类
\[[c]=o(E_0,i)\in H^2(X_0,I\otimes\mathcal End(E_0))\]
就是所求障碍; 可以验证它不依赖于各项选择. 若 $E_1,E_2$ 是 $E_0$ 在 $X$ 上的两个形变, 则由 (a), 从 $E_1$ 到 $E_2$ 的局部同构组成一个 $I\otimes\mathcal End(E_0)$ 下的挠子, 其类
\[[E_2]-[E_1]\in H^1(X_0,I\otimes\mathcal End(E_0))\]
仅依赖于 $E_1,E_2$ 的同构类 $[E_i]$, 并且为零当且仅当 $[E_1]=[E_2]$. 可以验证, 这就在形变的同构类集合上定义了所求的仿射结构. 这完成了 $X_0$ 分离时的证明. 在一般情形, 资料 $(U_i),(E_i),(g_{ij})$ 必须换成资料 $(U_i),(E_i),(g_{ij}^{\alpha})$, 其中对 $U_{ij}$ 的一个开覆盖 $(U_{ij}^{\alpha})_{\alpha\in A_{ij}}$, $g_{ij}^{\alpha}$ 是从 $E_i|U_{ij}^{\alpha}$ 到 $E_j|U_{ij}^{\alpha}$ 的同构. 此时各 $g_{ij}^{\alpha}$ 给出由 $(U_i),(U_{ij}^{\alpha})$ 定义的超覆盖的一个 2 余循环,
% Source PDF 223 / printed p.211.
其在 $H^2(X_0,I\otimes\mathcal End(E_0))$ 中的上同调类就是所求的障碍 (参见 [SGA4, V 7]); 证明的其余部分只需略作修改.

用 Giraud 的 gerbe 语言 [Gir71] 可更内在地呈现这一证明. $E_0$ 在 $X_0$ 的可变开子集 $U$ 上的形变 $U\mapsto\mathcal E(U)$ 组成一个群胚叠, 实际上是一个 gerbe, 即具有下列性质: $\mathcal E(U)$ 中任意两个对象局部同构, 且对任意 $U$, 存在 $U$ 的开覆盖 $(U_i)$, 使 $\mathcal E(U_i)$ 非空. $\mathcal E(U)$ 中对象的自同构层组成 $X_0$ 上的一个整体层, 即 $I\otimes\mathcal End(E_0)$, 称为 gerbe $\mathcal E$ 的\emph{带 (band, ``lien'')}. 障碍 $o(E_0,i)$ 是 $\mathcal E$ 的\emph{上同调类}. 当此类为零时, gerbe 是\emph{中性的}, 这意味着选定一个整体对象 $E$ ($E_0$ 在 $X$ 上的一个形变) 后, 就可将 $\mathcal E$ 与 $I\otimes\mathcal End(E_0)$ 下的挠子的 gerbe 等同 (在 $X_0$ 的可变开子集上). 关于把上述讨论推广到以光滑群概形 $G$ 取代 $\operatorname{GL}(n)$ (以 $G$ 下的挠子取代秩 $n$ 局部自由层) 的情形, 见 [Gir71, VII 1.3.1].

\begin{origstatement}{注}{8.5.4}\label{ch08:rem854}
\begin{enumerate}
\item[(a)] 8.5.3 (b) 中余循环 $c$ 的构造表明, 若 $L_0,M_0$ 是 $X_0$ 上的线丛, 则在 $H^2(X_0,I)$ 中有
\[o(L_0\otimes M_0,i)=o(L_0,i)+o(M_0,i).\]
因此在线丛上, 障碍的行为类似于第一 Chern 类. 事实上, 8.5.3 (b) 中的类 $o(E_0,i)$ 可视为一种 Atiyah 类, 类似于 Atiyah 在 [Ati57] 中为构造 Hodge 上同调中的 Chern 类而定义的类; 见 [Ill71, chap.~IV, V] 与 [KS04, 1.4.1].
\item[(b)] 实际中, 理想 $I$ 被一个更大的理想 $J$ 消灭. 更确切地, 改变一下记号, 设 $X_0\longrightarrow X_1\xrightarrow{i_1}X_2$ 为闭浸入, $I$ (相应地, $J$) 为 $X_1$ (相应地, $X_0$) 在 $X_2$ 中的理想, 假定 $I\cdot J=0$. 特别地, $I^2=0$, 而 $I$ 不仅可视为 $\mathcal O_{X_1}$ 模, 还可视为 $\mathcal O_{X_0}(=\mathcal O_{X_2}/J)$ 模. 对 $X_2$ 上的向量丛 $E,F$, 8.5.3 (a) 中出现的群 $H^q(X_1,I\otimes\mathcal Hom(E_1,F_1))$, 将 $i$ 换成 $i_1$, 并置 $E_1=i_1^*E$, $F_1=i_1^*F$, 可重写为
\[\begin{aligned}
H^q(X_1,I\otimes_{\mathcal O_{X_1}}\mathcal Hom(E_1,F_1))
=H^q(X_0,I\otimes_{\mathcal O_{X_0}}\mathcal Hom(E_0,F_0)),
\end{aligned}\]
其中 $E_0=\mathcal O_{X_0}\otimes E$, $F_0=\mathcal O_{X_0}\otimes F$. 类似地, 对 $X_1$ 上的向量丛 $E_1$, 8.5.3 (b) 中出现的群 $H^q(X_1,I\otimes_{\mathcal O_{X_1}}\mathcal End(E_1))$ (把 $i$ 换成 $i_1$), 可重写为
\[H^q(X_1,I\otimes_{\mathcal O_{X_1}}\mathcal End(E_1))=H^q(X_0,I\otimes_{\mathcal O_{X_0}}\mathcal End(E_0)),\]
其中 $E_0=\mathcal O_{X_0}\otimes E_1$.
\end{enumerate}
\end{origstatement}

\begin{origstatement}{推论}{8.5.5}\label{ch08:cor855}
设 $A$ 是完备局部 Noether 环, 极大理想为 $\mathfrak m$, 剩余域为 $k$. 令 $S=\Spec A$, $\widehat S=\operatorname{colim}S_n$, 其中 $S_n=\Spec A/\mathfrak m^{n+1}$. 设 $\mathcal X=\operatorname{colim}X_n$, $X_n=S_n\times_{\widehat S}\mathcal X$, 为 $\widehat S$ 上平坦的进局部 Noether 形式概形 (8.1.5), 假定 $H^2(X_0,\mathcal O_{X_0})=0$. 则 $X_0$ 上任意线丛 $L_0$ 均可提升到 $\mathcal X$ 上的线丛 $L$. 若此外 $H^1(X_0,\mathcal O_{X_0})=0$, 则这样的提升 $L$ 唯一, 相差一个 (在 $L_0$ 上诱导恒等的非唯一) 同构.
\end{origstatement}
% Source PDF 224 / printed p.212.
假定 $L_0$ 已提升到 $X_n$ 上的 $L_n$. 令 $I_n$ 为 $X_n$ 在 $X_{n+1}$ 中的理想. 由 $\mathcal X$ 在 $\widehat S$ 上平坦, 有
\[I_n=\mathcal O_{X_0}\otimes_k\mathfrak m^{n+1}/\mathfrak m^{n+2}.\]
结合 8.5.4 (b), 可知把 $L_n$ 提升到 $X_{n+1}$ 上的 $L_{n+1}$ 的障碍位于
\[H^2(X_0,I_n)=H^2(X_0,\mathcal O_{X_0})\otimes_k\mathfrak m^{n+1}/\mathfrak m^{n+2}=0,\]
于是得到第一个断言. 对第二个断言, 设 $L,L'$ 是 $L_0$ 在 $\mathcal X$ 上的两个提升. 假定对 $m\leq n$ 已构造同构 $u_m:L_m\xrightarrow\sim L'_m$, 且 $u_0=\id$. 因为 $H^1(X_0,\mathcal O_{X_0})=0$, 由 8.5.3 (a), 存在延拓 $u_n$ 的同构 $u_{n+1}:L_{n+1}\xrightarrow\sim L'_{n+1}$; $u=\lim u_n$ 就是从 $L$ 到 $L'$ 的在 $L_0$ 上诱导恒等的同构.

\begin{origstatement}{推论}{8.5.6}\label{ch08:cor856}
设 $\mathcal X$ 为 $\widehat S$ 上固有、平坦的进局部 Noether 形式概形. 则:

(a) 若 $X/S$ 为固有概形且 $\mathcal X=\widehat X$, 则 $X$ 在 $S$ 上平坦. 此外, 若 $H^2(X_0,\mathcal O_{X_0})=0$, 则 $X_0$ 上的任意线丛 $L_0$ 均可提升到 $X$ 上的线丛 $L$; 若 $H^1(X_0,\mathcal O_{X_0})=0$, 则此提升唯一 (相差一个同构).

(b) 若 $X_0$ 射影, 且 $X_0$ 上的某个丰沛线丛 $L_0$ 可提升到 $\mathcal X$ 上的线丛 $\mathcal L$, 则存在射影且平坦的概形 $X/S$, 使 $\mathcal X=\widehat X$, 以及 $X$ 上的丰沛线丛 $L$, 使 $\widehat L=\mathcal L$.
\end{origstatement}
先证明 (a). 设 $x$ 是 $X_0=X_s$ 的一点 ($s=S_0=\Spec k$). 对所有 $n\geq0$, $\mathcal O_{X,x}/\mathfrak m^{n+1}\mathcal O_{X,x}=\mathcal O_{X_n,x}$ 在 $A_n=A/\mathfrak m^{n+1}$ 上平坦, 故由通常的平坦性判别法, $\mathcal O_{X,x}$ 在 $A$ 上平坦. 因为态射平坦的点组成开集, 所以 $X$ 在特殊纤维 $X_s$ 的某个邻域中在 $S$ 上平坦; 又因为 $X$ 在 $S$ 上固有, 所以处处平坦. 第二个断言来自 8.4.2 与 8.5.5. 断言 (b) 来自 (a) 与 8.4.10.

\subsection*{B. 光滑概形的形变}
\paragraph{8.5.7.}
现在转向概形的形变问题. 设 $i:S_0\to S$ 是由平方为零的理想 $I$ 定义的 1 阶加厚 (8.1.3), $X_0$ 是 $S_0$ 上的平坦概形. 所谓 $X_0$ 在 $S$ 上的\emph{形变} (或\emph{提升}), 是指一个笛卡尔方块
\begin{equation}
\begin{tikzcd}X_0\arrow[r,"j"]\arrow[d]&X\arrow[d]\\S_0\arrow[r]&S\end{tikzcd}
\tag{8.5.7.1}\label{ch08:eq8571}\end{equation}
其中 $X$ 在 $S$ 上平坦. 平坦性条件可表达为自然映射
\begin{equation}f_0^*I\longrightarrow J,\tag{8.5.7.2}\label{ch08:eq8572}\end{equation}
是同构, 其中 $f_0:X_0\to S_0$ 为结构态射, $J$ 为 $X_0$ 在 $X$ 中的理想. 所谓形变的态射, 是指满足 $uj=j'$ 的 $S$ 态射 $u:X\to X'$ (其中 $j':X_0\to X'$). 这样的态射 $u$ 必为同构.

% Source PDF 225 / printed p.213.
\paragraph{8.5.8.}
先讨论初等的光滑情形. 设 $f:X\to Y$ 是概形态射. 回顾: 若 $f$ 局部有限呈示 (当 $Y$ 局部 Noether 时即为局部有限型), 且满足下列等价条件, 则称 $f$ \emph{光滑}:
\begin{enumerate}
\item[(i)] $f$ 平坦, 且 $X/Y$ 的几何纤维 $X_{\bar y}$ 正则 (这里 $\bar y\to y\in Y$ 遍历 $Y$ 的几何点, $k(\bar y)$ 代数闭);
\item[(ii)] (Jacobi 判别法) 对每个 $x\in X$, 存在 $x$ 的仿射开邻域 $U$ 及 $y=f(x)$ 的仿射开邻域 $V$, 使 $f(U)\subset V$, 且 $U$ 是标准仿射空间 $\mathbb A_V^n=\Spec A[t_1,\ldots,t_n]$ (其中 $V=\Spec A$) 中由方程 $g_1=\cdots=g_r=0$ ($g_i\in A[t_1,\ldots,t_n]$) 定义的闭子概形, 满足 $\operatorname{rk}(\partial g_i/\partial t_j)(x)=r$;
\item[(iii)] (形式光滑性) 对每个交换方块
\[\begin{tikzcd}S_0\arrow[r,"g_0"]\arrow[d,"i"']&X\arrow[d,"f"]\\S\arrow[r,"h"']&Y\end{tikzcd}\]
其中 $i$ 为 1 阶加厚, 在 $S$ 上 Zariski 局部存在延拓 $g_0$ 的 $Y$ 态射 $g:S\to X$, 即 $gi=g_0$.
\end{enumerate}
(这些条件的等价性以及光滑、平展态射的基本事实, 见 [BLR90]、[Ill72a]; 更全面的处理见 [EGAIV2, §17]、[SGA1, I, II, III].)

假定 $f$ 光滑. 则相对微分层 $\Omega^1_{X/Y}$ 是有限型局部自由层, 切层
\[T_{X/Y}=\mathcal Hom(\Omega^1_{X/Y},\mathcal O_X)\]
也是如此. 它们在 $X$ 的一点 $x$ 处的共同秩 $r(x)$, 等于纤维 $X_{f(x)}$ 在 $x$ 处的维数, 即 $X$ 在 $x$ 处的\emph{相对维数}. 它是 $x$ 的局部常值函数. 若态射 $f:X\to Y$ 光滑且在每一点相对维数为零, 则称 $f$ \emph{平展}; 换言之, $f$ 光滑且 $\Omega^1_{X/Y}=0$, 或等价地, $f$ 平坦、局部有限呈示且 $\Omega^1_{X/Y}=0$. 光滑性 (相应地, 平展性) 对复合与基变换稳定.

我们还需要形式概形背景下的光滑性定义. 设 $\mathcal Y=\operatorname{colim}Y_n$ 是局部 Noether 形式概形, 使用 8.1.5 的记号. 若进态射 $f:\mathcal X\to\mathcal Y$ 平坦 (8.1.5), 且每个 $X_n$ 都在 $Y_n$ 上光滑 (或等价地, 由 8.5.8 (i), $X_0$ 在 $Y_0$ 上光滑), 则称 $f$ \emph{光滑}. 我们将称 $\mathcal X$ 为\emph{提升 $X_0$ 的 $\mathcal Y$ 上光滑形式概形}.

关于光滑概形形变的主要结果概括于下述定理中.

\begin{origstatement}{定理}{8.5.9}\label{ch08:th859}
\begin{enumerate}
\item[(a)] 设 $X,Y$ 是概形 $S$ 上的概形, $Y$ 在 $S$ 上光滑, $j:X_0\to X$ 为一个由平方为零的理想 $J$ 定义的闭子概形. 设 $g:X_0\to Y$ 为 $S$ 态射. 存在一个障碍
\[o(g,j)\in H^1(X_0,J\otimes_{\mathcal O_{X_0}}g^*T_{Y/S}),\]
其消失是存在延拓 $g$ 的 $S$ 态射 $h:X\to Y$ (即 $hj=g$) 的充要条件. 当 $o(g,j)=0$ 时, $g$ 的各延拓 $h$ 的集合是 $H^0(X_0,J\otimes_{\mathcal O_{X_0}}g^*T_{Y/S})$ 下的仿射空间.
% Source PDF 226 / printed p.214.
\item[(b)] 设 $i:S_0\to S$ 为一个由平方为零的理想 $I$ 定义的 1 阶加厚, $X_0$ 为光滑 $S_0$ 概形. 存在一个障碍
\[o(X_0,i)\in H^2(X_0,f_0^*I\otimes T_{X_0/S_0})\]
(其中 $f_0:X_0\to S_0$ 为结构态射), 其消失是 $X_0$ 在 $S$ 上存在形变 $X$ 的充要条件 (8.5.7). 当 $o(X_0,i)=0$ 时, 这样的形变的同构类集合是 $H^1(X_0,f_0^*I\otimes T_{X_0/S_0})$ 下的仿射空间, 一个固定形变的自同构群同构于 $H^0(X_0,f_0^*I\otimes T_{X_0/S_0})$. 特别地, 若 $X_0$ 在 $S_0$ 上平展, 则存在 $X_0$ 在 $S$ 上的形变 $X$, 且它唯一, 相差一个唯一同构.
\end{enumerate}
\end{origstatement}
注意, 若 $X_0$ 光滑 (相应地, 平展), 则 $X_0$ 在 $S$ 上的任意形变都光滑 (相应地, 平展). 这来自 8.5.8 (i).

8.5.9 的证明类似于 8.5.3. 首先证明 (a). 因为 $Y$ 在 $S$ 上光滑, $g$ 的延拓 $h$ 在 $X_0$ 上局部存在. 而且, 两个这样的延拓之差是一个从 $\mathcal O_Y$ 到 $g_*J$ 的 $S$ 导子, 即 $J\otimes_{\mathcal O_{X_0}}g^*T_{Y/S}$ 的一个截面. 因此, 在 $X$ 的可变开子集上的延拓 $h$ 组成 $X$ 上 $J\otimes_{\mathcal O_{X_0}}g^*T_{Y/S}$ 下的一个挠子, $o(g,j)$ 就是该挠子的类. 为证明 (b), 首先注意到 $X_0$ 的形变在 $X_0$ 上局部存在. 这来自 8.5.8 (ii) (提升各多项式 $g_i$). 而且, (a) 蕴含任意两个形变局部同构, 并且对 $X_0$ 的任意开子集 $U_0$, $U_0$ 的一个形变 $U$ 的自同构层通过 $a\mapsto a-\id$ 与 $f_0^*I\otimes T_{X_0/S_0}$ 等同. 因此 (参见 [Gir71, VII 1.2]), 把 $X_0$ 的每个开子集 $U_0$ 对应到 $U_0$ 的形变群胚, 就定义一个 gerbe $\mathcal G=\mathcal G_{X_0}$, 其带为 $f_0^*I\otimes T_{X_0/S_0}$. $\mathcal G$ 在 $H^2(X_0,f_0^*I\otimes T_{X_0/S_0})$ 中的类 $o(X_0,i)$, 是 $\mathcal G(X_0)$ 中存在对象 (即 $X_0$ 的一个形变) 的障碍. 当 $o(X_0,i)=0$ 时, $\mathcal G$ 中性, 即存在 $X_0$ 的一个 (整体) 形变 $X$. 选定这样的 $X$ 后, 把 $X_0$ 的开子集 $U_0$ 的形变 $U$ 对应到 $U$ 与 $X|U_0$ 之间局部同构的挠子, 就可将 $\mathcal G$ 与 $X_0$ 上 $f_0^*I\otimes T_{X_0/S_0}$ 下的挠子 gerbe 等同. 特别地, 此时 $X_0$ 的形变同构类集合与 $H^1(X_0,f_0^*I\otimes T_{X_0/S_0})$ 等同.

如 8.5.3 的证明一样, 可以写出定义 $o(X_0,i)$ 的一个 2 余循环. 为简单起见, 假定 $X_0$ 分离. 选取 $(\mathcal U=((U_0)_i)_{i\in K},(U_i)_{i\in K})$, 其中 $\mathcal U$ 是 $X_0$ 的仿射开覆盖, $U_i$ 是 $(U_0)_i$ 的一个形变. 因为 $X_0$ 分离, $(U_0)_{ij}=(U_0)_i\cap(U_0)_j$ 仿射, 所以由 (a) 存在形变的同构 $g_{ij}:U_i|(U_0)_{ij}\xrightarrow\sim U_j|(U_0)_{ij}$. 于是
\[(i,j,k)\longmapsto c_{ijk}=g_{ij}g_{ik}^{-1}g_{jk}-\id\in H^0((U_0)_{ijk},f_0^*I\otimes T_{X_0/S_0})\]
为 $\mathcal U$ 的以 $f_0^*I\otimes T_{X_0/S_0}$ 为值的 2 余循环, 其在 $H^2(X_0,f_0^*I\otimes T_{X_0/S_0})$ 中的类表示 $o(X_0,i)$.

\begin{origstatement}{注}{8.5.10}\label{ch08:rem8510}
\begin{enumerate}
\item[(a)] 障碍 $o(X_0,i)$ 满足如下函子性. 设 $g_0:Y_0\to S_0$ 是光滑态射, $h_0:X_0\to Y_0$ 是 $S_0$ 态射 (于是 $f_0=g_0h_0:X_0\to S_0$). 则 $o(X_0,i)$ 与 $o(Y_0,i)$ 在典范映射
\[\begin{aligned}
H^2(X_0,f_0^*I\otimes T_{X_0/S_0})&\longrightarrow H^2(X_0,f_0^*I\otimes h_0^*T_{Y_0/S_0})\\
&\longleftarrow H^2(Y_0,g_0^*I\otimes T_{Y_0/S_0})
\end{aligned}\]
之下有相同的像.
% Source PDF 227 / printed p.215.
此外, 若 $X_0,Y_0$ 是光滑 $S_0$ 概形, 则将 $X_0\times_{S_0}Y_0$ 提升到 $S$ 的障碍 $o(X_0\times_{S_0}Y_0,i)$ 满足公式
\[o(X_0\times_{S_0}Y_0,i)=\operatorname{pr}_1^*o(X_0,i)+\operatorname{pr}_2^*o(Y_0,i),\]
其中 $\operatorname{pr}_1$ (相应地, $\operatorname{pr}_2$) 是从 $X_0\times_{S_0}Y_0$ 到 $X_0$ (相应地, $Y_0$) 的投影, $\operatorname{pr}_1^*$ 是函子性映射与第一个直和项的包含映射的复合
\[\begin{aligned}
H^2(X_0,I\otimes T_{X_0/S_0})&\longrightarrow H^2(X_0\times_{S_0}Y_0,I\otimes\operatorname{pr}_1^*T_{X_0/S_0})\\
&\longrightarrow H^2(X_0\times_{S_0}Y_0,I\otimes T_{X_0\times_{S_0}Y_0}),
\end{aligned}\]
$\operatorname{pr}_2^*$ 类似.

障碍 $o(g,j)$ 对态射复合满足一种相容性: 在 8.5.9 (b) 的情形, 若 $X,Y,Z$ 是 $S$ 上的光滑概形, $f_0:X_0\to Y_0$, $g_0:Y_0\to Z_0$ 是它们拉回到 $S_0$ 后所得概形之间的 $S_0$ 态射, 则将 $h_0=g_0f_0$ 提升为 $h:X\to Z$ 的障碍, 是将 $g_0$ 提升为 $g:Y\to Z$ 的障碍经 $g_0^*$ 的拉回.
% Original prints g_0^* in the final sentence; retained.
\item[(b)] 如 8.5.4 (b), 假定 $S_0\longrightarrow S_1\xrightarrow{i_1}S_2$ 是闭浸入, 其中 $i_1$ 的理想 $I$ 被 $S_0$ 在 $S_2$ 中的理想 $J$ 消灭. 则若 $f_1:X_1\to S_1$ 是光滑态射, 8.5.9 (b) 中与 $X_1$ 在 $S_2$ 上的形变有关的群可重写为 $H^q(X_0,f_0^*I\otimes T_{X_0/S_0})$, 其中 $f_0:X_0\to S_0$ 由 $f_1$ 经基变换得到.
\end{enumerate}
\end{origstatement}

\subsection*{C. 基本群的特殊化}
\paragraph{8.5.11.}
存在定理 8.4.2 与 8.5.3、8.5.9 的结合具有有力的应用. 先讨论与\emph{基本群}有关的应用.

设 $X$ 是局部 Noether 概形. $X$ 的一个\emph{平展覆盖} (``revêtement étale'' [SGA1, I 4.9]) 是指一个有限且平展的态射 $Y\to X$. 平展覆盖的态射 $Y'\to Y$ 定义为从 $Y'$ 到 $Y$ 的 $X$ 态射. 它自动是 $Y$ 的平展覆盖. 用
\begin{equation}\operatorname{Et}(X)\tag{8.5.11.1}\label{ch08:eq85111}\end{equation}
记 $X$ 的平展覆盖范畴. 假定 $X$ 连通, 并固定 $X$ 的一个\emph{几何点} $\bar x$, 位于某点 $x$ 上, 即态射 $\Spec k(\bar x)\to\Spec k(x)$, 其中 $k(\bar x)$ 为可分闭域. 此时定义了一个 profinite 群
\begin{equation}\pi_1(X,\bar x),\tag{8.5.11.2}\label{ch08:eq85112}\end{equation}
称为 $X$ 在 $\bar x$ 处的\emph{基本群}, 以及范畴等价
\begin{equation}\operatorname{Et}(X)\xrightarrow\sim\{\pi_1(X,\bar x)\text{--}fsets\},\tag{8.5.11.3}\label{ch08:eq85113}\end{equation}
其中 $\{\pi_1(X,\bar x)\text{--}fsets\}$ 表示带有 $\pi_1(X,\bar x)$ 连续作用的有限集范畴 [SGA1, V 7]. 更确切地, 函子
\[\operatorname{Et}(X)\longrightarrow\{fsets\},\qquad Y\longmapsto Y(\bar x)=Y_{\bar x},\]
将 $X$ 的平展覆盖 $Y$ 对应到它在 $\bar x$ 之上的点的有限集合, 称为 $\bar x$ \emph{处的纤维函子}, 它是 pro 可表示的: 存在 $\operatorname{Et}(X)$ 的一个 pro 对象 $P=(P_i)_{i\in I}$, 称为 $X$ 的\emph{万有 (pro) 平展覆盖}, 以及同构
\begin{equation}\Hom(P,Y)=\operatorname{colim}_i\Hom(P_i,Y)\xrightarrow\sim Y(\bar x)\tag{8.5.11.4}\label{ch08:eq85114}\end{equation}
% Source PDF 228 / printed p.216.
它对 $Y\in\operatorname{Et}(X)$ 具有函子性. $P$ 的恒等映射通过 (8.5.11.4) 对应于一点 $\xi\in P(\bar x)=\lim P_i(\bar x)$, 反过来, 该点通过 $(u:P\to Y)\mapsto u(\xi)\in Y(\bar x)$ 定义 (8.5.11.4). 那些为 Galois 的 $P_i$, 即连通、非空且使自然映射 $\Aut(P_i)=\Hom(P_i,P_i)\to\Hom(P,P_i)(\simeq P_i(\bar x))$ 双射的 $P_i$, 组成一个共尾系, 因而有
\[\Hom(P,P)=\Aut(P)=\lim_{i\in J}\Aut(P_i),\]
其中 $J$ 是 $I$ 中使 $P_i$ 为 Galois 的指标 $i$ 组成的子集. 按定义, $P$ 的自同构群 $\Aut(P)$ 的反群就是 $\pi_1(X,\bar x)$. 换言之, 它是 $\bar x$ 处纤维函子的自同构群. 它连续且函子性地 (从左边) 作用在 $Y(\bar x)$ 上, 这定义了等价 (8.5.11.3). 平展覆盖 $Y$ 连通当且仅当 $\pi_1(X,\bar x)$ 在 $Y(\bar x)$ 上传递作用.

若 $\bar a\to X$, $\bar b\to X$ 为两个几何点, 则因 $X$ 连通, $\bar a$ 处纤维函子 $F_{\bar a}$ 与 $\bar b$ 处纤维函子 $F_{\bar b}$ 同构 [SGA1, V 5.6]. 选定从 $F_{\bar a}$ 到 $F_{\bar b}$ 的一个同构, 称为从 $\bar a$ 到 $\bar b$ 的一条\emph{道路}. 这样的道路诱导同构
\begin{equation}\pi_1(X,\bar a)\xrightarrow\sim\pi_1(X,\bar b).\tag{8.5.11.5}\label{ch08:eq85115}\end{equation}
若不假定 $X$ 连通, 则将 $X$ 在 $\bar x$ 处的基本群定义为包含 $x$ 的连通分支的基本群. 基本群以自然方式成为带几何基点的局部 Noether 概形上的函子. 若 $f:X\to Z$ 是连通局部 Noether 概形之间的态射, 则逆像函子
\[f^*:\operatorname{Et}(Z)\longrightarrow\operatorname{Et}(X),\qquad Z'\longmapsto X\times_ZZ'\]
为等价当且仅当同态
\[f_*:\pi_1(X,\bar x)\longrightarrow\pi_1(Z,\bar z)\]
为同构, 其中 $\bar z$ 是 $\bar x$ 经 $f$ 所得到的几何点 $\bar x\to x\to z$. 同态 $f_*$ 满射当且仅当函子 $f^*$ 全忠实; 等价地, 对 $Z$ 的任意连通平展覆盖 $Z'$, $f^*Z'$ 连通 [SGA1, V 6.9]. 它单射当且仅当对 $X$ 的任意平展覆盖 $X'$, 存在 $Z$ 的平展覆盖 $Z'$ 以及从 $f^*Z'$ 的一个连通分支到 $X'$ 的映射 [SGA1, V 6.8].

下面的结果补充了 8.2.18.
\begin{origstatement}{定理}{8.5.12}\label{ch08:th8512}
[SGA1, IX 1.10] 设 $A$ 是完备局部 Noether 环, 极大理想为 $\mathfrak m$, 剩余域为 $k$. 令 $S=\Spec A$, $\widehat S=\operatorname{Spf}A=\operatorname{colim}S_n$, 其中 $S_n=\Spec A/\mathfrak m^{n+1}$. 设 $X$ 是 $S$ 上的固有概形. 则逆像函子
\[\operatorname{Et}(X)\longrightarrow\operatorname{Et}(X_s),\]
其中 $s=S_0=\Spec k$, 是一个等价. 换言之, 对 $X_s$ 的任意几何点 $\bar x$, 自然同态
\[\pi_1(X_s,\bar x)\longrightarrow\pi_1(X,\bar x)\]
是同构.
\end{origstatement}
% Source PDF 229 / printed p.217.
设 $\widehat X$ 是 $X$ 沿 $X_s$ 的形式完备化, 于是 $\widehat X=\operatorname{colim}_nX_n$, 其中 $X_n=S_n\times_SX$, $S_n=\Spec A/\mathfrak m^{n+1}$. 考虑自然态射
\[X_s\xrightarrow{i}\widehat X\xrightarrow{j}X.\]
有逆像函子
\[\operatorname{Et}(X)\xrightarrow{j^*}\operatorname{Et}(\widehat X)\xrightarrow{i^*}\operatorname{Et}(X_s),\]
其中 $\operatorname{Et}(\widehat X)$ 表示 $\widehat X$ 的平展覆盖范畴, 即 $\widehat X$ 上平展的有限形式概形 $\mathcal Y=\operatorname{colim}Y_n$ 的范畴 (8.4.4 (b)); 这里平展是说对所有 $n\geq0$, $Y_n$ 在 $X_n$ 上平展. 由 8.5.9 (b), $i^*$ 是等价. 另一方面, 由 8.4.6, $(ji)^*$ 全忠实. 还需证明 $j^*$ 本质满射. 设 $\mathcal Y$ 是 $\widehat X$ 的一个平展覆盖. 由 8.4.6, 存在唯一的 $X$ 上有限的概形 $Y$, 使 $\widehat Y=\mathcal Y$. 若 $y$ 是 $Y_s$ 的一点, $n\geq0$, 则 $\mathcal O_{Y,y}/\mathfrak m^{n+1}\mathcal O_{Y,y}=\mathcal O_{Y_n,y}$ 在 $\mathcal O_{X,x}/\mathfrak m^{n+1}\mathcal O_{X,x}=\mathcal O_{X_n,x}$ 上平坦, 其中 $x$ 是 $y$ 在 $X_s$ 中的像. 因而 $Y$ 在 $Y_s$ 的某个邻域中在 $X$ 上平坦; 又因 $Y$ 在 $S$ 上固有, 所以 $Y$ 在 $X$ 上平坦. 此外,
\[j^*\Omega^1_{Y/X}=\lim_n\Omega^1_{Y_n/X_n}=0,\]
因为 $Y_n$ 在 $X_n$ 上平展. 因而由 8.4.2, $\Omega^1_{Y/X}=0$, 所以 $Y$ 在 $X$ 上平展 (8.5.8).

\begin{origstatement}{注}{8.5.13}\label{ch08:rem8513}
由 Artin 逼近定理可知, 若只假定 $A$ 为 hensel 环而不要求完备, 8.5.12 的结论仍成立; 见 [SGA4-1/2, IV 2.2]. 断言 8.2.18 与 8.5.12 在平展上同调的\emph{固有基变换定理}的证明中至关重要 (同上及 [SGA4, XII]).
\end{origstatement}

\paragraph{8.5.14.}
定理 8.5.12 是 Grothendieck 的基本群特殊化理论 [SGA1, X] 的出发点. 设 $f:X\to Y$ 是局部 Noether 概形的固有态射, 其几何纤维连通. 设 $s,\eta$ 是 $Y$ 的点, 满足 $s\in\overline{\{\eta\}}$, $\bar s$ (相应地, $\bar\eta$) 是 $s$ (相应地, $\eta$) 之上的几何点, $a$ (相应地, $b$) 为 $X_{\bar s}$ (相应地, $X_{\bar\eta}$) 的几何点. 此时定义了一个同态 (同上 2.1, 2.4)
\begin{equation}\pi_1(X_{\bar\eta},b)\longrightarrow\pi_1(X_{\bar s},a),\tag{8.5.14.1}\label{ch08:eq85141}\end{equation}
称为\emph{特殊化同态}. 这个同态定义良好, 相差一个目标群的内自同构. 若 $Y$ 是 hensel 局部 Noether 环 $A$ 的谱, 闭点为 $s$ 且 $s=\bar s$, 则 (8.5.14.1) 是复合
\[\pi_1(X_{\bar\eta},b)\longrightarrow\pi_1(X,b)\xrightarrow\sim\pi_1(X,a)\xrightarrow\sim\pi_1(X_s,a),\]
其中第一个映射为函子性映射, 第二个为对应于从 $a$ 到 $b$ 的道路的同构 (8.5.11) (这样的道路存在, 因为由 8.2.18, 假设蕴含 $X$ 连通), 最后一个为 8.5.12、8.5.13 中同构的逆. 一般情形的定义较为细致, 见上述文献. 它使用如下事实: 对代数闭域 $k$ 上的固有连通概形 $X$, $X$ 的基本群在 $k$ 的代数闭扩张下不变 (这是 8.5.12 的一个非平凡推论). Grothendieck 关于 (8.5.14.1) 的主要结果如下:
% Source PDF 230 / printed p.218.
\begin{origstatement}{定理}{8.5.15}\label{ch08:th8515}
[SGA1, X 2.4, 3.8] (Grothendieck 特殊化定理) 设 $f:X\to Y$ 如 8.5.14 中所述.

(a) 若 $f$ 平坦且具有几何约化的纤维 (即对任意态射 $\bar y\to y\in Y$, 其中 $\bar y$ 为代数闭域的谱, $X_{\bar y}$ 都约化), 则 (8.5.14.1) 满射;

(b) 若 $f$ 光滑, $p$ 是 $s$ 的特征指数, 则 (8.5.14.1) 在基本群的最大素于 $p$ 的商上诱导同构
\[\pi_1^{(p')}(X_{\bar\eta},b)\xrightarrow\sim\pi_1^{(p')}(X_{\bar s},a).\]
\end{origstatement}
(我们用记号 $\pi_1^{(p')}$ 表示最大素于 $p$ 的商; 这一记号如今比上述文献中使用的记号 $\pi_1^{(p)}$ 更为常见.)

先在 $Y$ 为 hensel 局部 Noether 环 $A$ 的谱、剩余域 $k$ 代数闭且 $s=\Spec k$ 的情形证明 (a) (一般情形可归结到此). 需说明: 若 $Z$ 是 $X$ 的连通平展覆盖, 则 $Z_{\bar\eta}$ 连通. 注意, $Z$ 也在 $Y$ 上固有且平坦, 纤维几何约化. 因为 $Z$ 连通, 由 8.2.18, 特殊纤维 $Z_s$ 也连通. 所以 $H^0(Z_s,\mathcal O_{Z_s})$ 是剩余域为 $k$ 的 artin 局部 $k$ 代数. 因为 $Z_s$ 约化, $H^0(Z_s,\mathcal O_{Z_s})=k$. 典范映射的复合
\[\mathcal O_Y\otimes k\longrightarrow g_*\mathcal O_Z\otimes k\longrightarrow H^0(Z_s,\mathcal O_{Z_s})=k,\]
其中 $g:Z\to Y$ 是结构态射, 是恒等映射, 特别地,
\[g_*\mathcal O_Z\otimes k\longrightarrow H^0(Z_s,\mathcal O_{Z_s})\]
满射. 因为 $Z$ 在 $Y$ 上平坦, 由 8.3.11, 这个映射实际上是同构, 且 $g_*\mathcal O_Z$ 为秩 1 自由模, 它的形成与任意基变换交换; 换言之, $g_*\mathcal O_Z=\mathcal O_Y$ \emph{普遍地}成立 (即在任意基变换后成立). 特别地, $Z_{\bar\eta}$ 连通.

再概述 (b) 的证明, 其中 $Y$ 如上述 (a) 的证明, 并进一步假定 $A$ 是一致化参数为 $\pi$ 的离散赋值环 (细节见 [SGA1, X 3], 综述见 [OV00]). 由 (a), 限制函子 $\operatorname{Et}(X)\to\operatorname{Et}(X_{\bar\eta})$ 全忠实. 因而只需说明 $X_{\bar\eta}$ 的群为 $G$ 的任意 Galois 平展覆盖 $V$ (其中 $G$ 为阶素于 $p$ 的有限群), 都由 $X$ 的一个 $G$ Galois 平展覆盖诱导. 利用 8.5.12, 将 $Y$ 换成某个有限扩张后, 可假定 $V$ 是由 $X_\eta$ 的一个 $G$ 平展覆盖 $V_0$ 经基变换得到的. 令 $n$ 为 $G$ 的阶, $Y'=\Spec A'$, 其中 $A'=A[t]/(t^n-\pi)$; 令 $X'=Y'\times_YX$, $V'=\eta'\times_\eta V$ ($\eta'$ 为 $Y'$ 的一般点), 即由 $X$ 与 $V_0$ 经基变换得到的概形. 设 $R$ 为 $X$ 在 $X_s$ 的一般点 $\zeta$ 处的局部环. 它是以 $\pi$ 为一致化参数的离散赋值环. $X'$ 在 $\zeta$ 处的局部环为 $R'=R[t]/(t^n-\pi)$. 由 Abhyankar 引理, $V'$ 在 $\Spec R'$ 的一般点上的限制可延拓为 $\Spec R'$ 的平展 $G$ 覆盖, 进而延拓为 $X'$ 的某开子集 $U'$ 的平展 $G$ 覆盖 $W'$, 其补集是 $X_s$ 的一个在 $X_s$ 中余维至少为 1 的闭子集, 因而在 $X'$ 中余维至少为 2. 因为 $X'$ 在 $Y'$ 上光滑, 所以正则, Zariski--Nagata 纯性定理蕴含 $W'$ 可延拓为 $X'$ 的一个平展 $G$ 覆盖 $Z'$. 由 8.5.12, $Z'$ 来自 $X$ 的一个平展 $G$ 覆盖 $Z$, 证明完成.
% The source prints V'=eta' times_eta V, although V_0 was just introduced; retained.

(a) 的证明中概述的论证实际上给出如下结果 ([SGA1, 1.2]、[EGAIII2, 7.8.6]):
% Source PDF 231 / printed p.219.
\begin{origstatement}{命题}{8.5.16}\label{ch08:prop8516}
设 $f:X\to Y$ 是局部 Noether 概形的固有平坦态射, 其纤维几何约化, 设
\[X\longrightarrow Y'\longrightarrow Y\]
是其 Stein 分解 (8.2.11). 则 $Y'$ 是 $Y$ 的一个平展覆盖, 且它的形成与任意基变换交换. 特别地, $f$ 在次数零上同调平坦 (8.3.10), 而且下列条件等价:
\begin{enumerate}\item[(i)] $f_*\mathcal O_X=\mathcal O_Y$;\item[(ii)] $f$ 的几何纤维连通.\end{enumerate}
\end{origstatement}

\begin{origstatement}{注}{8.5.17}\label{ch08:rem8517}
\begin{enumerate}
\item[(a)] 在 8.5.15 (a) 的假设下, 即 $f:X\to Y$ 固有且平坦, 并具有几何约化且连通的纤维时, 特殊化同态 (8.5.14.1) 在近几年中得到广泛研究, 特别是在相对曲线的情形. 关于其中某些方面的讨论, 见 [BLR00].
\item[(b)] Fujiwara--Kato [FK95] 已在 ``对数几何'' 中构造了基本群理论的一个变体. 导论见 [Ill02]; 关于 Grothendieck 特殊化定理 8.5.15 在此背景下的推广, 以及对局部域上簇的素于 $p$ 的基本群的野惯性外自同构作用的应用 [Kis00], 见 [Vid01, I 2.2].
\end{enumerate}
\end{origstatement}

\subsection*{D. 曲线}
\paragraph{8.5.18.}
现在讨论对曲线提升的应用. 设 $Y$ 为局部 Noether 概形. 所谓 $Y$ 上的\emph{曲线}, 是指一个在 $Y$ 上平坦、分离且有限型, 相对维数为 1 的态射 $f:X\to Y$. 假定 $f$ 固有. 则对 $X$ 上的任意凝聚层 $F$, 由 8.2.10, 当 $q>1$ 时 $R^qf_*F=0$; 若此外 $F$ 在 $Y$ 上平坦, 例如 $F=\mathcal O_X$, 则复形 $Rf_*F$ 完美, 其完美幅度在 $[0,1]$ 中 (8.3.11). 简单例子表明 [Har77, III 12.9.2], 一般地 $f$ 在次数 0 和 1 都不必上同调平坦. 然而, 若 $f$ 的纤维几何约化, 则由 8.5.16, $f$ 在次数 0 上同调平坦, 从而由 8.3.11 在次数 1 也上同调平坦, 即对所有 $q$, $R^qf_*\mathcal O_X$ 为有限型局部自由模. 若此外 $f$ 的几何纤维连通, 从而 $f_*\mathcal O_X=\mathcal O_Y$, 则局部自由层 $R^1f_*\mathcal O_X$ 的秩称为 $Y$ 上曲线 $X$ 的\emph{(算术) 亏格}. 若 $f$ 固有且光滑, 则由 Grothendieck 对偶定理, 各层 $R^qf_*\Omega^1_{X/Y}$ 也是有限型局部自由层, 并定义了迹映射
\[\operatorname{Tr}:R^1f_*\Omega^1_{X/Y}\longrightarrow\mathcal O_Y,\]
而由到 $R^1f_*\Omega^1_{X/Y}$ 的自然配对与 $\operatorname{Tr}$ 复合而得的配对
\[R^qf_*\mathcal O_X\otimes R^{1-q}f_*\Omega^1_{X/Y}\longrightarrow\mathcal O_Y\]
是有限型局部自由层之间的完美配对. 特别地, 若 $f$ 具有连通几何纤维且亏格为 $g$, 则 $f_*\Omega^1_{X/Y}$ 为秩 $g$ 的局部自由层. 最后回顾: 域上的任何曲线都拟射影. 对曲线为代数闭域上固有曲线的情形, 见 [Har77, III Ex.~5.8] (一般情形可归结到此).

关于曲线提升的主要结果为下述定理:
% Source PDF 232 / printed p.220.
\begin{origstatement}{定理}{8.5.19}\label{ch08:th8519}
[SGA1, III 7.3] 设 $A$ 为完备局部 Noether 环, 剩余域为 $k$. 令 $S=\Spec A$, $s=\Spec k$, 设 $X_0$ 是 $s$ 上的射影光滑概形, 满足
\begin{equation}H^2(X_0,T_{X_0/s})=0.\tag{i}\label{ch08:eq8519i}\end{equation}
则存在 $\widehat S$ 上提升 $X_0$ 的固有光滑形式概形 $\mathcal X$ (8.5.8). 若 $X_0$ 除了满足 (i), 还满足
\begin{equation}H^2(X_0,\mathcal O_{X_0})=0,\tag{ii}\label{ch08:eq8519ii}\end{equation}
则存在 $S$ 上的射影光滑概形 $X$, 使 $X_s=X_0$.
\end{origstatement}
例如, 当 $X_0$ 是 $s$ 上的固有光滑曲线时, 条件 (i)、(ii) 均满足. 注意, 若 $X_0$ 是亏格为 $g$ 的固有、几何连通光滑曲线, 则 $X$ 在 $S$ 上的各纤维也有同样性质.

令 $\widehat S=\operatorname{Spf}A=\operatorname{colim}S_n$, 其中 $S_n=\Spec A/\mathfrak m^{n+1}$, $\mathfrak m$ 表示 $A$ 的极大理想. 先说明在假设 (i) 下, 存在 $\widehat S$ 上提升 $X_0$ 的 (固有) 光滑形式概形 $\mathcal X=\operatorname{colim}_nX_n$. 假定对 $m\leq n$, 已构造在 $S_m$ 上光滑的 $X_m$, 满足 $X_m=S_m\times_{S_n}X_n$; 令 $i_n:S_n\to S_{n+1}$ 为包含映射. 则由 8.5.9、8.5.10 (b), $X_n$ 在 $S_{n+1}$ 上存在光滑提升 $X_{n+1}$ 的障碍为
\[o(X_n,i_n)\in H^2(X_0,T_{X_0/s}\otimes\mathfrak m^{n+1}/\mathfrak m^{n+2}).\]
但
\[H^2(X_0,T_{X_0/s}\otimes\mathfrak m^{n+1}/\mathfrak m^{n+2})=H^2(X_0,T_{X_0/s})\otimes\mathfrak m^{n+1}/\mathfrak m^{n+2},\]
由 (i), 该群为零. 这证明了 $\mathcal X$ 的存在性. 对 8.5.19 的第二个断言, 由 8.5.6 得到存在 $S$ 上的射影平坦概形 $X$, 使 $X_s=X_0$. 此时 $X$ 在 $X_s$ 的每一点都在 $S$ 上光滑, 从而在 $X_s$ 的某个开邻域中光滑; 又因 $X$ 在 $S$ 上固有, 该邻域必须等于 $X$.

由 8.5.19, \emph{正特征固有光滑曲线可 ``无分歧地'' 提升到特征零}: 若 $k$ 的特征为 $p>0$, 可取 $A$ 为 $k$ 的 Cohen 环, 即剩余域为 $k$、分式域特征为零、极大理想由 $p$ 生成的完备离散赋值环 (当 $k$ 完美时, 就是 $k$ 上的 Witt 向量环 $W(k)$). 利用这一点及紧 Riemann 曲面 (拓扑) 基本群的已知结构, Grothendieck 从特殊化定理 8.5.15 导出了关于正特征固有光滑曲线的 (代数) 基本群的如下结果:

\begin{origstatement}{定理}{8.5.20}\label{ch08:th8520}
[SGA1, X, 3.9, 3.10] 设 $k$ 为特征指数为 $p$ 的代数闭域, $C$ 是 $k$ 上亏格为 $g$ 的固有、光滑连通曲线. 设 $x$ 为 $C$ 的有理点. 以 $\Pi_g$ 记由生成元 $a_i,b_i$ ($1\leq i\leq g$) 和关系 $\prod_{1\leq i\leq g}(a_i,b_i)=1$ 定义的群, 其中 $(a,b):=aba^{-1}b^{-1}$, 以 $\widehat\Pi_g$ 记其 profinite 完备化. 则存在一个满同态
\[\widehat\Pi_g\longrightarrow\pi_1(C,x),\]
它在最大素于 $p$ 的商上诱导同构
\[\widehat\Pi_g^{(p')}\xrightarrow\sim\pi_1(C,x)^{(p')}.\]
\end{origstatement}
% Source PDF 233 / printed p.221.
下面概述论证. 首先处理 $k=\mathbb C$ 的情形. 设 $C^{\mathrm{an}}$ 为对应于 $C$ 的 (紧、连通且亏格为 $g$ 的) Riemann 曲面. 由 Riemann 存在定理, 从 $C$ 的有限平展覆盖范畴到 $C^{\mathrm{an}}$ 的有限平展覆盖范畴的函子 $C'\mapsto C'^{\mathrm{an}}$ 为等价. 因而 $\pi_1(C,x)$ 是 $\pi_1(C^{\mathrm{an}},x)$ 的 profinite 完备化. 将 $C^{\mathrm{an}}$ 表示为具有 $4g$ 条边 $(e_i,e_i^{-1},f_i,f_i^{-1})$ ($1\leq i\leq g$) 的多边形按字 $\prod_{1\leq i\leq g}(e_i,f_i)$ 所指定的黏合所得的商, 拓扑论证表明 $\pi_1(C^{\mathrm{an}},x)=\Pi_g$. 因而结果在此情形成立. 当 $p=1$ 时, 利用 (固有概形的) (代数) 基本群在代数闭域的任意扩张下不变, 通过标准极限论证归结到这一情形. 最后, 设 $p\geq2$. 在 8.5.19 中, 取 $X_0=C$, $A=W(k)$ 为 $k$ 上的 Witt 向量环. 设 $X$ 为 $S$ 上的射影光滑概形, 使 $X_s=C$. 则由 8.5.18, $X/S$ 是具有亏格 $g$ 的连通几何纤维的射影光滑曲线. 于是结论来自 $p=1$ 的情形与 8.5.15.

\begin{origstatement}{注}{8.5.21}\label{ch08:rem8521}
\begin{enumerate}
\item[(a)] 若 $g=0$, 则 $C$ 同构于 $\mathbb P_k^1$, 因而单连通 (由 Riemann--Hurwitz). 更一般地, 所有 $\mathbb P_k^r$ 都单连通 [SGA1, XI 1.1].
\item[(b)] 若 $g=1$, 则 $C$ 为\emph{椭圆曲线}, $\pi_1(C)$ 是 $C$ 的 \emph{Tate 模}, $T(C)=\lim_n{}_nC(k)$, 其中 $n$ 遍历所有 $\geq1$ 的整数, ${}_nC(k)$ 表示 $C(k)$ 上乘 $n$ 的核, 当 $m=nd$ 时, ${}_mC(k)$ 通过乘 $d$ 映到 ${}_nC(k)$. 更一般地, 若 $A$ 是 $k$ 上的 Abel 簇, 则
\[\pi_1(A)=T(A),\]
其中 $T(-)$ 为类似定义的 Tate 模 [Mum70, IV 18].
\item[(c)] 由 8.5.20, $\pi_1(C,x)$ 为拓扑有限型. 如 Grothendieck 在 [SGA1, X 2.8] 中所指出的, $\pi_1(C,x)$ 似乎不太可能拓扑有限呈示, 但这个问题仍然未决. 使用超平面截面的某些 Lefschetz 型论证, Grothendieck 证明, 更一般地, 对 $k$ 上任何固有连通概形 $X$, $\pi_1(X)$ 都拓扑有限生成 [SGA1, X, 2.9].
\item[(d)] 8.5.20 的最后一个断言对仿射曲线也有一个变体. 更确切地说, 设 $C$ 为 $k$ 上亏格 $g\geq0$ 的固有光滑连通曲线. 设整数 $n\geq1$, $x_1,\ldots,x_n$ 为 $C$ 的互异有理点, 令 $X=C-\{x_1,\ldots,x_n\}$, 并选取 $X$ 的一个有理点 $x$. 则有同构
\[\widehat\Pi_{g,n}^{(p')}\xrightarrow\sim\pi_1(X,x)^{(p')},\]
其中 $\widehat\Pi_{g,n}^{(p')}$ 是 (自由) 群 $\Pi_{g,n}$ 的 profinite 完备化的素于 $p$ 的商, $\Pi_{g,n}$ 由生成元 $a_i,b_i$ ($1\leq i\leq g$)、$s_i$ ($1\leq i\leq n$) 及关系 $\prod_{1\leq i\leq g}(a_i,b_i)\prod_{1\leq i\leq n}s_i=1$ 定义 [SGA1, XIII 2.12]. 然而, 即使对 $X=\mathbb A_k^1$, $\pi_1(X,x)$ 也不是拓扑有限型. 有限群 $G$ 是 $\mathbb A_k^1$ 的一个连通平展覆盖的 Galois 群, 当且仅当它的最大素于 $p$ 的商平凡 (Abhyankar 猜想, 由 Raynaud 证明 [Ray94]).
\item[(e)] 当 $C$ 固有、连通且光滑, 亏格 $g\geq2$ 时, $C$ 的 (完整) 基本群编码了关于 $C$ 的极其深刻的信息. 例如, 提及下面这个 Tamagawa 的引人注目的结果:
% Source PDF 234 / printed p.222.
\begin{origstatement}{定理}{[Tam04]}
设 $k$ 为特征 $p>0$ 的代数闭域, $A=\Spec k[[t]]$, $X$ 为 $S$ 上具有亏格 $g\geq2$ 的连通几何纤维的固有光滑曲线. 令 $s=\Spec k$. 假定特殊纤维 $X_s$ 可定义在 $\Spec k_0$ 上, 其中 $k_0$ 为 $k$ 的一个有限子域. 设 $\bar\eta$ 为 $S$ 的一般点 $\eta$ 之上的几何点. 则若 (8.5.14.1) 的特殊化同态
\[\pi_1(X_{\bar\eta},b)\longrightarrow\pi_1(X_s,a)\]
为同构, 则 $X$ 在 $S$ 上为常值, 即同构于 $X_s\times_sS$.
\end{origstatement}
% Original prints A=Spec k[[t]] and subsequently uses S; retained, see review.
\end{enumerate}
\end{origstatement}

\subsection*{E. Abel 簇}
\paragraph{8.5.22.}
现在讨论 Abel 簇的提升. 设 $S$ 为概形. 回顾: $S$ 上的 \emph{Abel 概形} (当 $S$ 为域的谱时称 \emph{Abel 簇}) 是指一个固有且平坦的 $S$ 群概形, 其几何纤维约化且不可约. 设 $X$ 是 $S$ 上的 Abel 概形. 则 $X$ 自动\emph{光滑且交换}; 当 $S$ 为代数闭域的谱时, 见 [Mum70, II 4], 一般情形见 [Mum65, 6.5]. 还已知, 若 $S$ 正规, 甚至只要几何单支, 则 $X$ 在 $S$ 上射影 [Mur95]. Raynaud [Ray70] 已给出了不满足这些假设时的反例.

Grothendieck 证明了 (未发表) Abel 簇具有形式提升:
\begin{origstatement}{定理}{8.5.23}\label{ch08:th8523}
设 $\widehat S=\operatorname{Spf}A$ 如 8.5.12 中所述, $X_0$ 是 $s=\Spec k$ 上的 Abel 簇. 则:
\begin{enumerate}
\item[(a)] 存在 $\widehat S$ 上的固有光滑形式概形 $\mathcal X$, 使 $s\times_{\widehat S}\mathcal X=X_0$, 并有 $\mathcal X$ 在 $\widehat S$ 上的截面 $e$ 延拓 $X_0$ 的单位截面 $e_0$.
\item[(b)] 设 $(\mathcal X,e)$ 为如 (a) 中所述的 $(X_0,e_0)$ 在 $\widehat S$ 上的提升, 令 $X_n=S_n\times_{\widehat S}\mathcal X$, 其中 $S_n$ 如 8.2.12 中所述. 可唯一地归纳定义 $X_n$ 上的 $S_n$ Abel 概形结构, 使其以 $e_n$ 为单位截面, 且作为 Abel 概形有 $X_n=S_n\times_{S_{n+1}}X_{n+1}$.
\end{enumerate}
\end{origstatement}
假定对固定整数 $n$, 已构造提升 $X_0$ 的 Abel 概形 $X_n$ (单位截面为 $e_n$), 需说明:
\begin{enumerate}
\item[(i)] 存在 $S_{n+1}$ 上提升 $X_n$ 的光滑概形 $X_{n+1}$, 以及 $e_n$ 的提升 $e_{n+1}$;
\item[(ii)] 给定 $X_n$ 作为概形的光滑提升 $X_{n+1}$ 以及 $e_n$ 的提升 $e_{n+1}$, 存在唯一的 $S_{n+1}$ 群概形结构, 提升 $X_n$ 上的 $S_n$ 群概形结构, 并以 $e_{n+1}$ 为单位截面.
\end{enumerate}
(i)、(ii) 的证明类似. 在两种情形都会遇到位于某个非零上同调群中的障碍. 利用障碍对于适当态射的函子性 (8.5.10 (a)), 可说明该障碍为零.

概述 (i) 的证明 (参见 [Oor71, p.~238]). 考虑将 $X_n$ 提升到 $S_{n+1}$ 的障碍 (8.5.9 (b))
\[o(X_n)\in H^2(X_0,T_{X_0})\otimes I,\]
其中为简洁起见, 以 $T_{X_0}$ 记 $T_{X_0/s}$, 以 $I$ 记 $\mathfrak m^n/\mathfrak m^{n+1}$. 同时考虑将 $X_n\times X_n$ 提升到 $S_{n+1}$ 的障碍
\[o(X_n\times X_n)\in H^2(X_0\times X_0,T_{X_0\times X_0})\otimes I.\]
% Source PDF 235 / printed p.223.
由障碍与乘积的相容性 (8.5.10 (a)), 有
\begin{equation}o(X_n\times X_n)=\operatorname{pr}_1^*o(X_n)+\operatorname{pr}_2^*o(X_n).\tag{1}\label{ch08:eq85231}\end{equation}
令
\[s:X_n\times X_n\longrightarrow X_n,\qquad(x,y)\longmapsto x+y\]
为求和态射. 由障碍的函子性 (8.5.10 (a)), $o(X_n)$ 与 $o(X_n\times X_n)$ 在下列两个映射下有相同的像:
\[\begin{aligned}
H^2(X_0,T_{X_0})\otimes I&\longrightarrow H^2(X_0\times X_0,s^*T_{X_0})\otimes I\\
&\longleftarrow H^2(X_0\times X_0,T_{X_0\times X_0})\otimes I.
\end{aligned}\]
这两个映射可重写为
\[\begin{aligned}
H^2(X_0)\otimes t_{X_0}\otimes I&\xrightarrow{s^*\otimes\id}H^2(X_0\times X_0)\otimes t_{X_0}\otimes I\\
&\xleftarrow{\id\otimes s}H^2(X_0\times X_0)\otimes t_{X_0\times X_0}\otimes I,
\end{aligned}\]
其中以 $H^*(-)$ 代替 $H^*(-,\mathcal O)$, $t$ 表示原点处的切空间, 即切丛 $T$ 沿单位截面的拉回. 换言之, 有
\begin{equation}(s^*\otimes\id)(o(X_n))=(\id\otimes s)(o(X_n\times X_n)).\tag{2}\label{ch08:eq85232}\end{equation}
另一方面, 已知 [Ser59, VII 21]
\[H^q(X_0)=\Lambda^qH^1(X_0),\qquad H^q(X_0\times X_0)=\Lambda^qH^1(X_0\times X_0),\]
而
\[s^*:H^1(X_0)\longrightarrow H^1(X_0\times X_0)=H^1(X_0)\oplus H^1(X_0)\]
为对角映射. 选取 $H^1(X_0)$ 的一组基 $(e_i)$ ($1\leq i\leq g$, $g=\dim X_0$), 以及 $t_{X_0}$ 的一组基 $e_k$ ($1\leq k\leq g$) (事实上, $H^1(X_0)$ 与 $t_{X_0}$ 自然对偶, 可以选取对偶基, 但不需要这一点). 写作
\[o(X_n)=\sum_{\substack{1\leq i<j\leq g\\1\leq k\leq g}}a_{ij}^ke_i\wedge e_j\otimes e_k,\]
其中 $a_{ij}^k\in I$. 令 $e_i'=\operatorname{pr}_1^*e_i$, $e_i''=\operatorname{pr}_2^*e_i$, $e_k'=(e_k,0)$, $e_k''=(0,e_k)$. 由 (1), 有
\begin{equation}o(X_n\times X_n)=\sum a_{ij}^ke_i'\wedge e_j'\otimes e_k'+\sum a_{ij}^ke_i''\wedge e_j''\otimes e_k''.\tag{3}\label{ch08:eq85233}\end{equation}
利用 $s^*$ 为对角映射, 从而把 $e_i$ 映到 $e_i'+e_i''$, 得到
\[(s^*\otimes\id)(o(X_n))=\sum a_{ij}^k(e_i'\wedge e_j'+e_i''\wedge e_j'+e_i'\wedge e_j''+e_i''\wedge e_j'')\otimes e_k.\]
最后, 因为 $s:t_{X_0\times X_0}\to t_{X_0}$ 将 $e_k',e_k''$ 均映到 $e_k$, 由 (2) 得到
\[\begin{aligned}
&\sum a_{ij}^k(e_i'\wedge e_j'+e_i''\wedge e_j'+e_i'\wedge e_j''+e_i''\wedge e_j'')\otimes e_k\\
&\hspace{2em}=\sum a_{ij}^k(e_i'\wedge e_j'+e_i''\wedge e_j'')\otimes e_k.
\end{aligned}\]
因此对所有 $i,j,k$, 有
\[a_{ij}^k=0,\]
即 $o(X_n)=0$. 单位截面 $e_n$ 的提升 $e_{n+1}$ 的存在性是直接的.

(ii) 的证明见 [Mum65, 6.15]: 首先利用障碍与对角映射 $x\mapsto(x,x)$ 及映射 $x\mapsto(x,0)$ 的复合的相容性 (8.5.10 (a)), 说明提升差映射 $\mu_n:X_n\times_{S_n}X_n\to X_n$, $(x,y)\mapsto x-y$ 的障碍为零; 用 $e_{n+1}$ 将 $\mu_n$ 的提升规范化; 最后由一个刚性引理得出结论.

% Source PDF 236 / printed p.224.
\begin{origstatement}{注}{8.5.24}\label{ch08:rem8524}
\begin{enumerate}
\item[(a)] 使用上述论证, Grothendieck 实际证明了比 8.5.23 更一般且更精确的结果: 若 $S_0\to S$ 是仿射概形的闭浸入, 由平方为零的理想 $I$ 定义, 且 $X_0$ 是 $S_0$ 上的 Abel 概形, 则存在 $S$ 上提升 $X_0$ 的 Abel 概形 $X$; 此外, $S$ 上提升 $X_0$ 的 Abel 概形 $X$ 的同构类集合是 $\Gamma(S_0,t_{\widehat X_0}\otimes t_{X_0}\otimes I)$ 下的仿射空间, 其中 $\widehat X_0$ 是对偶 Abel 概形, 而任意提升 $X$ 的 (在 $X_0$ 上诱导恒等的) 自同构群为零. [Ill85, A 1.1] 使用余切复形理论给出了另一证明, 该理论给出把 $X_0$ 作为平坦交换群概形提升的障碍, 而该障碍所在的上同调群为零.
\item[(b)] 考虑如 8.5.23 (b) 中的形式 Abel 概形 $\mathcal X=\operatorname{colim}X_n$. 一般地, $\mathcal X$ 并不一定可代数化: 使用 Abel 簇的形式模理论, 甚至在 $k=\mathbb C$, $A=\mathbb C[[t]]$, $g=\dim X_0=2$ 时, 已可构造不可代数化的 $\mathcal X$ 的例子. 与曲线情形不同, $X_0$ 上的丰沛可逆层 $L_0$ 确实不总能提升到 $\mathcal X$ (甚至 $X_1$). 这种逐步提升的障碍位于形如 $H^2(X_0,\mathcal O)\otimes I$ 的群中, 当 $g\geq2$ 时该群非零, 而这些障碍也可能非零.

另一方面, Mumford 证明了\emph{正特征的任意 Abel 簇都可提升到特征零} [Mum69]. 更确切地, 若 $k$ 是特征 $p>0$ 的代数闭域, $X_0$ 是 $k$ 上的 Abel 簇, 则存在以 $k$ 为剩余域、分式域特征为零的完备离散赋值环 $A$, 以及 $\Spec A$ 上的一个 (射影) Abel 概形 $X$, 使 $X\otimes k=X_0$ (环 $A$ 是 $k$ 上 Witt 向量环 $W(k)$ 的一个有限扩张, 一般是分歧的).
\end{enumerate}
\end{origstatement}

\subsection*{F. 曲面}
\paragraph{8.5.25.}
设 $Y$ 为局部 Noether 概形. 所谓 $Y$ 上的\emph{曲面}, 是指一个在 $Y$ 上平坦、分离且有限型, 相对维数为 2 的概形 $X$. 我们只考虑固有光滑曲面. 由 Zariski 的一个定理 ([Zar58]、[Har70]), 域上的固有光滑曲面都是射影的. 与曲线和 Abel 簇的情形不同, 域上存在不能作形式提升的固有光滑曲面. 更确切地, 设 $k$ 是特征 $p>0$ 的代数闭域. 不可提升现象有两种类型.

(a) \emph{不能提升到 $W_2$.} 令 $W=W(k)$ 为 $k$ 上的 Witt 向量环, $W_n=W/p^nW$ 为长度 $n$ 的 Witt 向量环. 设 $X_0$ 是 $k$ 上的一个具有非闭的 1 次整体微分形式的固有光滑曲面. Mumford [Mum61a] 构造了这样的曲面, 后来 Lang [Lan79]、Raynaud 和 Szpiro (见 [Fos81]) 也给出了构造. 由 Deligne--Illusie 的一个定理 [DI87, 2.4], 这一病态现象阻止 $X_0$ 提升到 $W_2$.

(b) \emph{不能提升到特征零.} Mumford [Mum] 改进 Serre 的一个结果 [Ser61], 构造了 $s=\Spec k$ 上具有如下性质的固有光滑曲面 $X_0$ 的例子: 对任何剩余域为 $k$、分式域特征为零的整完备局部 Noether 环 $A$, 都不存在 $\Spec A$ 上的固有光滑概形 $X$, 使 $X_s=X_0$.

利用 Hodge--Witt 数 (用 de Rham--Witt 复形定义的 $X_0$ 的深刻不变量),
% Source PDF 237 / printed p.225.
Ekedhal [Eke86, p.~114] 观察到, (a) 中提到的适当 Raynaud 曲面也提供类似例子.

(a)、(b) 两类现象之间的关系还没有得到很好的理解.

\paragraph{8.5.26.}
下面给出一些肯定方向的结果. 如 8.5.12, 设 $A$ 为完备局部 Noether 环, 极大理想为 $\mathfrak m$, 剩余域为 $k$. 令 $S=\Spec A$, $\widehat S=\operatorname{Spf}A=\operatorname{colim}S_n$, 其中 $S_n=\Spec A/\mathfrak m^{n+1}$. 设 $X_0$ 是 $s=\Spec k$ 上的固有光滑曲面. 利用 8.5.19 及 [Har77, IV 2, 5] 的一般结果, 容易看出, 若 $X_0$ 有理或为直纹曲面, 则 $X_0$ 可提升到 $S$ 上的射影曲面. 另一方面, 已经看到若 $X_0$ 是 Abel 曲面, 则 $X_0$ 具有 $\widehat S$ 上的形式光滑提升 $\mathcal X$. 若 $X_0$ 是 K3 曲面, 即固有、光滑连通曲面, 满足 $\Omega^2_{X_0/s}$ 平凡且 $H^1(X_0,\mathcal O_{X_0})=0$, 则同样成立. 更确切地, 有下面这个 Rudakov--Shafarevitch 与 Deligne 的结果:

\begin{origstatement}{定理}{8.5.27}\label{ch08:th8527}
[Del81, 1.8] 用 8.5.26 的记号, 设 $X_0$ 是代数闭域 $k$ 上的 K3 曲面.

(a) 存在 $\widehat S$ 上提升 $X_0$ 的固有光滑形式概形 $\mathcal X$.

(b) 设 $L_0$ 是 $X_0$ 上的丰沛线丛. 则存在一个在 Witt 向量环 $W(k)$ 上有限的完备离散赋值环 $R$, 一个 $T=\Spec R$ 上提升 $X_0$ 的固有光滑概形 $X$, 以及 $L_0$ 到 $X$ 上某个丰沛线丛 $L$ 的提升.
\end{origstatement}
证明 (a). 由 Rudakov--Shafarevitch 的一个基本结果 [RS76] (另见 [Nyg79] 与 [LN80]), 有
\[H^0(X_0,T_{X_0/k})=0.\]
因为 $\Omega^2_{X_0/k}$ 平凡, 有 $T_{X_0/k}=\Omega^1_{X_0/k}$, 所以由 Serre 对偶得到 $H^2(X_0,T_{X_0/k})=0$. 因此结论来自 8.5.19. (b) 的证明困难得多, 因为 $H^2(X_0,\mathcal O)=k$, 不能应用 8.5.19. 细节见 [Del81].

\begin{origstatement}{注}{8.5.28}\label{ch08:rem8528}
\begin{enumerate}
\item[(a)] 与 Abel 簇的情形一样, 在 8.5.27 (a) 的情形下, $X_0$ 的给定极化可能无法提升到 $\mathcal X$; 更精确的陈述见 [Del81, 1.6].
\item[(b)] 当 $p=3$ 时, M. Hirokado [Hir99] 构造了一个 Calabi--Yau 三维簇 $X_0/k$ (即满足 $\Omega^3_{X_0/k}\simeq\mathcal O_{X_0}$ 及 $H^1(X_0,\mathcal O_{X_0})=H^2(X_0,\mathcal O_{X_0})=0$ 的三维光滑射影概形), 使 $b_3=0$, 其中
\[b_3=\dim H^3(X_0,\mathbb Q_\ell),\qquad\ell\ne p.\]
由 Hodge 理论, 这样的概形没有到特征零的光滑射影提升. 这个 Calabi--Yau 三维簇由 $\mathbb P_k^3$ 的一个爆破对某个向量场取商而构造. 因此, 虽然 Calabi--Yau 三维簇可视为 K3 曲面的类似物, Deligne 的结果 8.5.27 却不能推广到维数 3.
\end{enumerate}
\end{origstatement}

% Source PDF 238 / printed p.226.
\subsection*{G. 余切复形}
\paragraph{8.5.29.}
至此我们只讨论了光滑态射的形变. 要处理更一般的态射, 必须使用余切复形理论 [Ill71]. 全面的综述见 [Ill72b]. 这里只给出非常简短的说明.

设 $f:X\to Y$ 为概形态射. $f$ (或 $X/Y$) 的\emph{余切复形}, 记作
\[L_{X/Y},\]
是一个次数集中在 $\leq0$ 的 $\mathcal O_X$ 模复形, 定义如下.

一对函子: 由集合层生成的自由 $f^{-1}(\mathcal O_Y)$ 代数, 与 $f^{-1}(\mathcal O_Y)$ 代数的底集合层, 定义了一个 Godement 风格的、增广到 $\mathcal O_X$ 的标准单纯 $f^{-1}(\mathcal O_Y)$ 代数 $P$, 其分量是集合层上的自由 $f^{-1}(\mathcal O_Y)$ 代数, 而且底增广单纯 $f^{-1}(\mathcal O_Y)$ 模对应的链复形无上同调. 逐分量施行函子 $\Omega^1$ (Kähler 微分), 得到单纯 $f^{-1}(\mathcal O_Y)$ 模 $\Omega^1_{P/f^{-1}(\mathcal O_Y)}\otimes_P\mathcal O_X$, 其对应的链复形就是 $L_{X/Y}$. 此复形有一个到 $\Omega^1_{X/Y}$ 的自然增广, 定义同构 $\mathcal H^0(L_{X/Y})\xrightarrow\sim\Omega^1_{X/Y}$. 它的各分量都是平坦 $\mathcal O_X$ 模. 它对 $X/Y$ 具有函子性.

此外, 态射序列 $X\xrightarrow{f}Y\to Z$ 给出 $D(X)$ 中的一个区别三角, 称为\emph{传递性三角}
\[f^*L_{Y/Z}\longrightarrow L_{X/Z}\longrightarrow L_{X/Y}\longrightarrow.\]

假定 $f$ 是局部 Noether 概形之间的局部有限型态射. 则 $L_{X/Y}$ 伪凝聚 (8.3.6.1). 若 $f$ 光滑, 则增广 $L_{X/Y}\to\Omega^1_{X/Y}$ 为拟同构. 若 $f$ 是一个由理想 $I$ 定义的闭浸入, 则存在自然增广 $L_{X/Y}\to I/I^2[1]$; 当 $f$ 为正则浸入 (即局部由一个正则序列定义) 时, 该增广为拟同构; 此时 $I/I^2$ 局部自由. 若 $f$ 为局部完全交, 即局部地 (在 $X$ 上) 是一个正则浸入与一个光滑态射的复合, 则 $L_{X/Y}$ 完美, 其完美幅度在 $[-1,0]$ 中 (8.3.6.3).

\paragraph{8.5.30.}
余切复形与形变理论的关系来自下列事实. 设 $f:X\to Y$ 为概形态射. 若 $i:X\to X'$ 是到一个 $Y$ 概形中的闭浸入, 由平方为零的理想 $I$ 定义, 则 $I$ 是 $X$ 上的拟凝聚模. 我们称 $i$ (或 $X'$) 为 $X$ 以 $I$ 为核的一个 $Y$ \emph{扩张}. 对固定的 $I$, 这些 $Y$ 扩张组成一个 Abel 群, 可证明它典范同构于 $\operatorname{Ext}^1(L_{X/Y},I)$. 此同构对 $I$ 具有函子性. 利用传递性三角 (8.5.29), 容易导出 8.5.9 的如下推广 (见 [Ill71, III 2]):

\begin{origstatement}{定理}{8.5.31}\label{ch08:th8531}
(a) 设 $X,Y$ 是概形 $S$ 上的概形, $j:X_0\to X$ 为一个由平方为零的理想 $J$ 定义的闭子概形. 设 $g:X_0\to Y$ 是 $S$ 态射. 存在障碍
\[o(g,j)\in\operatorname{Ext}^1(g^*L_{Y/S},J),\]
其消失是存在延拓 $g$ 的 $S$ 态射 $h:X\to Y$ (即 $hj=g$) 的充要条件. 当 $o(g,j)=0$ 时, $g$ 的延拓 $h$ 的集合
% Source PDF 239 / printed p.227.
是 $\operatorname{Ext}^0(g^*L_{Y/S},J)=\Hom(g^*\Omega^1_{Y/S},J)$ 下的仿射空间.

(b) 设 $i:S_0\to S$ 为由平方为零的理想 $I$ 定义的 1 阶加厚, $X_0$ 为平坦 $S_0$ 概形. 存在障碍
\[o(X_0,i)\in\operatorname{Ext}^2(L_{X_0/S_0},f_0^*I)\]
(其中 $f_0:X_0\to S_0$ 为结构态射), 其消失是存在 $X_0$ 在 $S$ 上的形变 $X$ 的充要条件 (8.5.7). 当 $o(X_0,i)=0$ 时, 这样的形变的同构类集合是 $\operatorname{Ext}^1(L_{X_0/S_0},f_0^*I)$ 下的仿射空间, 一个固定形变的自同构群同构于 $\operatorname{Ext}^0(L_{X_0/S_0},f_0^*I)=\Hom(\Omega^1_{X_0/S_0},f_0^*I)$.
\end{origstatement}

下面给出对某些奇异曲线提升的一个应用 (推广 8.5.19 中处理的光滑情形):
\begin{origstatement}{推论}{8.5.32}\label{ch08:cor8532}
设 $S=\Spec A$ 如 8.5.19 中所述, $X_0$ 是 $s$ 上的固有曲线 (8.5.18). 假定 $X_0$ 在 $s$ 上为局部完全交, 且除有限个闭点外在 $s$ 上光滑. 则存在 $S$ 上的射影平坦曲线 $X$, 使 $X_s=X_0$.
\end{origstatement}
注意, 这样的提升 $X$ 自动在 $S$ 上为局部完全交 ([EGAIV3, 11.3.8]、[EGAIV4, 19.2.4]), 且除一个有限子概形外在 $S$ 上光滑 ($X/S$ 的非光滑轨迹是闭的, 其特殊纤维有限, 故由 8.2.16, 该轨迹有限).

如 8.5.19 的证明, 先说明存在 $\widehat S$ 上提升 $X_0$ 的 (固有) 平坦形式概形 $\mathcal X=\operatorname{colim}X_n$. 假定对 $m\leq n$, 已构造 $S_m$ 上平坦的 $X_m$, 满足 $X_m=S_m\times_{S_n}X_n$, 令 $i_n:S_n\to S_{n+1}$ 为包含映射. 则由 8.5.31, 存在障碍
\[\begin{aligned}
o(X_n,i_n)&\in\operatorname{Ext}^2(L_{X_0/s},\mathcal O_{X_0}\otimes\mathfrak m^{n+1}/\mathfrak m^{n+2})\\
&=\operatorname{Ext}^2(L_{X_0/s},\mathcal O_{X_0})\otimes_k\mathfrak m^{n+1}/\mathfrak m^{n+2},
\end{aligned}\]
它是存在 $X_n$ 在 $S_{n+1}$ 上的平坦提升 $X_{n+1}$ 的障碍. 因而只需说明
\begin{equation}\operatorname{Ext}^2(L_{X_0/s},\mathcal O_{X_0})=0.\tag{*}\label{ch08:eq8532star}\end{equation}
有
\[\operatorname{Ext}^2(L_{X_0/s},\mathcal O_{X_0})=H^2(X_0,R\mathcal Hom(L_{X_0/s},\mathcal O_{X_0})).\]
因为 $L_{X_0/s}$ 的完美幅度在 $[-1,0]$ 中, $R\mathcal Hom(L_{X_0/s},\mathcal O_{X_0})$ 的完美幅度在 $[0,1]$ 中, 特别地, 当 $i\ne0,1$ 时,
\[\mathcal Ext^i(L_{X_0/s},\mathcal O_{X_0})=0.\]
所以只需说明
\begin{equation}H^2(X_0,\mathcal Hom(L_{X_0/s},\mathcal O_{X_0}))=0;\tag{1}\label{ch08:eq85321}\end{equation}
\begin{equation}H^1(X_0,\mathcal Ext^1(L_{X_0/s},\mathcal O_{X_0}))=0.\tag{2}\label{ch08:eq85322}\end{equation}
(1) 显然成立, 因为 $X_0$ 的维数为 1. 因为 $X_0$ 在某个有限闭子集 $\Sigma$ 之外在 $s$ 上光滑, $\mathcal Ext^1(L_{X_0/s},\mathcal O_{X_0})$ 集中在 $\Sigma$ 上, 这蕴含 (2), 从而得到 (*).

还需说明 $\mathcal X$ 可代数化为 $S$ 上的射影概形. 若 $D$ 是支集包含于 $X_0$ 的光滑轨迹且与 $X_0$ 的每个不可约分支相交的任意有效除子, 则 $\mathcal O_{X_0}(D)$ 丰沛; 因为 $H^2(X_0,\mathcal O_{X_0})=0$, 结论来自 8.5.6.

% Source PDF 240 / printed p.228.
\begin{origstatement}{注}{8.5.33}\label{ch08:rem8533}
可以证明, 在 8.5.32 的假设下, 自然映射 $L_{X_0/s}\to\Omega^1_{X_0/s}$ 是 $D(X_0)$ 中的同构.
\end{origstatement}

\section{Serre 的例子}\label{ch08:sec86}
\paragraph{8.6.1.}
设 $k$ 为特征 $p>0$ 的代数闭域, $n\geq0$, $r\geq1$ 为整数, $G$ 为有限群, 并有表示
\[\rho_0:G\longrightarrow\operatorname{PGL}_{n+1}(k)\quad(=\operatorname{GL}_{n+1}(k)/k^*).\]
令 $P_0=\mathbb P_k^n$. 因为 $P_0$ 的 $k$ 自同构群为 $\operatorname{PGL}_{n+1}(k)$ [Har77, II 7.1.1], $\rho_0$ 定义 $G$ 在 $P_0$ 上的一个 (右) 作用. 对 $g\in G$, 以 $\operatorname{Fix}(g)$ 记 $g$ 的不动点 (闭) 子概形 (即 $P_0\times_kP_0$ 中 $g$ 的图与对角线之交). 设 $Q_0\subset P_0$ 为所有 $g\ne e$ 的 $\operatorname{Fix}(g)$ 的并. 考虑条件
\begin{equation}r+\dim(Q_0)<n.\tag{8.6.1.1}\label{ch08:eq8611}\end{equation}
Serre 构造的出发点是如下结果 [Ser58, Prop.~15]:

\begin{origstatement}{命题}{8.6.2}\label{ch08:prop862}
假定 (8.6.1.1) 成立. 则存在整数 $d_0\geq1$, 使对任何被 $d_0$ 整除的整数 $d$, 都可在 $P_0$ 中找到一个维数为 $r$ 的光滑完全交 $Y_0=V(h_1,\ldots,h_{n-r})$, 满足当 $1\leq i\leq r$ 时 $\deg(h_i)=d$, 它在 $G$ 下稳定且 $G$ 在其上自由作用.
\end{origstatement}
% Original upper bound on deg(h_i) is r, rather than n-r; retained.
由 [SGA2, V 1.8], $G$ 在 $P_0$ 上的作用是可容许的, 特别地, 商 $Z_0=P_0/G$ 存在. 投影 $f:P_0\to Z_0$ 有限, 且 $(f_*\mathcal O_{P_0})^G=\mathcal O_{Z_0}$. 由 [EGAII, 6.6.4], $Z_0$ 射影 (事实上, [EGAII, 6.5.1] 的条件 (II bis) 满足: 因为 $P_0$ 正规, $Z_0$ 也正规, 这来自上述 $\mathcal O_{Z_0}$ 的公式). 选取一个嵌入 $i:Z_0\to\mathbb P_k^s$. 则对某个整数 $d_0>0$, 有 $(if)^*\mathcal O_{\mathbb P_k^s}(1)=\mathcal O_{P_0}(d_0)$. 对任意整数 $m\geq1$, 以 $i_m:Z_0\to\mathbb P_k^{N(m)}$ ($N(m)=\binom{s+m}{s}-1$) 记 $i$ 的 $m$ 次倍嵌入. 则 $(i_mf)^*\mathcal O_{\mathbb P_k^{N(m)}}(1)=\mathcal O_{P_0}(d)$, 其中 $d=md_0$. 因为 $f$ 有限, $f(Q_0)$ 在 $Z_0$ 中闭且 $\dim(f(Q_0))=\dim(Q_0)$. 因为 (8.6.1.1) 成立, 由 Bertini 的一个定理 [Jou79, 6.11], 存在 $\mathbb P_k^{N(m)}$ 中余维为 $n-r$ 的线性子空间 $L_0=V(\ell_1,\ldots,\ell_{n-r})$ (其中 $\deg(\ell_i)=1$), 使 $L_0\cap Z_0$ 包含于 $U_0=Z_0-f(Q_0)$ 且 $L_0$ 与 $U_0$ 横截. 因为 $f|U_0:f^{-1}(U_0)\to U_0$ 平展, 齐次多项式 $h_i=(i_mf)^*\ell_i\in\Gamma(P_0,\mathcal O(d))$ ($1\leq i\leq n-r$) 在 $P_0$ 中定义一个光滑完全交 $Y_0$, 它在 $G$ 下稳定, 且不与 $Q_0$ 相交, 因而 $G$ 在其上自由作用.

\paragraph{8.6.3.}
设 $d,Y_0$ 如 8.6.2, 且 $d\geq2$, 令
\[X_0=Y_0/G\]
为 $Y_0$ 对 $G$ 的商. 因为 $G$ 在 $Y_0$ 上自由作用, $X_0/k$ 是维数 $r$ 的光滑射影概形, 投影
\[f:Y_0\longrightarrow X_0\]
是群为 $G$ 的平展覆盖 [SGA1, V 2.3]. 此外, 因为 $Y_0$ 是维数 $r\geq1$ 的完全交, $Y_0$ 连通 [Ser55, n\textsuperscript{o}78, Prop.~5].

Serre 构造的关键是下述结果.
% Source PDF 241 / printed p.229.
\begin{origstatement}{命题}{8.6.4}\label{ch08:prop864}
假定 $r\geq3$, 或者 $r=2$, $(p,n+1)=1$ 且 $p\mid d$. 设 $A$ 是剩余域为 $k$ 的完备局部 Noether 环. 设 $\mathcal X$ 是 $A$ 上提升 $X_0$ 的平坦形式概形. 则 $\mathcal X$ 可代数化, 即 (8.4.9) 存在 (唯一的) 固有概形 $X/A$, 使 $\widehat X=\mathcal X$. 此外, $X$ 在 $A$ 上射影且光滑, 而表示 $\rho_0$ (8.6.1) 可提升为表示
\[\rho:G\longrightarrow\operatorname{PGL}_{n+1}(A)\quad(=\operatorname{GL}_{n+1}(A)/A^*).\]
\end{origstatement}
$r\geq3$ 的情形已在 [Ser61] 中处理. $r=2$ 的情形归功于 Mumford [Mum].

由 8.5.9 (b), $Y_0$ 唯一地提升为 $\mathcal X=\operatorname{colim}X_m$ 的一个形式平展覆盖 $\mathcal Y=\operatorname{colim}Y_m$ (相差一个唯一同构), 即对所有 $m\geq0$, $Y_m$ 在 $X_m$ 上有限平展 (其中 $X_m=\mathcal X\otimes A/\mathfrak m^{m+1}$). 再由 8.5.9 (b), $G$ 在 $Y_0$ 上的作用 (唯一地) 延拓为 $G$ 在 $\mathcal Y$ 上的作用, 使 $\mathcal Y$ 成为 $\mathcal X$ 的群为 $G$ 的平展 Galois 覆盖 (即一个 $G$ Galois 平展覆盖 $Y_m\to X_m$ 的归纳系). 因为 $r\geq2$, 有 $H^1(Y_0,\mathcal O_{Y_0})=0$ 及 $H^0(Y_0,\mathcal O_{Y_0})=k$ ([Ser55, n\textsuperscript{o}78] 或 [SGA7, IX]), 所以 8.3.11.2 蕴含对所有 $m$, $H^0(Y_m,\mathcal O_{X_m})=A_m$.
% Original prints O_{X_m} here; retained for review.

令 $i:Y_0\to P_0$ 为包含映射, $L_0=\mathcal O_{Y_0}(1)=i^*\mathcal O_{P_0}(1)$. 我们将说明:
\begin{equation}\text{$L_0$ 可提升到 $\mathcal Y$ 上的可逆层 $\mathcal L$, 且唯一, 相差一个在 $L_0$ 上诱导恒等的非唯一同构.}\tag{*}\label{ch08:eq864star}\end{equation}
先假定 $r\geq3$. 则由上述文献, $H^2(Y_0,\mathcal O_{Y_0})=0$. 因为也有 $H^1(Y_0,\mathcal O_{Y_0})=0$, 由 8.5.5, (*) 成立. 现在假定 $r=2$. 此时 $H^2(Y_0,\mathcal O_{Y_0})=0$ 不再成立. 为说明 $L_0$ 可提升 (若可提升, 则因 $H^1(Y_0,\mathcal O_{Y_0})=0$ 而唯一), Mumford 作如下论证. 有
\begin{equation}\Omega^2_{Y_0/k}=\mathcal O_{Y_0}(N),\tag{**$_0$}\label{ch08:eq864star20}\end{equation}
其中 $N=(n-r)d-n-1$. 假设蕴含 $(p,N)=1$. 假定对 $m\geq0$, $L_0$ 已提升为 $X_m$ 上的可逆层 $L_m$, 且同构 (**$_0$) 已提升为同构
\begin{equation}L_m^{\otimes N}\simeq\Omega^2_{Y_m/A_m}.\tag{**$_m$}\label{ch08:eq864star2m}\end{equation}
% Original prints X_m rather than Y_m in the preceding sentence; retained.
令 $i_m:Y_m\to Y_{m+1}$ 为包含映射. 考虑把 $L_m$ 提升到 $Y_{m+1}$ 的障碍 $o(L_m,i_m)$ (8.5.3 (b)). 由 8.5.4 (a), 有
\[o(L_m^{\otimes N},i_m)=No(L_m,i_m).\]
因为 $\Omega^2_{Y_{m+1}/A_{m+1}}$ 提升 $\Omega^2_{Y_m/A_m}$, 同构 (**$_m$) 蕴含 $o(L_m^{\otimes N},i_m)=0$, 从而 $o(L_m,i_m)=0$, 因为 $p$ 不整除 $N$. 所以 $L_m$ 可提升为 $Y_{m+1}$ 上的可逆层 $L_{m+1}$. 因为 $H^1(Y_0,\mathcal O_{Y_0})=0$, 由 8.5.3 (a), 同构 (**$_m$) 可提升为同构 (**$_{m+1}$). 因而 $L_0$ 可提升为 $\mathcal Y$ 上的可逆层 $\mathcal L$.

因为 $L_0$ 丰沛, 由 8.5.6, 存在 $A$ 上的射影平坦概形 $Y$, 使 $\widehat Y=\mathcal Y$, 并有 $Y$ 上的丰沛线丛 $L$, 使 $\widehat L=\mathcal L$. 由 [EGAIII2, 6.6.1], $L_0$ 的范数 $E_0=N_{Y_0/X_0}L_0$ 是 $X_0$ 上的丰沛线丛. 对 $m\in\mathbb N$, 令 $E_m=N_{Y_m/X_m}L_m$, $\mathcal E=\lim E_m$. 则 $\mathcal E$ 提升 $E_0$, 所以由 8.4.10, 存在射影概形 $X/A$, 使 $\widehat X=\mathcal X$, 并有 $X$ 上的丰沛线丛 $E$, 使 $\widehat E=\mathcal E$. 由 8.5.6 (以及 8.5.19 证明末尾的论证), $X$ 在 $A$ 上光滑. 此外, 由 8.4.7, 平展 Galois 覆盖 $\widehat Y\to\widehat X$ 由一个 (唯一的) 群为 $G$ 的平展 Galois 覆盖 $Y\to X$ 的完备化得到; 由 8.4.2,
% Source PDF 242 / printed p.230.
$E=N_{Y/X}L$.

还需说明 $\rho_0$ 可提升到 $A$. 因为 $d\geq2$, 由 [FAC n\textsuperscript{o}78], 有 $H^0(Y_0,L_0)=k^{n+1}$, $H^1(Y_0,L_0)=0$. 因而由 8.3.11.2, 对所有 $m$, 有 $H^0(Y_m,L_m)=A_m^{n+1}$, 从而由 8.2.4, $H^0(Y,L)=A^{n+1}$. 设 $g\in G$. 因为 $H^1(Y_0,\mathcal O_{Y_0})=0$, 由 8.5.3 (a) 与 8.4.2, 存在 $g:Y\xrightarrow\sim Y$ 之上的同构 $a(g):L\xrightarrow\sim L$, 即同构 $g^*L\xrightarrow\sim L$, 它唯一到相差 $L$ 的一个自同构, 并使 $a(g)_0$ 为 $G$ 在 $Y_0$ 上的作用给出的 $g:Y_0\xrightarrow\sim Y_0$ 之上的同构 $L_0\xrightarrow\sim L_0$ (后者定义良好, 相差 $L_0$ 的一个自同构). 对 $G$ 中的 $g,h$, 有 $a(g)a(h)=a(gh)$ 与 $a(e)=\id$, 相差 $L$ 的一个自同构. 因而得到表示
\[\rho:G\longrightarrow\operatorname{PGL}(H^0(Y,L))=\operatorname{PGL}_{n+1}(A),\]
将 $g$ 对应到二元组 $(g,a(g))$ 诱导的 $H^0(Y,L)=A^{n+1}$ 的自同构 $\rho(g)$; 该自同构定义良好, 相差乘以 $A^*$ 的一个元素. 此表示提升 $\rho_0$.

\paragraph{8.6.5.}
现在设 $r,n$ 为满足 $1\leq r<n$ 的整数, $G$ 为阶 $p^s$ 的 $(p,\ldots,p)$ 型群, 即 $G\simeq\mathbb F_p^s$, 其中 $s\geq n+1$. 再假定 $p\geq n+1$. 选取一个单同态 $h:G\to k$ (这里 $k$ 视为加法群). 令 $N=(u_{ij})$ 为 $n+1$ 阶幂零矩阵, 定义为当 $j=i+1$ 时 $u_{ij}=1$, 否则 $u_{ij}=0$. 对 $g\in G$, 令
\[\widetilde\rho_0(g)=\exp(h(g)N)\in\operatorname{GL}_{n+1}(k)\]
(因为 $p\geq n+1$, 这有意义), 并令 $\rho_0(g)$ 为 $\widetilde\rho_0(g)$ 在 $\operatorname{PGL}_{n+1}(k)$ 中的像. 因而得到表示
\begin{equation}\rho_0:G\longrightarrow\operatorname{PGL}_{n+1}(k),\tag{8.6.5.1}\label{ch08:eq8651}\end{equation}
它是忠实的, 因为 $h$ 单射. 对任意 $g\in G$, $g\ne e$, $\operatorname{Fix}(g)$ 由 $P_0$ 的单个有理点 $(1,0,\ldots,0)$ 组成. 特别地, 用 8.6.1 的记号, $\dim Q_0=0$, 故条件 (8.6.1.1) 满足.

\begin{origstatement}{命题}{8.6.6}\label{ch08:prop866}
假定 $p>n+1$. 设 $A$ 为剩余域为 $k$、分式域 $K$ 的特征为零的整局部环. 则不存在提升 $\rho_0$ (8.6.5.1) 的同态 $\rho:G\to\operatorname{PGL}_{n+1}(A)$.
\end{origstatement}
下面的论证归功于 Serre (私人交流). 假定存在这样的同态 $\rho$. 因为 $\rho_0$ 单射, $\rho$ 也单射, 与包含映射 $\operatorname{PGL}_{n+1}(A)\to\operatorname{PGL}_{n+1}(K)$ 的复合也单射; 仍以 $\rho$ 记该复合. 因为 $K$ 的特征为零且 $p$ 不整除 $n+1$, 这个表示可提升为一个 (忠实) 表示 $\rho':G\to\operatorname{SL}(V)$, 其中 $V=K^{n+1}$. 因为 $G$ 交换且 $K$ 的特征为零, 将标量扩张到 $K$ 的一个有限扩张后, $V$ 分解为 $G$ 稳定的一维子空间之和
\[V=\bigoplus_{1\leq i\leq n+1}V_i,\]
对应于乘积为 1 的特征标 $\chi_i:G\to\Aut(V_i)=K^*$. $\rho$ 的核 $Z$ 是各 $\chi_i$ 的核 $H_i$ ($1\leq i\leq n$) 的交. 每个 $\chi_i$ 为从 $G$ 到 $\mu_p(K)$ 的同态, 所以可视为把 $G$ 当作 $\mathbb F_p$ 向量空间后的一个线性形式 $\varphi_i$. 因为 $\sum\varphi_i=0$, 而 $G$ 在 $\mathbb F_p$ 上的维数为 $s\geq n+1$, $Z$ 不可能为零, 与 $\rho$ 的忠实性矛盾.

% Source PDF 243 / printed p.231.
\begin{origstatement}{推论}{8.6.7}\label{ch08:cor867}
设 $r,n$ 为满足 $2\leq r<n$ 且 $p>n+1$ 的整数. 令 $G=\mathbb F_p^s$, 其中 $s\geq n+1$. 则 $P_0$ 中存在维数 $r$ 的光滑射影完全交 $Y_0$, 它在 (8.6.5.1) 所构造的表示 $\rho_0$ 定义的 $G$ 在 $P_0$ 上的作用下稳定, $G$ 在其上自由作用, 并且光滑射影概形 $X_0=Y_0/G$ 具有如下性质: 对任何剩余域为 $k$、分式域 $K$ 的特征为零的整完备局部 Noether 环 $A$, 都不存在 $A$ 上提升 $X_0$ 的平坦形式概形 $\mathcal X$.
\end{origstatement}
令 $d_0\geq1$ 为具有 8.6.2 所述性质的整数. 若 $r\geq3$, 取 $d_0$ 的任意非零倍数 $d\geq2$; 若 $r=2$, 取 $d_0$ 的任意被 $p$ 整除的非零倍数 $d$. 由 8.6.2, 在 $P_0$ 中选取多重次数为 $(d,\ldots,d)$ 的光滑完全交 $Y_0$, 使它在 (8.6.5.1) 所构造的表示 $\rho_0$ 定义的 $G$ 在 $P_0$ 上的作用下稳定, 且 $G$ 在其上自由作用. 令 $X_0=Y_0/G$. 假定存在 $A$ 上提升 $X_0$ 的平坦形式概形 $\mathcal X$. 因为 $p>n+1$, 而当 $r=2$ 时 $p$ 整除 $d$, 所以 8.6.4 的假设满足; 其结论与 8.6.6 合起来得到矛盾.

最小的例子由 $r=2$, $n=3$, $s=4$, $p=5$ 得到. (在 [Ser61] 中, 最小的例子对应于 $r=3$, $n=4$, $s=5$, $p=7$.)

\begin{origstatement}{注}{8.6.8}\label{ch08:rem868}
设 $X_0$ 为 8.6.7 考虑的概形. 设 $A$ 为剩余域为 $k$ 的完备局部 Noether 环, 它是 $X_0$ 的一个形式万有形变 $\mathcal X$ 的底 [Sch68]. 这样的环 $A$ 是一个形式有限型的 $W$ 代数, 其中 $W=W(k)$ 为 $k$ 上的 Witt 向量环. 令 $K_0=W[1/p]$ 为 $W$ 的分式域. 由 8.6.7 可知 $A\otimes_WK_0=0$; 换言之, 存在整数 $n_0\geq1$, 使 $p^{n_0}A=0$. 否则, 可找到 $\Spec A$ 的一个整闭子概形 $T=\Spec B$, 其一般点特征为零. 将 $\mathcal X$ 拉回到 $\operatorname{Spf}B$, 就会得到矛盾.
\end{origstatement}

\section{Serre 的一封信}\label{ch08:sec87}
在下面的信中, Serre 证明了实际上 $n_0=1$. 他的论证还表明, 在 8.6.7 中只需假定 $s\geq2$, 无需 $s\geq n+1$.

\begin{flushright}Paris, 2003 年 10 月 11 日\end{flushright}
亲爱的 Illusie:

下面证明你知道的那个形式模簇 ``被 $p$ 消灭''.

\emph{记号.} 我考虑一个局部环 $A$, 其剩余域 $k$ 的特征为 $p$. 设 $n>1$, $n<p$. 我们考虑 $\operatorname{GL}(n,k)$ 的一个 $(p,p)$ 型子群 $G$, 它具有如下性质: 对每个 $s\in G$, $s\ne1$, $s-1$ 的核的维数为 1 (换言之, 可以用一个长度为 $n$ 的 Jordan 块表示 $s$).

\begin{origstatement}{定理}{1}
若 $G$ 可提升到 $\operatorname{GL}(n,A)$ 中, 则在 $A$ 中有 $p=0$.
\end{origstatement}
\begin{origstatement}{定理}{2}
若 $p>n$, 且 $G$ 可提升到 $\operatorname{PGL}(n,A)$ 中, 则它可提升到 $\operatorname{SL}(n,A)$ 中 (从而由定理 1, 在 $A$ 中有 $p=0$).
\end{origstatement}
当然, 对形式模簇有用的是定理 2, 第一个有趣的情形为 $n=4$; 此时当 $p\geq5$ 时定理 2 适用.

定理 2 几乎是直接的: 若 $z$ 为 $\operatorname{PGL}(n,A)$ 中的一个 $p$ 阶元素, 则可唯一地将其提升为 $\operatorname{SL}(n,A)$ 中的一个 $p$ 阶元素 $z'$. 事实上,
% Source PDF 244 / printed p.232.
选取 $s$ 在 $\operatorname{GL}(n,A)$ 中的一个提升 $z''$; 有 $z''^p=c$, 其中 $c\in A^*$. 若 $d=\det(z'')$, 则 $d^p=c^n$; 因为 $n$ 不被 $p$ 整除, 这表明 $c$ 形如 $u^p$, 其中 $u\in A^*$. 元素 $z'=z''/u$ 属于 $\operatorname{SL}(n,A)$ 且是 $z$ 的一个提升. 将此应用于 $G$ 的各个元素 (把 $G$ 视为 $\operatorname{PGL}(n,A)$ 的子群). 所得的各提升彼此交换 (这来自提升的唯一性), 并给出 $G$ 到 $\operatorname{SL}(n,A)$ 的一个嵌入. 还需说明这个 ``$G$'' 在 $\operatorname{SL}(n,k)$ 中的像确为最初的群 $G$, 但这很清楚, 因为上述 $G$ 显然包含于 $\operatorname{SL}(n,k)$.
% Original says lifting of s rather than z in first sentence; retained.

现在进入正题, 即\emph{定理 1 的证明}.

假定 $G$ 已提升到 $\operatorname{GL}(n,A)$ 中, 我们希望证明在 $A$ 中 $p=0$. 令 $\mathfrak m=\ker(A\to k)$ 为 $A$ 的极大理想. 若能证明 $p$ 属于理想 $p\mathfrak m$, 则 (由一个显然的 Nakayama 论证) 即得 $p=0$. 换言之, 只需证明 $p$ 在 $A/p\mathfrak m$ 中的像为 $0$. \emph{因此可以假定 $p\mathfrak m=0$.}

设 $(s,s')$ 为 $G$ 的一组生成元, 把 $G$ 视为 $\operatorname{GL}(n,A)$ 的子群. 我将 $s$ 写成 $s=1+e$; 用 $\underline e$ 表示 $e$ 模 $\mathfrak m$ 的约化, 可假定 $\underline e$ 为秩 $n-1$ 的标准幂零矩阵 (单个 Jordan 块). $\underline e$ 的特征多项式为 $T^n$, 它同时也是最小多项式. $e$ 的特征多项式形如 $T^n+a_1T^{n-1}+\cdots+a_n$, 其中 $a_i\in\mathfrak m$. 因为当 $i=0,\ldots,n-1$ 时, 各 $\underline e^i$ 在 $k$ 上线性无关, 所以各 $e^i$ 在 $A$ 上也线性无关, 从而环 $A[e]$ 在 $A$ 上自由且秩为 $n$, 一组基为 $1,e,\ldots,e^{n-1}$. 而且, \emph{$A$ 模 $A^n$ 在 $A[e]$ 上自由且秩为 1}. 这蕴含 $A^n$ 的每个与 $e$ 交换的自同态都属于 $A[e]$. 这特别适用于 $G$ 的第二个生成元 $s'$: 可以将 $s'$ 写成一个 $A$ 系数的 $e$ 的多项式:
\[s'=a_0+a_1e+\cdots+a_{n-1}e^{n-1},\]
其中 $a_i\in A$. 我现在利用关系 $s'^p=1$ 来得到各 $a_i$ 之间的关系. 一般地, 若 $x$ 为 $A[e]$ 的元素, 以 $t_0(x),t_1(x),\ldots$ 记 $x$ 在 $A[e]$ 的基 $1,e,\ldots,e^{n-1}$ 下的各系数.

\begin{origstatement}{命题}{}
有
\[t_1(x^p)=pt_1(x)(t_0(x)^{p-1}-t_1(x)^{p-1}).\]
\end{origstatement}
暂且承认这一断言. 将其应用于 $x=s'$, 结合 $s'^p=1$, 有 $t_1(x)^p=0$, 从而 $pu=0$, 其中 $u$ 为
\[u=a_1(a_0^{p-1}-a_1^{p-1}).\]
但我们知道 $a_0,a_1$ 在 $k$ 中的像. $a_0$ 的像显然是 1, $a_1$ 的像则是 $k$ 中一个不属于 $\mathbb Z/p\mathbb Z$ 的元素 $t$ (若该像等于 $j$, 其中 $0\leq j<p$, 则 $G$ 的元素 $s^{-j}s'$ 将满足 $\dim\ker(1-s^{-j}s')>1$). 所以在 $k$ 中有 $t-t^p\ne0$, 意味着 $u$ 可逆, 方程 $pu=0$ 就蕴含 $p=0$, 如我们所愿.
% Original prints t_1(x)^p=0 here; retained and recorded for review.

还需证明上述命题. 可以直接作多项式展开, 但会导致一大堆指标. 我将采用一种较为温和的方法. 置 $f(x)=t_1(x^p)$.

\begin{origstatement}{引理}{}
若 $y\in A[e]$ 为 $e^2$ 的倍数, 则 $f(x+y)=f(x)$.
\end{origstatement}
% Source PDF 245 / printed p.233.
写 $(x+y)^p=x^p+v+y^p$, 其中 $v$ 是一些形如 $py'$ 的项之和, $y'$ 被 $e^2$ 整除. 现在若 $y'$ 为 $e^2$ 的倍数, 则 $t_1(y')=0\pmod{\mathfrak m}$; 另一方面, $s^p=1$ 意味着 $(1+e)^p=1$, 即 $e^p=-\sum\binom pi e^i$, 其中对满足 $0<i<p$ 的指标 $i$ 求和. 因为所有 $\binom pi$ 都被 $p$ 整除, 可见 $e^p$ 被 $p$ 整除, 且 $e^{2p}=0$ (因为 $p\mathfrak m=0$, 在 $A$ 中有 $p^2=0$). 因而 $y^p=0$, 得到所求结果.

这样, 为证明命题, 可以去掉 $x$ 中的 $e^2,e^3,\ldots$ 各项, 即假定 $x$ 形如 $a+be$. 然后用与上述相同的方法计算 $f(x)=t_1(x^p)$: 写作
\[x^p=a^p+pa^{p-1}be+\cdots+b^pe^p,\]
从而
\[f(x)=t_1(a^p)+pt_1(a^{p-1}be)+pt_1(\cdots)+t_1(b^pe^p).\]
这些项的值依次为
\[0,\quad pa^{p-1}b,\quad0,\ldots,0,\quad-pb^p\]
(注意, 由上面给出的公式 $e^p=-\sum_{1\leq i\leq p-1}\binom pi e^i$, 有 $t_1(e^p)=-p$).

所以确实得到
\[f(x)=pa^{p-1}b-pb^p,\]
正如所求.

呼, 证毕, 画个小方块, 等等.

\begin{flushright}祝好,\\J-P. Serre\end{flushright}
% Source PDF 246 / printed p.234 is blank. End of Part 4 / Chapter 8.
